\documentclass[11pt]{article}

\usepackage[T1]{fontenc}
\usepackage[margin=1in]{geometry}
\usepackage{amsmath,amssymb,amsthm,mathtools}
\usepackage{etoolbox}
\usepackage{libertinus}
\usepackage{libertinust1math}
\DeclareSymbolFont{cmlargesymbols}{OMX}{cmex}{m}{n}
\DeclareMathSymbol{\intop}{\mathop}{cmlargesymbols}{"52}
\usepackage{booktabs}
\usepackage{array}
\usepackage{enumitem}
\usepackage{microtype}
\usepackage{titlesec}
\usepackage{xcolor}
\usepackage{tikz}
\usetikzlibrary{arrows.meta,positioning,calc,patterns}
\IfFileExists{numerical_certificates.zip}{%
  \usepackage{embedfile}%
  \embedfile[desc={Exact rational certificates and verification scripts}]{numerical_certificates.zip}%
}{}

% Link and figure palette.
\definecolor{CitationBlue}{HTML}{205E9D}
\definecolor{InternalLink}{HTML}{7A3E65}
\definecolor{URLTeal}{HTML}{1B645B}
\definecolor{ContentsHeading}{HTML}{40546F}
\definecolor{AccentBlue}{HTML}{2C5286}
\colorlet{AccentTeal}{URLTeal}
\definecolor{AccentBrown}{HTML}{785135}

\usepackage{hyperref}
\urlstyle{sf}
\hypersetup{
  colorlinks=true,
  linkcolor=InternalLink,
  citecolor=CitationBlue,
  urlcolor=URLTeal,
  linktoc=all,
  pdftitle={The Exact Optimal Metric Distortion of g-Stable Lotteries},
  pdfauthor={},
  pdfsubject={Metric distortion in social choice, g-stable lotteries, variational optimization},
  pdfkeywords={metric distortion, social choice, stable lotteries, activation functions, variational analysis},
  pdflang={en-US}
}
\newcommand{\citelink}[2]{%
  \hyperlink{cite.#2}{\textcolor{CitationBlue}{#1}}%
}
\usepackage[all]{hypcap}
\usepackage[capitalise,noabbrev,nameinlink]{cleveref}
\newcommand{\creflastconjunction}{, and\nobreakspace}
\crefformat{equation}{#2Eq.~(#1)#3}
\crefmultiformat{equation}{#2Eqs.~(#1)#3}{ and~#2(#1)#3}{, #2(#1)#3}{, and~#2(#1)#3}

\numberwithin{equation}{section}
\theoremstyle{plain}
\newtheorem{theorem}{Theorem}[section]
\newtheorem{proposition}[theorem]{Proposition}
\newtheorem{lemma}[theorem]{Lemma}
\theoremstyle{definition}
\newtheorem{definition}[theorem]{Definition}
\newtheorem{remark}[theorem]{Remark}

\newcommand{\E}{\mathbb{E}}
\newcommand{\R}{\mathbb{R}}
\newcommand{\1}{\mathbf{1}}
\newcommand{\lowrho}{\underline{\rho}}
\newcommand{\cost}{\operatorname{cost}}
\newcommand{\rank}{\operatorname{rank}}
\newcommand{\supp}{\operatorname{supp}}
\newcommand{\calE}{\mathcal{E}}
\newcommand{\calG}{\mathcal{G}}
\newcommand{\calH}{\mathcal{H}}
\newcommand{\calJ}{\mathcal{J}}
\newcommand{\calM}{\mathcal{M}}
\newcommand{\calZ}{\mathcal{Z}}
\newcommand{\calK}{\mathcal{K}}
\newcommand{\calQ}{\mathcal{Q}}
\newcommand{\calU}{\mathcal{U}}

\allowdisplaybreaks
\clubpenalty=10000
\widowpenalty=10000
\setlength{\emergencystretch}{2em}
\setlength{\parindent}{0pt}
\setlength{\parskip}{0.55em plus 0.1em minus 0.1em}
\setlength{\abovedisplayskip}{7.5pt plus 2pt minus 2pt}
\setlength{\belowdisplayskip}{7.5pt plus 2pt minus 2pt}
\setlength{\abovedisplayshortskip}{4pt plus 2pt}
\setlength{\belowdisplayshortskip}{6pt plus 2pt minus 2pt}
\titlespacing*{\section}{0pt}{2.4ex plus 0.6ex minus 0.3ex}{1.0ex}
\titlespacing*{\subsection}{0pt}{1.8ex plus 0.5ex minus 0.2ex}{0.7ex}
\setlist[itemize]{leftmargin=2em,itemsep=0.25em,topsep=0.35em}
\setlist[enumerate]{leftmargin=2.2em,itemsep=0.25em,topsep=0.35em}

\setcounter{tocdepth}{2}

\begin{document}

\hypersetup{pageanchor=false}
\pagenumbering{gobble}
\thispagestyle{empty}

\begin{center}
  {\LARGE\bfseries\boldmath The Exact Optimal Metric Distortion of $g$-Stable Lotteries\par}
  \vspace{0.9em}
  {\small\itshape This manuscript was AI-generated with limited human input.\par}
\end{center}
\vspace{0.6em}

\begin{abstract}
\noindent
How much can the choice of activation function reduce the metric distortion of stable lotteries? For a continuous activation $g \colon [0, 1] \to [0, 1]$, let $M(g)$ be the supremum metric distortion over all $g$-stable lotteries in all finite elections, and let $\calG$ be the class of continuous activations $g \colon [0, 1] \to [0, 1]$. We prove that a continuous strictly increasing activation $g^\star \in \calG$ attains the exact minimum
\[
  M(g^\star) \;=\; \min_{g \in \calG} M(g) \;=\; 1 + 2\lambda^\star \;\approx\; 2.1246569,
  \qquad 0.56232844 \;<\; \lambda^\star \;<\; 0.56232847.
\]
Every $g^\star$-stable lottery has metric distortion at most $1 + 2\lambda^\star$, and this bound is tight: for every $g \in \calG$ and every $\varepsilon > 0$, some finite election has a $g$-stable lottery with distortion greater than $1 + 2\lambda^\star - \varepsilon$.

For the lower bound, three families of elections with candidate distances in $\{0, 1, 2\}$ from an optimal candidate $a^\star$ reduce the search for an activation with distortion at most $1 + 2\lambda$ to a variational problem over an interval of low ranks. Compactness gives a minimizer at $\lambda^\star$, and local perturbations give supporting lines that prove Lipschitz continuity at low ranks and concavity at high ranks; completing this minimizer to $[0, 1]$ defines $g^\star$. For the matching upper bound, the same supporting lines determine the weights and potential function needed to telescope slice bounds across distance thresholds. Exact rational and outward-rounded interval certificates verify all numerical bounds in the proof.
\end{abstract}

\clearpage
\thispagestyle{empty}
\begingroup
\setlength{\parskip}{0pt}
\makeatletter
\patchcmd{\l@section}{1.0em}{0.55em}{}{}
\makeatother
\tableofcontents
\endgroup
\clearpage

\pagenumbering{arabic}
\hypersetup{pageanchor=true}

\section{Introduction and Main Result}
\label{sec:intro}

In the metric distortion model of \citelink{Anshelevich, Bhardwaj, and Postl}{AnshelevichBhardwajPostl15}~\cite{AnshelevichBhardwajPostl15} and \citelink{Anshelevich, Bhardwaj, Elkind, Postl, and Skowron}{AnshelevichEtAl18}~\cite{AnshelevichEtAl18}, voters and candidates lie in a common pseudometric space with distance $d(\cdot, \cdot)$. Each voter $v \in V$ ranks the candidates in $C$ in increasing order of distance from $v$, and a voting rule chooses a candidate, or a probability distribution (lottery) $D$ over $C$, using only the rankings. The \emph{distortion} of $D$ is the worst-case ratio between its expected social cost $\cost_d(D) = \E_{b \sim D}[\cost_d(b)]$ and the minimum candidate social cost $\min_{a \in C} \cost_d(a)$ over all pseudometrics $d$ consistent with the rankings.

A natural family of randomized voting rules is parameterized by a continuous activation function $g \colon [0, 1] \to [0, 1]$~\cite{Kreweras65,Fishburn84,ChengEtAl20,CharikarEtAl25}. Fix a lottery $D$ and place the candidates along each voter's \emph{rank line} $[0, 1]$ from worst ($0$) to best ($1$), giving each candidate $a \in C$ an interval of length $D(a)$ starting at $\rank_v(a; D) = \Pr_{b \sim D}[a \succ_v b]$. Candidate $a$'s activation on voter $v$'s rank line is the average of $g$ over this interval (or $g(\rank_v(a; D))$ when $D(a) = 0$). A lottery $D$ is \emph{$g$-stable} if no candidate's expected activation across voters exceeds the mean activation $\tau = \int_0^1 g(r)\,\mathrm{d}r$:
\[
  \E_{v \sim V}\!\left[\int_0^1 g\bigl(\rank_v(a; D) + u\,D(a)\bigr)\,\mathrm{d}u\right] \;\le\; \tau
  \qquad \text{for all } a \in C.
\]
Every finite election has at least one $g$-stable lottery, and every supported candidate in a $g$-stable lottery has expected activation equal to $\tau$.

Let $\calG$ be the class of continuous activations $g \colon [0, 1] \to [0, 1]$, and let $M(g)$ be the supremum distortion over all $g$-stable lotteries in all finite elections. Taking $g(r) = r$ gives maximal lotteries~\cite{Kreweras65,Fishburn84}, whose distortion is $3$~\cite{CharikarEtAl24}; taking $g(r) = r^k$ gives stable $k$-lotteries~\cite{ChengEtAl20,CharikarEtAl25b}, whose distortion is $7/3 \approx 2.333$ for $k \in \{2, 3\}$~\cite{CaiEtAl26b}. \citelink{Cai, Charikar, Frank, Hastings, Ramakrishnan, Wang, and Ye}{CaiEtAl26b}~\cite{CaiEtAl26b} proved that
\[
  2.1232 \;\le\; \inf_{g \in \calG} M(g) \;\le\; 2.1248.
\]
For comparison, among all randomized voting rules, \citelink{Pulyassary and Swamy}{PulyassarySwamy21}~\cite{PulyassarySwamy21} refuted the distortion-$2$ conjecture with a lower bound of $2.063164$, and \citelink{Charikar and Ramakrishnan}{CharikarRamakrishnan22}~\cite{CharikarRamakrishnan22} proved a lower bound of $2.11264$.

In this paper, we determine the exact minimum of $M(g)$ over $g \in \calG$ and show that it is attained by a continuous strictly increasing activation $g^\star$.

\begin{theorem}[Exact optimal distortion of $g$-stable lotteries]
\label{thm:main}
There is a continuous strictly increasing activation $g^\star \in \calG$, with $g^\star(0) = 0$ and $g^\star(1) = 1$, such that
\[
  M(g^\star) \;=\; \min_{g \in \calG} M(g) \;=\; 1 + 2\lambda^\star \;\approx\; 2.1246569,
\]
where $\lambda^\star \in (0.56232844,\, 0.56232847)$ is defined in \cref{eq:lambda-star-def}. Every $g^\star$-stable lottery has distortion at most $1 + 2\lambda^\star$, and this bound is tight: for every $g \in \calG$ and every $\varepsilon > 0$, some finite election has a $g$-stable lottery with distortion greater than $1 + 2\lambda^\star - \varepsilon$.
\end{theorem}

Because $M(g)$ takes the supremum over all $g$-stable lotteries, \cref{thm:main} bounds the distortion of every $g^\star$-stable lottery from above and shows that for any $g \in \calG$, some finite election has a $g$-stable lottery with distortion arbitrarily close to $1 + 2\lambda^\star$. It does not determine the minimax distortion over arbitrary randomized voting rules, nor the distortion achievable by selecting the most favorable $g$-stable lottery in each election.

\subsection{Technical Overview}
\label{sec:overview}

By the biased metric reduction of \citelink{Charikar and Ramakrishnan}{CharikarRamakrishnan22}~\cite{CharikarRamakrishnan22} (\cref{sec:preliminaries}), the worst-case consistent pseudometric for a lottery $D$ is determined by the distances $\rho(b) = d(a^\star, b)$ from an optimal candidate $a^\star \in C$ to the other candidates. Along voter $v$'s ranking, the running minimum $\lowrho_v(b) = \min_{c \preceq_v b} \rho(c)$ gives $d(v, b) - d(v, a^\star)$, while the largest drop $\Delta_v = \max_{b \in C}(\rho(b) - \lowrho_v(b))$ gives $2\,d(v, a^\star)$. Thus $D$ has distortion at most $1 + 2\lambda$ if and only if $\E_v[N_v(D)] \le \lambda\,\E_v[\Delta_v]$ for every choice of $a^\star$ and $\rho$, where $N_v(D) = \E_{b \sim D}[\lowrho_v(b)]$ (\cref{lem:biased-reduction}).

To write the stability conditions in terms of ranks, we center and rescale $g$ around its mean $\tau$ by setting $\phi(r) = c(\tau - g(r))$ and $J(r) = \int_0^r \phi(u)\,\mathrm{d}u$ for a constant $c > 0$, so that $J(0) = J(1) = 0$. On voter $v$'s rank line, an unsupported candidate at rank $r$ has centered activation $-\phi(r)$, while a block of candidates occupying $[u, w]$ has total centered activation $-\int_u^w \phi(r)\,\mathrm{d}r = J(u) - J(w)$. By \cref{lem:stationarity}, a lottery $D$ is $g$-stable when every supported candidate has expected centered activation zero across voters and every unsupported candidate has expected centered activation at most zero.

\paragraph{1. Three adversarial elections (\cref{sec:elections}).}
Because a $g$-stable lottery is defined implicitly, starting from an election and solving for $D$ would be unwieldy. Instead, we fix $D$ and $\rho$ first, and choose the voters' rankings so that $D$ is $g$-stable while making $\E_v[N_v]$ as large as possible relative to $\E_v[\Delta_v]$. Three families of elections constrain $\phi$ and $J$ across the rank line (\cref{fig:optimal-activation}(a)):
\begin{itemize}
  \item \textbf{Election 1 (Pointwise upper bound on $\phi$):} Suppose $a^\star \notin \supp(D)$ and every supported candidate has $\rho = 1$. A voter placing $a^\star$ at rank $r < 1$ has $N_v = r$ and $\Delta_v = 1$, whereas a voter placing $a^\star$ at the top ($r = 1$) has $N_v = 1$ and $\Delta_v = 0$. Balancing two ranks for $a^\star$ on opposite sides of $\tau$ gives $\phi(r) \le \lambda - r$ on $[0, 1]$ and $\phi(1) \le -1$ (\cref{lem:election-1}).
  \item \textbf{Election 2 (Two-class inequality for $J$):} Partition $\supp(D)$ into class $\calH$ ($\rho = 1$, mass $p$) and class $\calZ$ ($\rho = 0$, mass $1 - p$). Voters who rank all of $\calH$ below $\calZ$ have $N_v = p$ and $\Delta_v = 0$, with centered activation $-J(p)$ on $\calH$. Balancing them against voters who split $\calH$ into a bottom block of mass $q$ and a top block of mass $p - q$ around $\calZ$ gives $p\,[J(1 - p + q) - J(q)] \le (\lambda - q)J(p)$ (\cref{lem:election-2}).
  \item \textbf{Election 3 (Differential inequality at high ranks):} Partition $\supp(D)$ into class $\calH$ ($\rho = 2$, mass $p$) and class $\calM$ ($\rho = 1$, mass $1 - p$), with $a^\star \notin \supp(D)$. Mixing three voter types that place $a^\star$ at two low ranks $x, y < \lambda$ and one high rank $z > \lambda$ while placing $\calH$ at the top, on $[z - p, z]$, and at the bottom gives a linear differential inequality bounding the high-rank derivative $\phi(z) = J'(z)$ in terms of $J(z)$ and the low-rank values $\phi(x), \phi(y), J(p), J(1 - p)$ (\cref{lem:election-3}).
\end{itemize}

\begin{figure}[!t]
\centering
\begin{tikzpicture}[y=1.0cm,font=\small,
    axis/.style={-{Latex[length=4.5pt,width=3.2pt]},draw=ContentsHeading,line width=0.55pt},
    ptr/.style={-{Latex[length=3.5pt,width=2.5pt]},line width=0.45pt}]
  % Panel (a): Activation g*(r)
  \begin{scope}[shift={(0,0)}]
    \node[font=\fontsize{9.5}{11}\selectfont\bfseries,text=ContentsHeading] at (3.0,4.55) {(a) Optimal activation $g^\star(r)$};

    % Shaded regions 2, 1, 3
    \fill[AccentTeal!12] (0.55,0) rectangle (1.85,4.0);
    \fill[AccentBlue!12] (1.85,0) rectangle (4.30,4.0);
    \fill[AccentBrown!12] (4.30,0) rectangle (4.85,4.0);

    \node[font=\scriptsize\bfseries,text=AccentTeal] at (1.20,3.72) {Election 2};
    \node[font=\scriptsize\bfseries,text=AccentBlue] at (3.07,3.72) {Election 1};
    \node[font=\scriptsize\bfseries,text=AccentBrown] at (4.58,3.72) {3};

    % Grid and tau line
    \draw[ContentsHeading!30,line width=0.4pt] (0,4.0) -- (5.5,4.0);
    \draw[ContentsHeading!30,line width=0.4pt] (5.5,0) -- (5.5,4.0);
    \draw[densely dashed,ContentsHeading!55,line width=0.45pt] (0,1.23) -- (5.5,1.23);
    \node[left,font=\scriptsize,text=ContentsHeading] at (0,1.23) {$\tau$};
    \node[left,font=\scriptsize,text=ContentsHeading] at (0,0) {$0$};
    \node[left,font=\scriptsize,text=ContentsHeading] at (0,4.0) {$1$};

    % Sloped line encoding phi <= lambda - r
    \draw[densely dotted,AccentBlue!80,line width=0.8pt] (0,0.22) -- (5.5,2.15);

    % Curve g*(r)
    \draw[AccentBrown!90!black,line width=1.25pt]
      (0,0) -- (0.55,0.13)
      .. controls (1.10,0.27) and (1.45,0.42) .. (1.85,0.87)
      -- (4.30,1.73)
      .. controls (4.65,1.87) and (5.05,2.55) .. (5.5,4.0);

    % Ticks
    \draw[ContentsHeading,line width=0.5pt] (3.09,0) -- (3.09,-0.08)
      node[below,font=\scriptsize,text=ContentsHeading] {$\lambda^\star$};
    \draw[ContentsHeading,line width=0.5pt] (5.5,0) -- (5.5,-0.08)
      node[below,font=\scriptsize,text=ContentsHeading] {$1$};
    \node[below,font=\scriptsize,text=ContentsHeading] at (0,-0.08) {$0$};

    \draw[axis] (0,0) -- (5.85,0) node[right,text=black] {$r$};
    \draw[axis] (0,0) -- (0,4.25);
  \end{scope}

  % Panel (b): Primitive J*(r) and variational bounds
  \begin{scope}[shift={(7.8,0)}]
    \node[font=\fontsize{9.5}{11}\selectfont\bfseries,text=ContentsHeading] at (3.2,4.55) {(b) Primitive $J^\star(r)$ and variational bounds};

    % Upper bound C_{lambda,S}(r) at low ranks
    \draw[densely dashed,AccentTeal,line width=0.85pt]
      plot[domain=0.45:2.35,smooth,samples=25] (\x,{0.18 + 1.42*\x - 0.19*\x*\x});
    \node[font=\scriptsize,text=AccentTeal,anchor=west] at (0.12,3.42) {upper bound $C_{\lambda^\star,S}(r)$};
    \draw[ptr,AccentTeal] (1.35,3.22) to[out=-80,in=120] (1.45,1.85);

    % Middle quadratic Q_{lambda}(r) - S
    \draw[densely dotted,AccentBlue,line width=0.85pt]
      plot[domain=1.50:4.80,smooth,samples=30] (\x,{-0.22 + 1.82*\x - 0.295*\x*\x});
    \node[font=\scriptsize,text=AccentBlue] at (3.65,2.92) {$Q_{\lambda^\star}(r) - S$};

    % Two-class high lower bound T_{lambda,S}(z, q_*)
    \draw[densely dashed,AccentBrown,line width=0.85pt]
      plot[domain=4.25:5.15,smooth,samples=20] (\x,{2.18 - 0.34*(\x-4.25) - 0.95*(\x-4.25)*(\x-4.25)});
    \node[font=\scriptsize,text=AccentBrown,anchor=east] at (4.55,0.82) {lower bound $T_{\lambda^\star,S}(z, q_\star)$};
    \draw[ptr,AccentBrown] (4.10,1.02) to[out=25,in=-135] (4.68,1.65);

    % Actual primitive J*(r)
    \draw[AccentBlue!90!black,line width=1.25pt]
      (0,0) .. controls (0.55,0.72) and (1.10,1.38) .. (1.65,1.86)
      .. controls (1.88,2.05) and (2.05,2.16) .. (2.20,2.24)
      plot[domain=2.20:4.30,smooth,samples=20] (\x,{-0.22 + 1.82*\x - 0.295*\x*\x})
      .. controls (4.55,1.98) and (4.95,1.35) .. (5.5,0);

    % Contact points c and z_*
    \fill[AccentTeal,draw=white,line width=0.5pt] (1.65,1.86) circle (2.1pt);
    \fill[AccentBrown,draw=white,line width=0.5pt] (4.72,1.72) circle (2.1pt);

    \draw[ContentsHeading,line width=0.5pt] (1.65,0) -- (1.65,-0.08)
      node[below,font=\scriptsize,text=AccentTeal] {$c$};
    \draw[ContentsHeading,line width=0.5pt] (3.09,0) -- (3.09,-0.08)
      node[below,font=\scriptsize,text=ContentsHeading] {$\lambda^\star$};
    \draw[ContentsHeading,line width=0.5pt] (4.72,0) -- (4.72,-0.08)
      node[below,font=\scriptsize,text=AccentBrown] {$z_\star$};
    \draw[ContentsHeading,line width=0.5pt] (5.5,0) -- (5.5,-0.08)
      node[below,font=\scriptsize,text=ContentsHeading] {$1$};
    \node[below,font=\scriptsize,text=ContentsHeading] at (0,-0.08) {$0$};

    \draw[axis] (0,0) -- (5.85,0) node[right,text=black] {$r$};
    \draw[axis] (0,0) -- (0,4.25);
  \end{scope}
\end{tikzpicture}
\caption{Structure of the optimal activation $g^\star$ and its primitive $J^\star(r) = \int_0^r \phi^\star(u)\,\mathrm{d}u$. (a)~In the middle region, $g^\star$ runs along the dotted line encoding $\phi \le \lambda^\star - r$ (Election~1); at low ranks, $g^\star$ starts with an affine piece and then follows $J^\star = C_{\lambda^\star,S}$ up to $c$ (Election~2); at high ranks, $g^\star$ follows the differential inequality of Election~3 before rising steeply to $g^\star(1) = 1$. (b)~The high-rank continuation $Y(z)$ of $J^\star$ touches the two-class lower bound $T_{\lambda^\star,S}(z, q)$ tangentially at an interior pair $(z_\star, q_\star)$.}
\label{fig:optimal-activation}
\end{figure}

\paragraph{2. The variational problem defining $\lambda^\star$ (\cref{sec:variational}).}
When $\phi(r) = \lambda - r$ is tight across an interval of middle ranks, integrating gives $J(r) = Q_\lambda(r) - S$, where $Q_\lambda(r) = \lambda r - r^2/2$ and $S$ is the total slack $\int (\lambda - u - \phi(u))\,\mathrm{d}u$ accumulated at smaller ranks. For low ranks $r$ whose complement $1 - r$ lies in this middle interval, the $q = 0$ two-class inequality $(1 - r)J(r) \le \lambda J(1 - r)$ bounds $J(r)$ from above by a rational curve $C_{\lambda, S}(r)$ (\cref{fig:optimal-activation}(b)).

Once $S$ and the low-rank function $\phi$ are chosen, the differential inequality of Election~3 bounds $J'(z)$ from above at high ranks $z$; taking the infimum of these upper bounds over all tested triples $(p, x, y)$ defines a maximal continuation $Y(z)$ by an ordinary differential equation. Meanwhile, the two-class inequality at high mass $z$ and split $q$ bounds $J(z)$ from below by an explicit function $T_{\lambda, S}(z, q)$. The optimal parameter $\lambda^\star$ is the smallest $\lambda$ for which some low-rank pair $(S, \phi)$ produces a continuation $Y$ that stays above every lower bound $T_{\lambda, S}(z, q)$.

Any activation achieving distortion at most $1 + 2\lambda$ can be modified into a feasible pair $(S, \phi)$ for this variational problem (\cref{prop:reduction}). Weak-* compactness in $L^\infty$ shows that $\lambda^\star$ is attained and that at $\lambda^\star$, $Y(z)$ touches at least one lower bound $T_{\lambda^\star, S}(z_\star, q_\star)$ (\cref{prop:attainment}; \cref{fig:optimal-activation}(b)). Local perturbations at $\lambda^\star$ then show that almost every low rank between the last contact $c$ with $C_{\lambda^\star, S}$ and the line $\lambda^\star - r$ is active in the Election~3 differential equation (\cref{prop:shape}). Varying $x$ or $y$ at an active triple produces supporting lines to $\phi$ that rule out jumps, prove Lipschitz continuity, and show that the high-rank derivative $Y'(z)$ is concave with slope strictly below $-1$ after leaving $\lambda^\star - z$.

\paragraph{3. The matching upper bound from supporting lines (\cref{sec:upper-bound}).}
To prove the matching upper bound, we complete a minimizing function $\phi$ at $\lambda^\star$ to $[0, 1]$ with an affine piece near $0$ and a quadratic piece near $1$ matching the tangential contact at $(z_\star, q_\star)$. Following \citelink{Cai, Charikar, Frank, Hastings, Ramakrishnan, Wang, and Ye}{CaiEtAl26b}~\cite{CaiEtAl26b}, we slice $N_v(D)$ by distance threshold $t \ge 0$, offset each slice using the stability of supported candidates and of $a^\star$, and telescope the offset slice widths across thresholds separated by $\Delta_v$ using a potential function (\cref{lem:verifier}). At $\lambda^\star$, the supporting lines from \cref{prop:shape} directly supply the slice weights and corrections needed for this telescoping argument: substituting the maximizing low-rank positions $x$ and $y$ into the Election~3 differential inequality reduces it to the high-rank verifier inequality at $q = 0$. Concavity and monotonicity extend the verification to all masses $p$ and splits $q > 0$ (\cref{app:upper-all}), while a second-variation comparison with an explicit reference function (\cref{sec:reference-comparison,app:numerical}) certifies $\lambda^\star \in (0.56232844,\, 0.56232847)$.

\section{Preliminaries: \texorpdfstring{$g$}{g}-Stable Lotteries and Biased Metrics}
\label{sec:preliminaries}

An \emph{election} (or preference profile) is a tuple $\calE = (V, C, \{\succ_v\}_{v \in V})$, where $V$ is a finite nonempty set of voters, $C$ is a finite nonempty set of candidates, and each voter $v \in V$ has a strict linear order $\succ_v$ over $C$. We write $a \succ_v b$ when voter $v$ strictly prefers candidate $a$ to candidate $b$, and $a \succeq_v b$ (or $b \preceq_v a$) when $a \succ_v b$ or $a = b$. We write $v \sim V$ for a voter sampled uniformly at random from $V$; when voters with identical rankings are grouped into voter types, a type's weight $w_v$ is its fraction of the electorate. A \emph{lottery} is a probability distribution $D$ over $C$, with probability mass $D(a)$ on candidate $a \in C$ and support $\supp(D) = \{a \in C : D(a) > 0\}$.

We treat the voters and candidates $V \cup C$ as points in a pseudometric space with distance $d(\cdot, \cdot)$ satisfying nonnegativity ($d(x, y) \ge 0$ with $d(x, x) = 0$), symmetry ($d(x, y) = d(y, x)$), and the triangle inequality ($d(x, y) + d(y, z) \ge d(x, z)$). A pseudometric $d$ is \emph{consistent} with $\calE$ if $a \succ_v b$ implies $d(v, a) \le d(v, b)$ for every voter $v \in V$ and candidates $a, b \in C$. The social cost of a candidate $a$ and of a lottery $D$ are
\[
  \cost_d(a) \;:=\; \E_{v \sim V}\bigl[d(v, a)\bigr]
  \qquad\text{and}\qquad
  \cost_d(D) \;:=\; \E_{b \sim D}\bigl[\cost_d(b)\bigr],
\]
where we omit the subscript $d$ when the pseudometric is clear from context. Fix an optimal candidate $a^\star \in \arg\min_{a \in C} \cost_d(a)$. The \emph{distortion} of a lottery $D$ in $\calE$ is $\sup_d \cost_d(D) / \cost_d(a^\star)$ over all consistent pseudometrics $d$ with $\cost_d(a^\star) > 0$.

\subsection{The Rank Line and Stationarity of \texorpdfstring{$g$}{g}-Stable Lotteries}
\label{sec:rank-line}

Following \citelink{Charikar, Lassota, Ramakrishnan, Vetta, and Wang}{CharikarEtAl25}~\cite{CharikarEtAl25} and \citelink{Cai, Charikar, Frank, Hastings, Ramakrishnan, Wang, and Ye}{CaiEtAl26b}~\cite{CaiEtAl26b}, we visualize each voter's ranking with respect to a lottery $D$ on the unit interval $[0, 1]$, called the \emph{rank line} of voter $v$ (\cref{fig:rank-line-boxes}(b)). Each candidate $a \in C$ is assigned an interval $I_v(a) := [\rank_v(a; D),\, \rank_v(a; D) + D(a))$ of width $D(a)$, arranged in increasing order of $v$'s preference from worst ($0$) to best ($1$), so that the left endpoint of $I_v(a)$ is
\[
  \rank_v(a; D) \;:=\; \Pr_{b \sim D}[a \succ_v b] \;=\; \sum_{b:\, a \succ_v b} D(b).
\]
Sampling a candidate $b \sim D$ is equivalent to drawing a uniform rank $r \sim \operatorname{Unif}(0, 1)$ and outputting the candidate whose interval $I_v(b)$ contains $r$. We adopt the continuum-of-copies convention~\cite{CaiEtAl26b}: each candidate $a \in C$ corresponds to a continuum of copies indexed by the points of $I_v(a)$, with ties broken uniformly. Under a continuous activation function $g \colon [0, 1] \to [0, 1]$, the expected activation of a random copy of $a$ on voter $v$'s rank line is its average over $I_v(a)$:
\[
  \psi_v(a; D) \;:=\; \int_0^1 g\bigl(\rank_v(a; D) + u\,D(a)\bigr)\,\mathrm{d}u.
\]
When $D(a) = 0$, this equals $g(\rank_v(a; D))$; when $D(a) > 0$, it equals $\frac{1}{D(a)}\int_{I_v(a)} g(r)\,\mathrm{d}r$. Drawing $b \sim D$ and then a uniform copy in $I_v(b)$ produces a uniform rank $r \sim \operatorname{Unif}(0, 1)$, whose expected activation is
\[
  \tau \;:=\; \int_0^1 g(r)\,\mathrm{d}r.
\]

\begin{definition}[$g$-stable lottery]
\label{def:g-stable}
For a continuous activation function $g \colon [0, 1] \to [0, 1]$ with mean $\tau = \int_0^1 g(r)\,\mathrm{d}r$, a lottery $D$ is \emph{$g$-stable} if
\[
  \E_{v \sim V}\bigl[\psi_v(a; D)\bigr] \;\le\; \tau
  \qquad \text{for all } a \in C.
\]
We write $M(g)$ for the supremum distortion of a $g$-stable lottery over all finite elections, and $\calG$ for the class of continuous activation functions $g \colon [0, 1] \to [0, 1]$.
\end{definition}

In game-theoretic terms, if an attacker chooses a candidate $a$ and a defender chooses a lottery $D$, a random voter $v \sim V$ and a random copy of $a$ in $I_v(a)$ give the attacker payoff $\E_v[\psi_v(a; D)]$. Drawing $a \sim D$ gives payoff $\tau$ against every $D$, and a $g$-stable lottery $D$ ensures that no attacking candidate gets payoff exceeding $\tau$. Because each interval $I_v(a)$ depends nonlinearly on $D$, we prove existence below using Brouwer's fixed-point theorem. Adding a constant to $g$ or multiplying $g$ by a positive constant preserves the set of $g$-stable lotteries~\cite{CharikarEtAl25,CaiEtAl26b}. Neither the definition of $g$-stable lotteries nor our lower bound requires $g$ to be monotone; our optimal activation $g^\star \in \calG$ for the matching upper bound is continuous and strictly increasing.

\begin{lemma}[Existence and stationarity of $g$-stable lotteries]
\label{lem:stationarity}
For every continuous activation function $g \colon [0, 1] \to [0, 1]$, every finite election has at least one $g$-stable lottery $D$. For every voter $v$ and every lottery $D$, we have $\sum_{a \in C} D(a)\,\psi_v(a; D) = \tau$, and thus every supported candidate $a \in \supp(D)$ in a $g$-stable lottery $D$ satisfies
\[
  \E_{v \sim V}\bigl[\psi_v(a; D)\bigr] \;=\; \tau.
\]
When $\tau > 0$, if we sample a voter $v \sim V$ uniformly and a rank $r \in [0, 1]$ with density $g(r)/\tau$, the candidate whose interval $I_v(a)$ contains $r$ has distribution $D$.
\end{lemma}

The identity $\sum_{a \in C} D(a)\,\psi_v(a; D) = \tau$ simply integrates $g$ over the partition $\{I_v(a)\}_{a \in C}$ of $[0, 1]$. Taking expectations over $v \sim V$, stability bounds each term $\E_v[\psi_v(a; D)]$ by $\tau$, so equality of their $D$-weighted average implies $\E_{v \sim V}[\psi_v(a; D)] = \tau$ whenever $D(a) > 0$. Existence follows from Brouwer's fixed-point theorem and is proved in \cref{app:existence}.

\subsection{Biased Metrics and the Running-Minimum Reduction}
\label{sec:biased-metrics}

To analyze the distortion of a lottery $D$ in a fixed election $\calE$, we use the biased metric framework of \citelink{Charikar and Ramakrishnan}{CharikarRamakrishnan22}~\cite{CharikarRamakrishnan22}, which characterizes the worst-case consistent pseudometrics in terms of the distances $\rho(b) := d(a^\star, b)$ from an optimal candidate $a^\star$ to the other candidates (\cref{fig:rank-line-boxes}(a)).

For any consistent pseudometric $d$ and any candidates $b \succeq_v c$ (meaning voter $v$ ranks $b$ at or above $c$), consistency and the triangle inequality give
\[
  d(v, b) \;\le\; d(v, c) \;\le\; d(v, a^\star) + \rho(c),
  \qquad
  \rho(b) - \rho(c) \;\le\; d(a^\star, v) + d(v, b) - \bigl(d(v, c) - d(a^\star, v)\bigr) \;\le\; 2\,d(v, a^\star).
\]
Minimizing the first bound over all $c \preceq_v b$ and maximizing the second over all $b \succeq_v c$ motivates the following definition~\cite{CharikarRamakrishnan22,CaiEtAl26b}.

\begin{definition}[Biased metric]
\label{def:biased-metric}
Fix a candidate $a^\star \in C$ and a function $\rho \colon C \to \R_{\ge 0}$ with $\rho(a^\star) = 0$. For each voter $v \in V$ and candidate $b \in C$, define the running minimum and maximum gap
\[
  \lowrho_v(b) \;:=\; \min_{c \preceq_v b} \rho(c),
  \qquad
  \Delta_v \;:=\; \max_{b \succeq_v c}\bigl(\rho(b) - \rho(c)\bigr) \;=\; \max_{b \in C}\bigl(\rho(b) - \lowrho_v(b)\bigr).
\]
The \emph{biased pseudometric} $d$ corresponding to $\rho$ sets $d(b, c) = \max\{\rho(b), \rho(c)\}$ for distinct candidates $b \ne c$ and, for each voter $v \in V$ and candidate $b \in C$,
\[
  d(v, a^\star) \;=\; \frac{1}{2}\Delta_v,
  \qquad
  d(v, b) - d(v, a^\star) \;=\; \lowrho_v(b).
\]
\end{definition}

For any lottery $D$ and any consistent pseudometric $d$ with $d(a^\star, b) = \rho(b)$, averaging the two inequalities $d(v, b) - d(v, a^\star) \le \lowrho_v(b)$ and $\Delta_v \le 2\,d(v, a^\star)$ over $b \sim D$ and $v \sim V$ gives
\[
  \cost_d(D) - \cost_d(a^\star) \;\le\; \E_{v \sim V}\bigl[N_v(D)\bigr]
  \qquad\text{and}\qquad
  \E_{v \sim V}[\Delta_v] \;\le\; 2\cost_d(a^\star),
\]
where
\[
  N_v(D) \;:=\; \E_{b \sim D}\bigl[\lowrho_v(b)\bigr] \;=\; \sum_{b \in C} D(b)\,\lowrho_v(b)
\]
is voter $v$'s expected distance difference $d(v, b) - d(v, a^\star)$ under $b \sim D$ in the biased pseudometric. Conversely, the biased pseudometric (\cref{def:biased-metric}) is itself consistent with $\calE$ and attains both bounds with equality: $\cost_d(a^\star) = \frac{1}{2}\E_v[\Delta_v]$ and $\cost_d(D) - \cost_d(a^\star) = \E_v[N_v(D)]$. We write $N_v$ for $N_v(D)$ when $D$ is fixed.

\begin{figure}[!t]
\centering
\begin{tikzpicture}[y=1.0cm,font=\small,
    axis/.style={-{Latex[length=4.5pt,width=3.2pt]},draw=ContentsHeading,line width=0.55pt},
    ptr/.style={-{Latex[length=3.5pt,width=2.5pt]},line width=0.55pt},
    dptr/.style={{Latex[length=3.2pt,width=2.3pt]}-{Latex[length=3.2pt,width=2.3pt]},line width=0.55pt}]
  % Panel (a): Biased metric geometry
  \begin{scope}[shift={(0,0.2)}]
    \node[font=\fontsize{9.5}{11}\selectfont\bfseries,text=ContentsHeading] at (2.5,3.45) {(a) Biased metric distances from $\rho$};

    \coordinate (astar) at (3.25,1.55);
    \coordinate (v) at (1.50,2.25);
    \coordinate (b) at (3.45,2.95);
    \coordinate (c) at (0.45,0.45);
    \coordinate (cp) at (4.55,0.45);

    % Solid segments from a*
    \draw[ContentsHeading!85,line width=0.8pt] (astar) -- (b)
      node[midway,right=2pt,font=\scriptsize,text=ContentsHeading] {$\rho(b)$};
    \draw[ContentsHeading!85,line width=0.8pt] (astar) -- (c)
      node[midway,above left=-1pt,font=\scriptsize,text=ContentsHeading] {$\rho(c)$};
    \draw[ContentsHeading!85,line width=0.8pt] (astar) -- (cp)
      node[midway,above right=-1pt,font=\scriptsize,text=ContentsHeading] {$\rho(c')$};

    % Dashed voter distances
    \draw[densely dashed,AccentTeal!90!black,line width=0.75pt,-{Latex[length=3.5pt,width=2.5pt]}]
      (v) -- ($(v)!0.88!(astar)$)
      node[midway,above=1pt,sloped,font=\scriptsize,text=AccentTeal!90!black] {shorten};
    \draw[densely dashed,AccentBrown!90!black,line width=0.75pt,-{Latex[length=3.5pt,width=2.5pt]}]
      (v) -- ($(v)!0.90!(b)$)
      node[midway,above=1pt,sloped,font=\scriptsize,text=AccentBrown!90!black] {lengthen};
    \draw[densely dashed,AccentBrown!90!black,line width=0.75pt,-{Latex[length=3.5pt,width=2.5pt]}]
      (v) -- ($(v)!0.92!(c)$)
      node[midway,above=1pt,sloped,font=\scriptsize,text=AccentBrown!90!black] {lengthen};
    \draw[densely dashed,AccentBrown!90!black,line width=0.75pt,-{Latex[length=3.5pt,width=2.5pt]}]
      (v) -- ($(v)!0.94!(cp)$)
      node[pos=0.68,below=1pt,sloped,font=\scriptsize,text=AccentBrown!90!black] {lengthen};

    % Nodes
    \fill[AccentTeal,draw=white,line width=0.6pt] (v) circle (2.5pt)
      node[above left=1pt,font=\footnotesize\bfseries,text=AccentTeal] {$v$};
    \fill[AccentBrown!90!black,draw=white,line width=0.6pt] (astar) circle (2.5pt)
      node[right=2pt,font=\footnotesize\bfseries,text=AccentBrown!90!black] {$a^\star$};
    \fill[AccentBlue,draw=white,line width=0.6pt] (b) circle (2.3pt)
      node[above=1pt,font=\footnotesize\bfseries,text=AccentBlue] {$b$};
    \fill[AccentBlue,draw=white,line width=0.6pt] (c) circle (2.3pt)
      node[left=1pt,font=\footnotesize\bfseries,text=AccentBlue] {$c$};
    \fill[AccentBlue,draw=white,line width=0.6pt] (cp) circle (2.3pt)
      node[right=1pt,font=\footnotesize\bfseries,text=AccentBlue] {$c'$};

    \node[font=\scriptsize,text=ContentsHeading,align=center] at (2.50,-0.15)
      {$d(v, a^\star) = \tfrac{1}{2}\Delta_v$, \quad $d(v, b) - d(v, a^\star) = \lowrho_v(b)$};
  \end{scope}

  % Panel (b): Step function rho_v and running minimum on the rank line
  \begin{scope}[shift={(6.4,0)}]
    \node[font=\fontsize{9.5}{11}\selectfont\bfseries,text=ContentsHeading] at (4.1,3.65) {(b) $\rho_v$ and running minimum $\lowrho_v$ on $v$'s rank line};

    % Shaded running minimum area
    \fill[pattern=north east lines,pattern color=AccentBrown!55]
      (0,0) -- (0,2.50) -- (1.25,2.50) -- (1.25,1.75) -- (3.30,1.75) -- (3.30,0.98) -- (5.85,0.98) -- (5.85,0.20) -- (7.90,0.20) -- (7.90,0) -- cycle;
    \fill[AccentBrown!10,opacity=0.6]
      (0,0) -- (0,2.50) -- (1.25,2.50) -- (1.25,1.75) -- (3.30,1.75) -- (3.30,0.98) -- (5.85,0.98) -- (5.85,0.20) -- (7.90,0.20) -- (7.90,0) -- cycle;

    % Boxes for candidates on the rank line
    \draw[ContentsHeading!80,line width=0.6pt] (0,0) rectangle (1.25,2.50);
    \draw[ContentsHeading!80,line width=0.6pt] (1.25,0) rectangle (2.25,1.75);
    \draw[ContentsHeading!80,line width=0.6pt] (2.25,0) rectangle (3.30,2.88);
    \draw[ContentsHeading!80,line width=0.6pt] (3.30,0) rectangle (4.05,0.98);
    \draw[ContentsHeading!80,line width=0.6pt] (4.05,0) rectangle (5.30,2.35);
    \draw[ContentsHeading!80,line width=0.6pt] (5.30,0) rectangle (5.85,2.62);
    \draw[ContentsHeading!80,line width=0.6pt] (5.85,0) rectangle (6.90,0.20);
    \draw[ContentsHeading!80,line width=0.6pt] (6.90,0) rectangle (7.90,1.60);

    % Running minimum line
    \draw[AccentBrown!90!black,line width=1.25pt]
      (0,2.50) -- (1.25,2.50) -- (1.25,1.75) -- (3.30,1.75) -- (3.30,0.98) -- (5.85,0.98) -- (5.85,0.20) -- (7.90,0.20);

    % Top profile rho_v(r)
    \draw[AccentBlue!90!black,line width=1.15pt]
      (0,2.50) -- (1.25,2.50) -- (1.25,1.75) -- (2.25,1.75) -- (2.25,2.88) -- (3.30,2.88) -- (3.30,0.98) -- (4.05,0.98) -- (4.05,2.35) -- (5.30,2.35) -- (5.30,2.62) -- (5.85,2.62) -- (5.85,0.20) -- (6.90,0.20) -- (6.90,1.60) -- (7.90,1.60);

    % Colored rank intervals along bottom axis
    \fill[AccentTeal!75] (0,-0.04) rectangle (1.25,0.04);
    \fill[AccentBlue!70] (1.25,-0.04) rectangle (2.25,0.04);
    \fill[AccentBrown!75] (2.25,-0.04) rectangle (3.30,0.04);
    \fill[AccentTeal!55] (3.30,-0.04) rectangle (4.05,0.04);
    \fill[AccentBlue!55] (4.05,-0.04) rectangle (5.30,0.04);
    \fill[AccentBrown!55] (5.30,-0.04) rectangle (5.85,0.04);
    \fill[AccentTeal!85] (5.85,-0.04) rectangle (6.90,0.04);
    \fill[AccentBlue!85] (6.90,-0.04) rectangle (7.90,0.04);

    % Annotations for rho(b), lowrho_v(b), and Delta_v
    \draw[dptr,AccentBlue] (4.55,0.05) -- (4.55,2.33)
      node[midway,left=1pt,font=\scriptsize,text=AccentBlue,fill=white,inner sep=1pt] {$\rho(b)$};
    \draw[dptr,AccentBrown!90!black] (4.95,0.05) -- (4.95,0.96)
      node[midway,right=1pt,font=\scriptsize,text=AccentBrown!90!black,fill=white,inner sep=1pt] {$\lowrho_v(b)$};
    \draw[dptr,AccentTeal!90!black] (5.58,1.00) -- (5.58,2.60)
      node[midway,right=2pt,font=\scriptsize\bfseries,text=AccentTeal!90!black,fill=white,inner sep=1pt] {$\Delta_v$};

    \node[font=\scriptsize,text=AccentBrown!85!black,fill=white,inner sep=1.5pt] at (1.85,0.80)
      {$N_v(D) = \E_{b \sim D}[\lowrho_v(b)]$};
    \node[above,font=\scriptsize,text=AccentBlue!90!black] at (2.78,2.90) {$\rho_v(\cdot)$};

    \node[below,font=\scriptsize,text=ContentsHeading] at (0,-0.08) {$0$};
    \node[below,font=\scriptsize,text=ContentsHeading] at (4.68,-0.08) {$b$};
    \node[below,font=\scriptsize,text=ContentsHeading] at (7.90,-0.08) {$1$};

    \draw[axis] (0,0) -- (8.30,0) node[right,text=black] {$r$};
    \draw[axis] (0,0) -- (0,3.25);
  \end{scope}
\end{tikzpicture}
\caption{Biased metrics and the rank-line representation~\cite{CharikarRamakrishnan22,CaiEtAl26b}. (a)~Fixing $\rho(b) = d(a^\star, b)$, the biased pseudometric minimizes $d(v, a^\star) = \frac{1}{2}\Delta_v$ and maximizes $d(v, b) - d(v, a^\star) = \lowrho_v(b)$ subject to consistency with voter $v$'s ranking. (b)~On voter $v$'s rank line $[0, 1]$, each candidate $b \in C$ corresponds to a box of width $D(b)$ and height $\rho(b)$. The hatched area under the running minimum $\lowrho_v$ equals $N_v(D) = \E_{b \sim D}[\lowrho_v(b)]$, and the largest vertical gap between $\rho_v$ and $\lowrho_v$ equals $\Delta_v = \max_{b \in C}(\rho(b) - \lowrho_v(b))$.}
\label{fig:rank-line-boxes}
\end{figure}

\begin{lemma}[Biased metric reduction]
\label{lem:biased-reduction}
For $\lambda \ge 0$, a lottery $D$ has distortion at most $1 + 2\lambda$ if and only if
\begin{equation}
  \E_{v \sim V}\bigl[N_v(D)\bigr] \;\le\; \lambda\,\E_{v \sim V}[\Delta_v]
  \label{eq:height-inequality}
\end{equation}
for every candidate $a^\star \in C$ and every function $\rho \colon C \to \R_{\ge 0}$ with $\rho(a^\star) = 0$.
\end{lemma}

We give the proof of \cref{lem:biased-reduction}~\cite{CharikarRamakrishnan22,CaiEtAl26b} in \cref{app:existence}. Placing a box of width $D(b)$ and height $\rho(b)$ over each interval $I_v(b)$ on voter $v$'s rank line traces out a step function $\rho_v \colon [0, 1] \to \R_{\ge 0}$ whose running minimum is $\lowrho_v$ (\cref{fig:rank-line-boxes}(b)). Both the running minimum $\lowrho_v(b) = \min_{c \preceq_v b} \rho(c)$ and the maximum gap $\Delta_v = \max_{b \in C}(\rho(b) - \lowrho_v(b))$ include candidates outside $\supp(D)$ whose intervals on the rank line have width zero.

\section{Three Elections and Necessary Inequalities}
\label{sec:elections}

Suppose an activation $g \in \calG$ guarantees distortion at most $1 + 2\lambda < \infty$ in every finite election. Because a $g$-stable lottery is defined implicitly, starting from an election and solving for its $g$-stable lottery would be unwieldy. Instead, we work in reverse~\cite{CaiEtAl26b}: we fix a lottery $D$ and distances $\rho \colon C \to \R_{\ge 0}$ with $\rho(a^\star) = 0$ first, and then design the voters' rankings so that $D$ is $g$-stable while making $\E_v[N_v]$ as large as possible relative to $\E_v[\Delta_v]$.\footnote{In this section we describe the constructions using real voter-type weights and continuum candidate classes. In \cref{app:normalization,app:finite-two-three,app:equal-weights}, we prove that any strict violation of the resulting inequalities is realized by a finite election with equally weighted voters and an exactly $g$-stable lottery.}

Because stability compares $g(r)$ against its mean $\tau = \int_0^1 g(r)\,\mathrm{d}r$, we center $g$ around $\tau$. For a scaling factor $c > 0$ chosen in \cref{lem:election-1} below, define
\begin{equation}
  \phi(r) \;:=\; c\bigl(\tau - g(r)\bigr),
  \qquad
  J(r) \;:=\; \int_0^r \phi(u)\,\mathrm{d}u.
  \label{eq:score-primitive-def}
\end{equation}
Since $g \in \calG$ is continuous, $\phi$ is continuous and integrates to zero on $[0, 1]$, so $J(0) = J(1) = 0$; we show in \cref{lem:election-1} that $M(g) \le 1 + 2\lambda < \infty$ also implies $J(r) > 0$ for all $r \in (0, 1)$.

We refer to $-\phi(r) = c\bigl(g(r) - \tau\bigr)$ as the \emph{centered activation} at rank $r$. On any voter's rank line, a candidate outside $\supp(D)$ at rank $r$ has centered activation $-\phi(r)$, while a block of candidates occupying the rank interval $[u, w]$ has centered activation
\[
  -\int_u^w \phi(r)\,\mathrm{d}r \;=\; J(u) - J(w).
\]
In particular, a block $[0, p]$ at the bottom of a voter's ranking has centered activation $-J(p) < 0$, a block $[1 - p, 1]$ at the top has centered activation $J(1 - p) > 0$, and the entire rank line $[0, 1]$ has centered activation $J(0) - J(1) = 0$.

Throughout the three constructions below, we group candidates in $\supp(D)$ that share the same $\rho$-value into a class. To treat all candidates within a class alike, we split the class's probability mass under $D$ equally among $m$ candidates and replicate each voter's ranking across all $m!$ permutations of the class. Every candidate in the class then receives an equal share of the class's centered activation, so by \cref{lem:stationarity}, $D$ is $g$-stable if and only if each class in $\supp(D)$ has expected centered activation zero and each candidate outside $\supp(D)$ has expected centered activation at most zero.

\subsection{Election 1: A Pointwise Bound on \texorpdfstring{$\phi$}{phi}}
\label{sec:election-1}

Our first construction tests how $g$ responds when $a^\star \notin \supp(D)$ has $\rho(a^\star) = 0$ and every candidate $b \in \supp(D)$ has $\rho(b) = 1$. Since $\supp(D)$ forms a single class occupying the entire rank line $[0, 1]$, every candidate in $\supp(D)$ automatically has expected centered activation zero. Thus $D$ is $g$-stable whenever the distribution of $a^\star$'s rank across voters gives $a^\star$ expected centered activation $\E_v[-\phi(r_v)] \le 0$.

Consider a voter $v$ who places $a^\star$ at rank $r \in [0, 1]$. Every candidate ranked below $a^\star$ has $\lowrho_v(b) = 1$, while every candidate ranked above $a^\star$ has $\lowrho_v(b) = \rho(a^\star) = 0$:
\begin{itemize}
  \item If $r < 1$, the candidates below $a^\star$ have total mass $r$, so $N_v = r$, while the candidates in $\supp(D)$ ranked above $a^\star$ have $\rho(b) - \lowrho_v(b) = 1$, so $\Delta_v = 1$.
  \item If $r = 1$ (meaning $v$ ranks $a^\star$ above all of $\supp(D)$), then $N_v = 1$, but $\lowrho_v(b) = \rho(b) = 1$ for every $b \in \supp(D)$ and $\lowrho_v(a^\star) = \rho(a^\star) = 0$, so $\Delta_v = 0$.
\end{itemize}
Therefore, a voter who places $a^\star$ at rank $r$ contributes $f_\lambda(r)$ to $N_v - \lambda\Delta_v$, where
\[
  f_\lambda(r) \;:=\;
  \begin{cases}
    r - \lambda, & 0 \le r < 1, \\
    1, & r = 1.
  \end{cases}
\]
By \cref{lem:biased-reduction}, every distribution of ranks $r_v$ for $a^\star$ with $\E_v[g(r_v) - \tau] \le 0$ must satisfy $\E_v[f_\lambda(r_v)] \le 0$. Mixing two ranks whose activations lie on opposite sides of $\tau$ turns this expectation condition into a pointwise bound for a suitable scaling factor $c > 0$ (\cref{app:normalization}):

\begin{lemma}[Election 1 bound]
\label{lem:election-1}
If $M(g) \le 1 + 2\lambda$, there is a constant $c > 0$ such that $\phi(r) = c(\tau - g(r))$ and $J(r) = \int_0^r \phi(u)\,\mathrm{d}u$ satisfy
\begin{equation}
  \begin{gathered}
    J(0) = J(1) = 0, \qquad J(r) > 0 \quad (0 < r < 1), \\
    \phi(r) \le \lambda - r \quad (0 \le r \le 1), \qquad \phi(1) \le -1.
  \end{gathered}
  \label{eq:election-1-bound}
\end{equation}
\end{lemma}

The jump from $\lambda - 1$ to $-1$ at $r = 1$ reflects the voters who rank $a^\star$ at the very top: because they contribute $N_v = 1$ while paying nothing toward $\Delta_v = 0$, the activation $g(1)$ at rank $1$ must be large enough relative to $\tau$ to cap the fraction of such voters. Writing
\[
  Q_\lambda(r) \;:=\; \lambda r - \frac{r^2}{2},
  \qquad
  \sigma(r) \;:=\; \lambda - r - \phi(r),
\]
for the primitive of $\lambda - r$ and the nonnegative gap between $\lambda - r$ and $\phi(r)$, \cref{eq:election-1-bound} becomes
\begin{equation}
  \sigma(r) \;\ge\; 0,
  \qquad
  \int_0^1 \sigma(u)\,\mathrm{d}u \;=\; Q_\lambda(1) - J(1) \;=\; \lambda - \frac{1}{2},
  \qquad
  J(r) \;=\; Q_\lambda(r) - \int_0^r \sigma(u)\,\mathrm{d}u.
  \label{eq:sigma-relations}
\end{equation}

\begin{figure}[!t]
\centering
\begin{tikzpicture}[y=1.0cm,font=\small,
    axis/.style={-{Latex[length=4.5pt,width=3.2pt]},draw=ContentsHeading,line width=0.55pt}]
  % Panel (a): Election 2 voter types
  \begin{scope}[shift={(0,0)}]
    \node[font=\fontsize{9.5}{11}\selectfont\bfseries,text=ContentsHeading] at (3.1,5.25) {(a) Election 2: two candidate classes $\calH$ and $\calZ$};

    % Top sub-diagram: Type 1 (H_q, Z_{1-p}, H_{p-q})
    \begin{scope}[shift={(0,2.75)}]
      \node[right,font=\scriptsize\bfseries,text=ContentsHeading] at (0,1.65) {Type 1: $(\calH_q,\; \calZ_{1-p},\; \calH_{p-q})$ \quad $[N_v = q,\; \Delta_v = 1]$};
      % Hatched running minimum on [0, q]
      \fill[pattern=north east lines,pattern color=AccentBrown!55] (0,0) rectangle (1.35,1.15);
      \fill[AccentBrown!10,opacity=0.6] (0,0) rectangle (1.35,1.15);
      % Block H_q: [0, 1.35], height 1.15
      \draw[AccentTeal!85!black,line width=0.7pt] (0,0) rectangle (1.35,1.15);
      \node[font=\scriptsize,text=AccentTeal!90!black] at (0.67,0.58) {$\calH_q$};
      % Block Z_{1-p}: [1.35, 4.30], height 0
      \fill[AccentBlue!25] (1.35,0) rectangle (4.30,0.12);
      \node[above,font=\scriptsize,text=AccentBlue] at (2.82,0.10) {$\calZ_{1-p}$ ($\rho = 0$)};
      % Block H_{p-q}: [4.30, 5.80], height 1.15
      \draw[AccentTeal!85!black,line width=0.7pt] (4.30,0) rectangle (5.80,1.15);
      \node[font=\scriptsize,text=AccentTeal!90!black] at (5.05,0.58) {$\calH_{p-q}$};
      % Running minimum line
      \draw[AccentBrown!90!black,line width=1.25pt] (0,1.15) -- (1.35,1.15) -- (1.35,0) -- (5.80,0);
      % Ticks
      \node[below,font=\scriptsize,text=ContentsHeading] at (0,-0.04) {$0$};
      \node[below,font=\scriptsize,text=ContentsHeading] at (1.35,-0.04) {$q$};
      \node[below,font=\scriptsize,text=ContentsHeading] at (4.30,-0.04) {$1{-}p{+}q$};
      \node[below,font=\scriptsize,text=ContentsHeading] at (5.80,-0.04) {$1$};
      \node[left,font=\scriptsize,text=ContentsHeading] at (0,1.15) {$1$};
      \draw[axis] (0,0) -- (6.10,0);
      \draw[axis] (0,0) -- (0,1.45);
    \end{scope}

    % Bottom sub-diagram: Type 2 (H_p, Z_{1-p})
    \begin{scope}[shift={(0,0)}]
      \node[right,font=\scriptsize\bfseries,text=ContentsHeading] at (0,1.65) {Type 2: $(\calH_p,\; \calZ_{1-p})$ \quad $[N_v = p,\; \Delta_v = 0]$};
      \fill[pattern=north east lines,pattern color=AccentBrown!55] (0,0) rectangle (2.85,1.15);
      \fill[AccentBrown!10,opacity=0.6] (0,0) rectangle (2.85,1.15);
      \draw[AccentTeal!85!black,line width=0.7pt] (0,0) rectangle (2.85,1.15);
      \node[font=\scriptsize,text=AccentTeal!90!black] at (1.42,0.58) {$\calH_p$};
      \fill[AccentBlue!25] (2.85,0) rectangle (5.80,0.12);
      \node[above,font=\scriptsize,text=AccentBlue] at (4.32,0.10) {$\calZ_{1-p}$ ($\rho = 0$)};
      \draw[AccentBrown!90!black,line width=1.25pt] (0,1.15) -- (2.85,1.15) -- (2.85,0) -- (5.80,0);
      \node[below,font=\scriptsize,text=ContentsHeading] at (0,-0.04) {$0$};
      \node[below,font=\scriptsize,text=ContentsHeading] at (2.85,-0.04) {$p$};
      \node[below,font=\scriptsize,text=ContentsHeading] at (5.80,-0.04) {$1$};
      \node[left,font=\scriptsize,text=ContentsHeading] at (0,1.15) {$1$};
      \draw[axis] (0,0) -- (6.10,0) node[right,text=black] {$r$};
      \draw[axis] (0,0) -- (0,1.45);
    \end{scope}
  \end{scope}

  % Panel (b): Election 3 voter types
  \begin{scope}[shift={(7.7,0)}]
    \node[font=\fontsize{9.5}{11}\selectfont\bfseries,text=ContentsHeading] at (3.2,5.25) {(b) Election 3: three classes $\calH, \calM$, and $a^\star \notin \supp(D)$};

    % Type 1: M_x, a*, M_{1-p-x}, H_p
    \begin{scope}[shift={(0,3.60)}]
      \node[right,font=\scriptsize\bfseries,text=ContentsHeading] at (0,1.08) {Type 1: $(\calM_x, a^\star, \calM_{1-p-x}, \calH_p)$ \ $[N_v = x, \Delta_v = 2]$};
      \fill[pattern=north east lines,pattern color=AccentBrown!55] (0,0) rectangle (1.65,0.40);
      \draw[AccentBlue!85,line width=0.65pt] (0,0) rectangle (1.65,0.40);
      \draw[AccentBlue!85,line width=0.65pt] (1.65,0) rectangle (4.05,0.40);
      \draw[AccentTeal!85!black,line width=0.65pt] (4.05,0) rectangle (5.80,0.80);
      \draw[AccentBrown!90!black,line width=1.15pt] (0,0.40) -- (1.65,0.40) -- (1.65,0) -- (5.80,0);
      \fill[AccentBrown!90!black] (1.65,0) circle (1.6pt) node[below,font=\scriptsize] {$x$};
      \node[below,font=\scriptsize,text=ContentsHeading] at (4.05,-0.02) {$1{-}p$};
      \node[font=\scriptsize,text=AccentTeal!90!black] at (4.92,0.40) {$\calH_p$};
      \draw[axis] (0,0) -- (6.10,0);
      \draw[axis] (0,0) -- (0,0.92);
    \end{scope}

    % Type 2: M_{z-p}, H_p, a*, M_{1-z}
    \begin{scope}[shift={(0,1.80)}]
      \node[right,font=\scriptsize\bfseries,text=ContentsHeading] at (0,1.08) {Type 2: $(\calM_{z-p}, \calH_p, a^\star, \calM_{1-z})$ \ $[N_v = z, \Delta_v = 1]$};
      \fill[pattern=north east lines,pattern color=AccentBrown!55] (0,0) rectangle (4.65,0.40);
      \draw[AccentBlue!85,line width=0.65pt] (0,0) rectangle (2.90,0.40);
      \draw[AccentTeal!85!black,line width=0.65pt] (2.90,0) rectangle (4.65,0.80);
      \draw[AccentBlue!85,line width=0.65pt] (4.65,0) rectangle (5.80,0.40);
      \draw[AccentBrown!90!black,line width=1.15pt] (0,0.40) -- (4.65,0.40) -- (4.65,0) -- (5.80,0);
      \fill[AccentBrown!90!black] (4.65,0) circle (1.6pt) node[below,font=\scriptsize] {$z$};
      \node[below,font=\scriptsize,text=ContentsHeading] at (2.90,-0.02) {$z{-}p$};
      \node[font=\scriptsize,text=AccentTeal!90!black] at (3.78,0.60) {$\calH_p$};
      \draw[axis] (0,0) -- (6.10,0);
      \draw[axis] (0,0) -- (0,0.92);
    \end{scope}

    % Type 3: H_p, M_{y-p}, a*, M_{1-y}
    \begin{scope}[shift={(0,0)}]
      \node[right,font=\scriptsize\bfseries,text=ContentsHeading] at (0,1.08) {Type 3: $(\calH_p, \calM_{y-p}, a^\star, \calM_{1-y})$ \ $[N_v = p{+}y, \Delta_v = 1]$};
      \fill[pattern=north east lines,pattern color=AccentBrown!55]
        (0,0) -- (0,0.80) -- (1.75,0.80) -- (1.75,0.40) -- (2.35,0.40) -- (2.35,0) -- cycle;
      \draw[AccentTeal!85!black,line width=0.65pt] (0,0) rectangle (1.75,0.80);
      \draw[AccentBlue!85,line width=0.65pt] (1.75,0) rectangle (2.35,0.40);
      \draw[AccentBlue!85,line width=0.65pt] (2.35,0) rectangle (5.80,0.40);
      \draw[AccentBrown!90!black,line width=1.15pt] (0,0.80) -- (1.75,0.80) -- (1.75,0.40) -- (2.35,0.40) -- (2.35,0) -- (5.80,0);
      \fill[AccentBrown!90!black] (2.35,0) circle (1.6pt) node[below,font=\scriptsize] {$y$};
      \node[below,font=\scriptsize,text=ContentsHeading] at (1.75,-0.02) {$p$};
      \node[font=\scriptsize,text=AccentTeal!90!black] at (0.88,0.40) {$\calH_p$};
      \draw[axis] (0,0) -- (6.10,0) node[right,text=black] {$r$};
      \draw[axis] (0,0) -- (0,0.92);
    \end{scope}
  \end{scope}
\end{tikzpicture}
\caption{Rank-line block layouts and running minimums $\lowrho_v$ (hatched) for the voter types in Elections~2 and~3~\cite{CaiEtAl26b}. (a)~In Election~2, $\calH$ has $\rho = 1$ and $\calZ$ has $\rho = 0$. (b)~In Election~3, $\calH$ has $\rho = 2$, $\calM$ has $\rho = 1$, and $a^\star \notin \supp(D)$ has $\rho(a^\star) = 0$ (marked by the brown dot at ranks $x, z, y$, where $\lowrho_v$ drops to $0$).}
\label{fig:three-elections}
\end{figure}

\subsection{Election 2: The Two-Class Inequality}
\label{sec:election-2}

In Election~1, where $a^\star \notin \supp(D)$, a voter could achieve $\Delta_v = 0$ only by ranking $a^\star$ at the very top. If instead $D$ places probability mass $1 - p$ on candidates with $\rho = 0$ (including $a^\star$) and mass $p$ on candidates with $\rho = 1$, a voter can achieve $\Delta_v = 0$ simply by ranking all the $\rho = 1$ candidates below all the $\rho = 0$ candidates.

To turn this observation into a constraint on $J$, partition $\supp(D)$ into two classes (\cref{fig:three-elections}(a)): class $\calH$ with $\rho = 1$ and total probability $p \in (0, 1)$, and class $\calZ$ with $\rho = 0$ and total probability $1 - p$. Writing a subscript on a class for the probability mass of that block on the rank line, voters of Type~2 have ranking $(\calH_p, \calZ_{1-p})$ with $N_v = p$ and $\Delta_v = 0$, and give $\calH$ centered activation $-J(p) < 0$. To balance this negative activation, we introduce voters of Type~1 who split $\calH$ around $\calZ$: for $0 \le q < p$, they place mass $q$ of $\calH$ on $[0, q]$, all of $\calZ$ on $[q, 1 - p + q]$, and the other $p - q$ of $\calH$ on $[1 - p + q, 1]$:
\[
  \begin{array}{cccc}
    \toprule
    \text{Preference order (worst to best)} & \text{Centered activation of } \calH & N_v & \Delta_v \\
    \midrule
    \calH_q,\; \calZ_{1-p},\; \calH_{p-q} & J(1 - p + q) - J(q) & q & 1 \\
    \calH_p,\; \calZ_{1-p} & -J(p) & p & 0 \\
    \bottomrule
  \end{array}
\]
For Type~1 voters, $\lowrho_v$ drops from $1$ to $0$ at rank $q$ where $\calZ_{1-p}$ begins, so only $\calH_q$ contributes to $N_v = q$, while the top block $\calH_{p-q}$ gives $\Delta_v = 1$.

Write $\delta_p(q) := J(1 - p + q) - J(q)$ for the centered activation of $\calH$ in Type~1. Whenever $\delta_p(q) > 0$, giving Type~1 and Type~2 voters weights proportional to $J(p)$ and $\delta_p(q)$ makes the expected centered activation of $\calH$ (and hence of $\calZ$) equal to zero, so $D$ is $g$-stable with
\[
  \frac{\E_v[N_v]}{\E_v[\Delta_v]} \;=\; q + \frac{p\,\delta_p(q)}{J(p)}.
\]
By \cref{lem:biased-reduction}, this ratio cannot exceed $\lambda$; when $\delta_p(q) \le 0$, the resulting inequality holds immediately since $q \le \lambda$ and $J(p) > 0$.

\begin{lemma}[Election 2 bound]
\label{lem:election-2}
For every $0 < p < 1$ and $0 \le q \le \min\{p, \lambda\}$,
\begin{equation}
  p\,\bigl[J(1 - p + q) - J(q)\bigr] \;\le\; (\lambda - q)J(p).
  \label{eq:two-class}
\end{equation}
In particular, setting $q = 0$ gives
\begin{equation}
  p\,J(1 - p) \;\le\; \lambda J(p)
  \qquad\text{and}\qquad
  (1 - p)J(p) \;\le\; \lambda J(1 - p).
  \label{eq:two-class-zero}
\end{equation}
\end{lemma}

\subsection{Election 3: A Differential Inequality}
\label{sec:election-3}

Election~1 constrained $\phi$ pointwise by placing $a^\star$ outside $\supp(D)$, while Election~2 constrained $J$ by balancing two classes in $\supp(D)$. To make $\phi$ and $J$ interact, we now combine both ideas (\cref{fig:three-elections}(b)). Partition $\supp(D)$ into class $\calH$ with $\rho = 2$ and probability $p \in (0, 1)$, and class $\calM$ with $\rho = 1$ and probability $1 - p$, while keeping $a^\star \notin \supp(D)$ with $\rho(a^\star) = 0$ and probability $0$.

Choose ranks $x, y, z$ with $0 \le x < \min\{\lambda, 1 - p\}$, $p \le y < \lambda$, and $\lambda + p/2 < z \le 1$. We consider three types of voters, who place the block $\calH_p$ at the top of their ranking, on the intermediate interval $[z - p, z]$, and at the bottom, respectively, while placing $a^\star$ at ranks $x, z$, and $y$. Writing
\[
  R \;:=\; J(p),
  \qquad
  P \;:=\; J(1 - p),
  \qquad
  K \;:=\; J(z - p) - J(z),
\]
the three voter types have the following centered activations and values of $(N_v, \Delta_v)$:
\[
  \begin{array}{cccccc}
    \toprule
    & \multicolumn{2}{c}{\text{Centered activation}} & & & \\
    \cmidrule(lr){2-3}
    \text{Preference order (worst to best)} & \calH & a^\star & N_v & \Delta_v & N_v - \lambda\Delta_v \\
    \midrule
    \calM_x,\; a^\star,\; \calM_{1-p-x},\; \calH_p & P & -\phi(x) & x & 2 & -(2\lambda - x) \\
    \calM_{z-p},\; \calH_p,\; a^\star,\; \calM_{1-z} & K & -\phi(z) & z & 1 & z - \lambda \\
    \calH_p,\; \calM_{y-p},\; a^\star,\; \calM_{1-y} & -R & -\phi(y) & p + y & 1 & p + y - \lambda \\
    \bottomrule
  \end{array}
\]
To read off $N_v$ and $\Delta_v$ (\cref{fig:three-elections}(b)), note that $a^\star$ splits each voter's rank line in two: every candidate ranked above $a^\star$ has $\lowrho_v(b) = \rho(a^\star) = 0$, contributing nothing to $N_v$ and its full $\rho$-value to $\rho(b) - \lowrho_v(b)$, whereas below $a^\star$, $\lowrho_v$ is the running minimum of $\rho$:
\begin{itemize}
  \item In Type~1, only $\calM_x$ lies below $a^\star$, so $N_v = x$; above $a^\star$, $\calH_p$ has $\rho = 2$, so $\Delta_v = 2$.
  \item In Type~2, both $\calM_{z-p}$ and $\calH_p$ lie below $a^\star$ with running minimum $1$ across both, so $N_v = z$; the gap $\rho - \lowrho_v$ is $1$ on $\calH_p$ and on $\calM_{1-z}$ above $a^\star$, so $\Delta_v = 1$.
  \item In Type~3, the running minimum is $2$ on $\calH_p$ and $1$ on $\calM_{y-p}$, so $N_v = 2p + (y - p) = p + y$, while $\calM_{1-y}$ above $a^\star$ gives $\Delta_v = 1$.
\end{itemize}

To balance the centered activations of both $\calH$ and $a^\star$ to zero, we choose voter-type weights $(w_1, w_2, w_3)$ orthogonal to both columns $(P, K, -R)$ and $(-\phi(x), -\phi(z), -\phi(y))$, namely their $2 \times 2$ minors:
\begin{equation}
  \begin{aligned}
    w_1 &\;:=\; -K\phi(y) - R\phi(z), \\
    w_2 &\;:=\; R\phi(x) + P\phi(y), \\
    w_3 &\;:=\; K\phi(x) - P\phi(z).
  \end{aligned}
  \label{eq:three-class-weights}
\end{equation}
Direct substitution gives $P w_1 + K w_2 - R w_3 = 0$ and $-\phi(x)w_1 - \phi(z)w_2 - \phi(y)w_3 = 0$. When $w_1, w_2, w_3 > 0$, stability and \cref{lem:biased-reduction} require the weighted sum of $N_v - \lambda\Delta_v$ across the three voter types to be nonpositive:
\[
  -(2\lambda - x)w_1 + (z - \lambda)w_2 + (p + y - \lambda)w_3 \;\le\; 0.
\]
Substituting \cref{eq:three-class-weights}, the coefficient of $\phi(z)$ is $\Delta := (2\lambda - x)R - (p + y - \lambda)P$, which bounds $\phi(z)$ at high rank in terms of $\phi(x)$ and $\phi(y)$ at low ranks:

\begin{lemma}[Election 3 bound]
\label{lem:election-3}
Suppose
\[
  0 < p < 1,
  \qquad
  0 \le x < \min\{\lambda, 1 - p\},
  \qquad
  p \le y < \lambda,
  \qquad
  \lambda + \frac{p}{2} < z \le 1,
\]
and $(z - \lambda)R + (p + y - \lambda)K > 0$. Then $K > 0$, $\Delta := (2\lambda - x)R - (p + y - \lambda)P > 0$, and
\begin{equation}
  \bigl[(z - \lambda)R + (p + y - \lambda)K\bigr]\phi(x)
  + \bigl[(z - \lambda)P + (2\lambda - x)K\bigr]\phi(y)
  + \Delta\,\phi(z) \;\le\; 0.
  \label{eq:three-class}
\end{equation}
\end{lemma}

If one of the weights in \cref{eq:three-class-weights} is nonpositive, either \cref{eq:three-class} already holds or two of the three voter types violate \cref{eq:height-inequality} on their own. We prove these cases, together with the finite realization of Elections~2 and~3, in \cref{app:finite-two-three,app:equal-weights}.

\section{The Variational Problem Defining \texorpdfstring{$\lambda^\star$}{lambda*}}
\label{sec:variational}

We next show how the necessary inequalities from Elections~1, 2, and~3 determine the optimal parameter $\lambda^\star$ and the shape of a minimizing activation. The choice is made on an interval of low ranks; the inequalities then govern how far the primitive $J$ can be continued at high ranks.

\subsection{Constraining \texorpdfstring{$J$}{J} at Low and Middle Ranks}
\label{sec:low-ranks}

Suppose first that $\phi(r) = \lambda - r$ on an interval of middle ranks. By \cref{eq:sigma-relations}, integrating gives
\[
  J(r) \;=\; Q_\lambda(r) - S,
\]
where the scalar $S \ge 0$ is the cumulative slack below the line $\lambda - u$ accumulated at smaller ranks:
\[
  S \;=\; \int_0^r \bigl[\lambda - u - \phi(u)\bigr]\,\mathrm{d}u.
\]
Now apply the $q = 0$ two-class inequality~\eqref{eq:two-class-zero} at a low rank $r$ whose complement $1 - r$ lies in the middle interval:
\[
  (1 - r)J(r) \;\le\; \lambda J(1 - r) \;=\; \lambda\bigl[Q_\lambda(1 - r) - S\bigr].
\]
Dividing by $1 - r$ gives an explicit rational upper bound on the primitive $J$ at low ranks:
\[
  C_{\lambda, S}(r) \;:=\; \frac{\lambda\bigl[Q_\lambda(1 - r) - S\bigr]}{1 - r}
  \;=\; \lambda^2 - \frac{\lambda}{2} + \frac{\lambda r}{2} - \frac{\lambda S}{1 - r}.
\]

At low ranks, $J$ must balance two constraints: its derivative $J'(r) = \phi(r)$ lies in $[0, \lambda - r]$, and its value cannot exceed $C_{\lambda, S}(r)$. To isolate the region where $\phi$ transitions from $C'_{\lambda, S}$ to $\lambda - r$ (\cref{fig:optimal-activation}(b)), we partition $[p_-, z_1]$ by fixed rational endpoints (\cref{app:rank-intervals})
\[
  0 \;<\; p_- \;<\; \ell \;<\; x_- \;<\; p_+ \;<\; x_+ \;<\; t_e \;<\; \lambda \;<\; z_0 \;<\; z_B \;<\; z_1 \;<\; 1
\]
into four intervals:
\begin{itemize}
  \item On $[p_-, \ell]$, we set $J(r) = C_{\lambda, S}(r)$ (where the $q = 0$ Election~2 bound is tight).
  \item On $[\ell, t_e]$, we choose $\phi = J'$ connecting $C_{\lambda, S}$ to $Q_\lambda - S$ subject to
  \begin{equation}
    \begin{gathered}
      J(\ell) = C_{\lambda, S}(\ell), \qquad J(t_e) = Q_\lambda(t_e) - S, \\
      0 \le J'(r) \le \lambda - r, \qquad J(r) \le C_{\lambda, S}(r) \quad (\ell \le r \le t_e).
    \end{gathered}
    \label{eq:low-constraints}
  \end{equation}
  \item On $[t_e, z_0]$, we set $J(r) = Q_\lambda(r) - S$ (where the Election~1 bound $\phi(r) = \lambda - r$ is tight).
  \item On $[z_0, z_1]$, $J$ is continued by the Election~3 differential inequality (\cref{sec:high-ranks}).
\end{itemize}
In \cref{prop:reduction}, we show that restricting to functions satisfying these boundary conditions and $S \le S_+(\lambda)$ (\cref{app:rank-intervals}) does not weaken the lower bound.

\subsection{Determining the High-Rank Continuation from Election 3}
\label{sec:high-ranks}

Once $S$ and the low-rank primitive $J$ on $[\ell, t_e]$ have been chosen, the three-class inequality~\eqref{eq:three-class} of Election~3 bounds $\phi(z) = J'(z)$ from above at each high rank $z \in [z_0, z_1]$. For the triples $(p, x, y)$ used in our variational problem, both $1 - p$ and $z - p$ lie in the middle interval $[t_e, z_0]$ where $J = Q_\lambda - S$, so their primitive values are already fixed by $S$. Writing
\[
  R \;=\; J(p),
  \qquad
  P \;=\; Q_\lambda(1 - p) - S,
  \qquad
  \Delta \;=\; (2\lambda - x)R - (p + y - \lambda)P,
\]
positivity of $\Delta$ (by \cref{lem:election-3}) lets us solve \cref{eq:three-class} for $\phi(z)$. When the primitive at $z$ takes the value $u$, the resulting upper bound on $\phi(z)$ is
\begin{equation}
  H(z, u; p, x, y) \;:=\; -\frac{1}{\Delta}\Bigl\{
    (z - \lambda)\bigl[R\phi(x) + P\phi(y)\bigr]
    + \bigl[Q_\lambda(z - p) - S - u\bigr]\bigl[(p + y - \lambda)\phi(x) + (2\lambda - x)\phi(y)\bigr]
  \Bigr\}.
  \label{eq:H-def}
\end{equation}
Because \cref{eq:three-class} must hold for every admissible triple $(p, x, y)$ and $\phi(z) \le \lambda - z$ must also hold by Election~1, the derivative $J'(z)$ cannot exceed the essential infimum of these bounds.

To make the primitive at high ranks as large as possible, we set the derivative equal to this smallest upper bound. Starting at the left endpoint $z_0$ of the high-rank interval, this defines the maximal high-rank continuation $Y(z)$ by the ordinary differential equation
\begin{equation}
  Y'(z) \;=\; \min\Bigl\{\lambda - z,\; \operatorname*{ess\,inf}_{(p, x, y) \in \calU} H\bigl(z, Y(z); p, x, y\bigr)\Bigr\},
  \qquad
  Y(z_0) \;=\; Q_\lambda(z_0) - S,
  \label{eq:high-ode}
\end{equation}
where $\calU$ is the box of triples $(p, x, y)$ specified in \cref{app:variational,app:rank-intervals}. Because each function $u \mapsto H(z, u; p, x, y)$ is affine in $u$ with uniformly bounded slope, \cref{eq:high-ode} has a unique solution $Y$ on $[z_0, z_1]$, and standard ODE comparison shows that any primitive $J$ satisfying the same differential inequalities lies weakly below $Y$.

We take the largest possible continuation $Y(z)$ because the two-class inequalities from Election~2 bound $J(z)$ from below at high ranks. Specifically, in \cref{eq:two-class}, set the class mass to $z \in [z_B, z_1]$ and the split rank to $q \in \calQ \subset [p_-, \ell]$:
\[
  J(z) \;\ge\; \frac{z\bigl[J(1 - z + q) - J(q)\bigr]}{\lambda - q}.
\]
For every pair $(z, q) \in [z_B, z_1] \times \calQ$, the rank $1 - z + q$ lies in the middle interval $[t_e, z_0]$ (where $J = Q_\lambda - S$) and $q$ lies in $[p_-, \ell]$ (where $J = C_{\lambda, S}$). Thus the right-hand side is the explicit function
\[
  T_{\lambda, S}(z, q) \;:=\; \frac{z\bigl[Q_\lambda(1 - z + q) - S - C_{\lambda, S}(q)\bigr]}{\lambda - q}.
\]
The upper bound $Y(z)$ from Election~3 and the lower bound $T_{\lambda, S}(z, q)$ from Election~2 are compatible if and only if
\begin{equation}
  Y(z) \;\ge\; T_{\lambda, S}(z, q)
  \qquad \text{for all } (z, q) \in [z_B, z_1] \times \calQ.
  \label{eq:high-binary-compatibility}
\end{equation}

\subsection{The Least Feasible Value and Shape of the Minimizer}
\label{sec:least-feasible}

Combining these steps reduces the lower bound to a problem in the scalar $S$ and the function $\phi$ on $[\ell, t_e]$. Let $\calK_\lambda$ be the set of pairs $(S, \phi)$ satisfying $S \le S_+(\lambda)$ and \cref{eq:low-constraints}, with the endpoint and essential bounds specified in \cref{app:variational}, and let $Y$ be the maximal continuation in \cref{eq:high-ode}. We define
\begin{equation}
  \lambda^\star \;:=\; \inf\left\{\lambda \in [\lambda_0, \lambda_1] :
  \begin{array}{l}
  \text{there exists } (S, \phi) \in \calK_\lambda \text{ such that}\\
  Y(z) \ge T_{\lambda, S}(z, q) \text{ for all } (z, q) \in [z_B, z_1] \times \calQ
  \end{array}\right\}.
  \label{eq:lambda-star-def}
\end{equation}
The full definitions of $\calK_\lambda$, $\calU$, and $\calQ$ appear in \cref{app:variational}.

To prove that $\lambda^\star$ is the exact threshold for $g$-stable lotteries, we first show that any activation achieving distortion $1 + 2\lambda$ gives a feasible solution to this reduced problem.

\begin{proposition}[Reduction to the variational problem]
\label{prop:reduction}
Suppose an activation $g \in \calG$ has distortion $M(g) \le 1 + 2\lambda$ for some $\lambda \in [\lambda_0, \lambda_1]$. Then there exist $S \le S_+(\lambda)$ and a function $\phi$ on $[\ell, t_e]$ satisfying \cref{eq:low-constraints} whose maximal continuation $Y$ in \cref{eq:high-ode} satisfies all two-class bounds $Y(z) \ge T_{\lambda, S}(z, q)$ in \cref{eq:high-binary-compatibility}.
\end{proposition}

The proof of \cref{prop:reduction} (\cref{app:reduction}) has two steps. First, we remove the slack $\sigma(r) = \lambda - r - \phi(r)$ from the middle ranks so that $J = Q_\lambda - S$ holds identically there; on the subinterval that affects the high-rank two-class bounds, we simultaneously remove $\kappa_f$ times as much slack below $p_-$, raising $J$ in a ratio that preserves both $q = 0$ two-class inequalities~\eqref{eq:two-class-zero} while an ODE comparison controls the change in the high-rank continuation. Second, we raise the low-rank primitive to $C_{\lambda, S}$ on $[p_-, \ell]$ and to its nondecreasing envelope above $C_{\lambda, S}(\ell)$ on $[\ell, t_e]$, which increases $J(p)$ and preserves both \cref{eq:three-class} and \cref{eq:high-binary-compatibility}.

Next, to construct an optimal activation $g^\star$ for the upper bound, we show that the infimum in \cref{eq:lambda-star-def} is attained and makes at least one high-rank two-class inequality tight.

\begin{proposition}[Attainment and active two-class bound]
\label{prop:attainment}
The infimum defining $\lambda^\star$ in \cref{eq:lambda-star-def} is attained. Furthermore, for any minimizing pair $(S, \phi)$ at $\lambda^\star$, its maximal continuation $Y$ touches at least one two-class lower bound:
\[
  Y(z_\star) \;=\; T_{\lambda^\star, S}(z_\star, q_\star)
  \qquad \text{for some } (z_\star, q_\star) \in [z_B, z_1] \times \calQ.
\]
\end{proposition}

In \cref{app:attainment}, we work with bounded measurable functions $\phi \in L^\infty([\ell, t_e])$, where weak-* convergence gives uniform convergence of the primitives $J$. To pass to the limit in the ODE~\eqref{eq:high-ode} (where $\phi$ is evaluated at individual ranks $x, y$), we average its controls against probability densities in $x$ and $y$. Each averaged continuation is weak-* continuous, and the maximal continuation $Y$ is the infimum of these averaged continuations, making the maximum two-class violation weak-* lower semicontinuous. Mixing with a polynomial trial strictly inside the low-rank constraints then shows that any minimizer at $\lambda^\star$ must make at least one two-class bound tight; otherwise a small perturbation would allow a smaller $\lambda$.

Finally, we use local perturbations at $\lambda^\star$ to prove the regularity and shape of a minimizing $\phi$. If a positive-measure set of ranks after the last contact $c$ with $C_{\lambda^\star, S}$ where $\phi(r) < \lambda^\star - r$ were uniformly separated from all tight three-class inequalities, raising $\phi$ on that set while lowering $S$ to keep $J(t_e) = Q_{\lambda^\star}(t_e) - S$ fixed would strictly improve the high-rank continuation. Thus almost every such rank appears in arbitrarily tight three-class triples; their supporting inequalities prove Lipschitz continuity, after which compactness gives an attained tight triple through almost every such rank.

\begin{proposition}[Shape of a minimizing $\phi$]
\label{prop:shape}
At $\lambda = \lambda^\star$, a minimizing $\phi$ can be chosen continuous and strictly decreasing on $[\ell, t_e]$. It coincides with $C'_{\lambda^\star, S}$ up to a last contact rank $c \in (x_-, p_+)$, where it leaves $C'_{\lambda^\star, S}$ and later meets the line $\lambda^\star - r$ before $x_+$. Between these two contacts, $\phi$ is Lipschitz and satisfies, almost everywhere,
\[
  -1 \;\le\; \phi'(r) \;\le\; -\underline{\theta} \;<\; 0.
\]
It then follows $\lambda^\star - r$ into the high ranks. After leaving $\lambda^\star - r$ inside $(z_0, z_B)$, the high-rank derivative $Y'(z)$ is concave and has derivative strictly less than $-1$.
\end{proposition}

\begin{figure}[!t]
\centering
\begin{tikzpicture}[y=1.0cm,font=\small,
    axis/.style={-{Latex[length=4.5pt,width=3.2pt]},draw=ContentsHeading,line width=0.55pt},
    ptr/.style={-{Latex[length=3.5pt,width=2.5pt]},line width=0.45pt}]
  % Panel (a): Low-rank phi* and supporting lines
  \begin{scope}[shift={(0,0)}]
    \node[font=\fontsize{9.5}{11}\selectfont\bfseries,text=ContentsHeading] at (3.0,4.35) {(a) Low-rank $\phi^\star(r)$ and supporting lines};

    % Ceiling line lambda* - r
    \draw[densely dotted,AccentBlue!80,line width=0.8pt] (0,3.60) -- (5.40,0)
      node[pos=0.82,above right=-1pt,font=\scriptsize,text=AccentBlue] {$\lambda^\star - r$};

    % Cap derivative C'(r) continuing past c
    \draw[densely dashed,AccentTeal!75,line width=0.75pt]
      plot[domain=0.85:3.10,smooth,samples=20] (\x,{2.78 - 0.12*\x - 0.085*\x*\x});
    \node[font=\scriptsize,text=AccentTeal,anchor=west] at (0.08,1.95) {$C'_{\lambda^\star,S}(r)$};
    \draw[ptr,AccentTeal] (0.95,2.12) to[out=45,in=-120] (1.35,2.44);

    % Actual low-rank function phi*(r):
    % [0, 0.85]: affine head from (0, 2.88) to (0.85, 2.62)
    % [0.85, 1.85]: C'(r) from (0.85, 2.62) to c=(1.85, 2.27)
    % [1.85, 3.30]: transition curve from (1.85, 2.27) to r_c=(3.30, 1.40) on lambda*-r
    % [3.30, 5.40]: line lambda* - r from (3.30, 1.40) to (5.40, 0)
    \draw[AccentBlue!90!black,line width=1.25pt]
      (0,2.88) -- (0.85,2.62)
      plot[domain=0.85:1.85,smooth,samples=15] (\x,{2.78 - 0.12*\x - 0.085*\x*\x})
      .. controls (2.30,2.05) and (2.80,1.74) .. (3.30,1.40)
      -- (5.40,0);

    % Supporting line of slope -theta touching at x = 2.45, y_val = 1.93
    \draw[AccentBrown!90!black,line width=0.85pt]
      (1.35,2.56) -- (3.75,1.18);
    \node[font=\scriptsize,text=AccentBrown!90!black,anchor=west] at (2.68,2.15) {slope $-\theta$};
    \draw[ptr,AccentBrown!90!black] (2.95,1.98) to[out=-110,in=45] (2.52,1.95);

    % Join points x_h, c, x, r_c
    \fill[ContentsHeading,draw=white,line width=0.5pt] (0.85,2.62) circle (1.9pt);
    \fill[AccentTeal,draw=white,line width=0.5pt] (1.85,2.27) circle (2.1pt);
    \fill[AccentBrown!90!black,draw=white,line width=0.5pt] (2.45,1.93) circle (2.1pt);
    \fill[AccentBlue,draw=white,line width=0.5pt] (3.30,1.40) circle (2.0pt);

    \draw[densely dashed,ContentsHeading!45,line width=0.4pt] (0.85,0) -- (0.85,2.62);
    \draw[densely dashed,AccentTeal!55,line width=0.4pt] (1.85,0) -- (1.85,2.27);
    \draw[densely dashed,AccentBrown!65,line width=0.4pt] (2.45,0) -- (2.45,1.93);
    \draw[densely dashed,AccentBlue!55,line width=0.4pt] (3.30,0) -- (3.30,1.40);

    \node[below,font=\scriptsize,text=ContentsHeading] at (0,-0.06) {$0$};
    \node[below,font=\scriptsize,text=ContentsHeading] at (0.85,-0.06) {$x_h$};
    \node[below,font=\scriptsize,text=AccentTeal] at (1.85,-0.06) {$c$};
    \node[below,font=\scriptsize,text=AccentBrown!90!black] at (2.45,-0.06) {$x, y$};
    \node[below,font=\scriptsize,text=AccentBlue] at (3.30,-0.06) {$r_c$};
    \node[below,font=\scriptsize,text=ContentsHeading] at (5.40,-0.06) {$\lambda^\star$};

    \draw[axis] (0,0) -- (5.80,0) node[right,text=black] {$r$};
    \draw[axis] (0,0) -- (0,3.85) node[left,text=black] {$\phi$};
  \end{scope}

  % Panel (b): High-rank phi* and concave envelope
  \begin{scope}[shift={(7.8,0)}]
    \node[font=\fontsize{9.5}{11}\selectfont\bfseries,text=ContentsHeading] at (3.1,4.35) {(b) High-rank $\phi^\star(z)$ and quadratic piece};

    % Axes: r-axis at y = 3.40 (since phi <= 0 on [lambda*, 1])
    \draw[axis] (0,3.40) -- (5.80,3.40) node[right,text=black] {$z$};
    \draw[axis] (0,0) -- (0,3.85) node[left,text=black] {$\phi$};

    % Ceiling line lambda* - z: from (0, 3.40) to (5.40, 1.95)
    \draw[densely dotted,AccentBlue!80,line width=0.8pt] (0,3.40) -- (5.40,1.95)
      node[right,font=\scriptsize,text=AccentBlue] {$\lambda^\star - z$};

    % Individual fixed-p ODE upper bounds H_p(z)
    \draw[densely dashed,AccentTeal!65,line width=0.65pt]
      (2.40,2.85) .. controls (3.30,2.30) and (4.10,1.65) .. (4.85,0.92);
    \draw[densely dashed,AccentTeal!65,line width=0.65pt]
      (2.75,2.72) .. controls (3.60,2.05) and (4.30,1.32) .. (5.05,0.48);
    \node[font=\scriptsize,text=AccentTeal,anchor=west] at (3.65,2.35) {ODE bounds $H_p(z)$};

    % Actual high-rank function phi*(z):
    % [0, 2.65]: follows lambda* - z from (0, 3.40) to z_d=(2.65, 2.69)
    % [2.65, 4.15]: concave ODE envelope from (2.65, 2.69) to z_*=(4.15, 1.55)
    % [4.15, 5.40]: quadratic piece from (4.15, 1.55) to (5.40, 0.20) where phi(1) = -1
    \draw[AccentBrown!90!black,line width=1.25pt]
      (0,3.40) -- (2.65,2.69)
      .. controls (3.20,2.38) and (3.70,1.98) .. (4.15,1.55)
      .. controls (4.60,1.10) and (5.02,0.62) .. (5.40,0.20);

    \fill[AccentBlue,draw=white,line width=0.5pt] (2.65,2.69) circle (2.0pt);
    \fill[AccentBrown!90!black,draw=white,line width=0.5pt] (4.15,1.55) circle (2.1pt);
    \fill[ContentsHeading,draw=white,line width=0.5pt] (5.40,0.20) circle (2.0pt);

    \draw[densely dashed,AccentBlue!55,line width=0.4pt] (2.65,3.40) -- (2.65,2.69);
    \draw[densely dashed,AccentBrown!65,line width=0.4pt] (4.15,3.40) -- (4.15,1.55);
    \draw[densely dashed,ContentsHeading!45,line width=0.4pt] (5.40,3.40) -- (5.40,0.20) -- (0,0.20);

    \node[above,font=\scriptsize,text=ContentsHeading] at (0.18,3.42) {$\lambda^\star$};
    \node[above,font=\scriptsize,text=AccentBlue] at (2.65,3.42) {$z_d$};
    \node[above,font=\scriptsize,text=AccentBrown!90!black] at (4.15,3.42) {$z_\star$};
    \node[above,font=\scriptsize,text=ContentsHeading] at (5.40,3.42) {$1$};
    \node[left,font=\scriptsize,text=ContentsHeading] at (0,3.40) {$0$};
    \node[left,font=\scriptsize,text=ContentsHeading] at (0,0.20) {$-1$};

    \node[font=\scriptsize,text=AccentBrown!90!black,anchor=east] at (4.00,0.88) {quadratic piece on $[z_\star, 1]$};
    \draw[ptr,AccentBrown!90!black] (4.05,0.88) to[out=0,in=-135] (4.78,0.86);
  \end{scope}
\end{tikzpicture}
\caption{Schematic shape of the minimizing function $\phi^\star$ (\cref{prop:shape} and \cref{sec:completing-score}). (a)~On $[0, \lambda^\star]$, $\phi^\star$ begins with an affine piece on $[0, x_h]$, follows $C'_{\lambda^\star,S}$ on $[x_h, c]$, transitions with supporting lines of slope $-\theta \in [-1, -\underline{\theta}]$ at active positions $x, y$ to the line $\lambda^\star - r$ at $r_c < x_+$, and stays on $\lambda^\star - r$ through $\lambda^\star$. (b)~On $[\lambda^\star, 1]$, $\phi^\star$ follows $\lambda^\star - z$ up to $z_d \in [z_-, z_+]$, follows the concave envelope of the fixed-$p$ bounds $H_p(z)$ up to the tangential contact $z_\star$, and finishes with a concave quadratic piece reaching $\phi^\star(1) = -1$.}
\label{fig:supporting-lines}
\end{figure}

To see why active ranks imply this shape (\cref{fig:supporting-lines} and \cref{app:shape}), fix a mass $p$ for which a three-class inequality is tight in \cref{eq:high-ode}, and write
\[
  \theta \;:=\; \frac{R\phi(x) + P\phi(y)}{\Delta} \;\in\; (\underline{\theta}, 1),
  \qquad
  \alpha \;:=\; \frac{(p + y - \lambda^\star)\phi(x) + (2\lambda^\star - x)\phi(y)}{\Delta} \;>\; 0,
\]
so that $H\bigl(z, Y(z); p, x, y\bigr) = -\theta(z - \lambda^\star) - \alpha\bigl[Q_{\lambda^\star}(z - p) - S - Y(z)\bigr]$. In \cref{eq:H-def}, moving $x$ to another rank $r \in [x_-, x_+]$ changes $H$ by a negative multiple of $\phi(r) + \theta r - \bigl(\phi(x) + \theta x\bigr)$, and similarly for $y \in [\max\{x_-, p\}, x_+]$ (\cref{app:shape}). Thus, minimizing $H$ over $(x, y)$ requires $x$ to maximize $\phi(r) + \theta r$ on $[x_-, x_+]$ and $y$ to maximize $\phi(r) + \theta r$ on $[\max\{x_-, p\}, x_+]$. At every active position $x$ (or $y > p$), the line of slope $-\theta$ through $(x, \phi(x))$ is therefore a supporting line lying above the graph of $\phi$. Comparing neighboring positions rules out jumps and bounds the slope of $\phi$ between $-1$ and $-\underline{\theta}$. Crucially, $\theta$ and $\alpha$ depend only on $p$ and $\phi$ on $[\ell, t_e]$, and \emph{not} on the high rank $z$. Each fixed-$p$ bound therefore imposes a linear differential inequality $Y'(z) \le -\theta(z - \lambda^\star) - \alpha\bigl[Q_{\lambda^\star}(z - p) - S - Y(z)\bigr]$, which makes $Y'(z)$ concave and prevents it from returning to $\lambda^\star - z$ after departure.

\section{The Matching Upper Bound}
\label{sec:upper-bound}

For any continuous, strictly decreasing function $\phi \colon [0, 1] \to \R$ with $\int_0^1 \phi(r)\,\mathrm{d}r = 0$ and $\phi(1) = -1$, reversing the normalization of \cref{eq:score-primitive-def} defines a continuous, strictly increasing activation $g_\phi \colon [0, 1] \to [0, 1]$ by
\[
  g_\phi(r) \;:=\; \frac{\phi(0) - \phi(r)}{1 + \phi(0)}.
\]
It satisfies $g_\phi(0) = 0$ and $g_\phi(1) = 1$, and has centered activation $-\phi$ after positive rescaling. Starting from a minimizing pair $(S, \phi)$ at $\lambda = \lambda^\star$, we complete $\phi$ to such a function $\phi^\star$ on $[0, 1]$ and set $g^\star := g_{\phi^\star}$. Throughout this section, we write $\lambda = \lambda^\star$, $\phi = \phi^\star$, and $J(r) = \int_0^r \phi(u)\,\mathrm{d}u$.

\subsection{Slicing by Threshold and the Verifier Lemma}
\label{sec:slicing-verifier}

Fix a finite election, a $g_\phi$-stable lottery $D$, an optimal candidate $a^\star \in C$, and distances $\rho \colon C \to \R_{\ge 0}$ with $\rho(a^\star) = 0$. By \cref{lem:biased-reduction}, proving that $D$ has distortion at most $1 + 2\lambda$ is equivalent to proving $\E_v[N_v(D)] \le \lambda\,\E_v[\Delta_v]$. Because this inequality cannot hold for every voter $v$ individually, we slice $N_v(D)$ by threshold $t \ge 0$, offset each voter's slice width using stability, and telescope the resulting per-threshold bounds across $t$ with a potential function~\cite{CaiEtAl26b}.

\begin{figure}[!t]
\centering
\begin{tikzpicture}[y=1.0cm,font=\small,
    axis/.style={-{Latex[length=4.5pt,width=3.2pt]},draw=ContentsHeading,line width=0.55pt},
    ptr/.style={-{Latex[length=3.5pt,width=2.5pt]},line width=0.45pt},
    dptr/.style={{Latex[length=3.2pt,width=2.3pt]}-{Latex[length=3.2pt,width=2.3pt]},line width=0.5pt}]
  % Panel (a): S_t and S_{t+Delta_v} on the rank line
  \begin{scope}[shift={(0,0)}]
    \node[font=\fontsize{9.5}{11}\selectfont\bfseries,text=ContentsHeading] at (3.2,4.45) {(a) Upper-level sets $S_t$ and $S_{t+\Delta_v}$ on $v$'s rank line};

    % Threshold lines t and t + Delta_v
    \draw[densely dashed,AccentBrown!85,line width=0.65pt] (0,1.50) -- (6.15,1.50)
      node[right,font=\scriptsize,text=AccentBrown!90!black] {$t$};
    \draw[densely dashed,AccentTeal!85,line width=0.65pt] (0,2.55) -- (6.15,2.55)
      node[right,font=\scriptsize,text=AccentTeal!90!black] {$t + \Delta_v$};

    % Boxes in S_t (orange hatching) and S_{t+Delta_v} (teal hatching)
    % Box 1: [0, 1.35], height 3.20 (in S_{t+Delta_v} and S_t)
    \fill[pattern=north east lines,pattern color=AccentBrown!55] (0,0) rectangle (1.35,3.20);
    \fill[pattern=north west lines,pattern color=AccentTeal!60] (0,1.50) rectangle (1.35,3.20);
    \draw[ContentsHeading!85,line width=0.6pt] (0,0) rectangle (1.35,3.20);

    % Box 2: [1.35, 2.40], height 2.15 (in S_t only)
    \fill[pattern=north east lines,pattern color=AccentBrown!55] (1.35,0) rectangle (2.40,2.15);
    \draw[ContentsHeading!85,line width=0.6pt] (1.35,0) rectangle (2.40,2.15);

    % Box 3: [2.40, 3.30], height 1.10 (outside S_t)
    \draw[ContentsHeading!85,line width=0.6pt] (2.40,0) rectangle (3.30,1.10);

    % Box 4: [3.30, 4.35], height 2.30 (in S_t, above q_v(t))
    \fill[pattern=north east lines,pattern color=AccentBrown!55] (3.30,0) rectangle (4.35,2.30);
    \draw[ContentsHeading!85,line width=0.6pt] (3.30,0) rectangle (4.35,2.30);

    % Box 5: [4.35, 5.15], height 0.85 (outside S_t)
    \draw[ContentsHeading!85,line width=0.6pt] (4.35,0) rectangle (5.15,0.85);

    % a* at s_v = 5.15 (height 0) and Box 6: [5.15, 5.85], height 0.65
    \draw[ContentsHeading!85,line width=0.6pt] (5.15,0) rectangle (5.85,0.65);

    % Ticks and labels
    \draw[ContentsHeading,line width=0.5pt] (1.35,0) -- (1.35,-0.08)
      node[below,font=\scriptsize,text=AccentTeal!90!black] {$z_v(t{+}\Delta_v)$};
    \draw[ContentsHeading,line width=0.5pt] (2.40,0) -- (2.40,-0.08)
      node[below,font=\scriptsize,text=AccentBrown!90!black] {$q_v(t)$};
    \draw[ContentsHeading,line width=0.5pt] (4.35,0) -- (4.35,-0.08)
      node[below,font=\scriptsize,text=AccentBrown!90!black] {$z_v(t)$};
    \draw[ContentsHeading,line width=0.5pt] (5.15,0) -- (5.15,-0.08)
      node[below,font=\scriptsize,text=AccentBlue] {$s_v$};
    \draw[ContentsHeading,line width=0.5pt] (5.85,0) -- (5.85,-0.08)
      node[below,font=\scriptsize,text=ContentsHeading] {$1$};
    \node[below,font=\scriptsize,text=ContentsHeading] at (0,-0.08) {$0$};

    \draw[dptr,AccentBrown!85!black] (0,-0.62) -- (2.40,-0.62)
      node[midway,below=1pt,font=\scriptsize,text=AccentBrown!85!black] {all in $S_t$};

    \draw[axis] (0,0) -- (6.15,0);
    \draw[axis] (0,0) -- (0,3.75) node[above,text=black] {$\rho_v$};
  \end{scope}

  % Panel (b): Potential W_x(z)
  \begin{scope}[shift={(8.1,0.45)}]
    \node[font=\fontsize{9.5}{11}\selectfont\bfseries,text=ContentsHeading] at (2.8,4.00) {(b) Potential $W_x(z)$};

    % Axes: r-axis at y = 1.85
    \draw[axis] (0,1.85) -- (5.85,1.85) node[right,text=black] {$r$};
    \draw[axis] (0,-0.25) -- (0,3.60) node[left,text=black] {$\phi$};

    % Curve phi(r)
    \draw[AccentTeal!90!black,line width=1.2pt]
      (0,3.25) .. controls (0.95,3.02) and (1.85,2.45) .. (2.85,1.85)
      .. controls (3.55,1.43) and (4.20,0.82) .. (4.95,0.05);
    \node[right,font=\scriptsize,text=AccentTeal!90!black] at (4.90,0.12) {$\phi(r)$};

    % Line through (x, phi(x)) and (lambda, 0)
    \draw[AccentBlue!85,line width=0.95pt] (0,3.46) -- (5.25,0.49);
    \node[font=\scriptsize,text=AccentBlue,align=center] at (4.15,2.75)
      {line through $(x, \phi(x))$\\[-1pt] and $(\lambda, 0)$};
    \draw[ptr,AccentBlue] (3.65,2.45) to[out=-120,in=55] (2.15,2.25);

    % Horizontal gap W_x(z) at z = 4.18 (where curve has y = 0.85, line hits y = 0.85 at r = 4.62)
    \draw[AccentBrown!95!black,line width=1.6pt] (4.18,0.85) -- (4.62,0.85);
    \node[below,font=\scriptsize\bfseries,text=AccentBrown!95!black] at (4.42,0.82) {$W_x(z)$};

    % Dashed projections
    \draw[densely dashed,ContentsHeading!65,line width=0.45pt] (0.85,1.85) -- (0.85,2.98) -- (0,2.98);
    \draw[densely dashed,ContentsHeading!65,line width=0.45pt] (4.18,1.85) -- (4.18,0.85) -- (0,0.85);
    \draw[densely dashed,ContentsHeading!65,line width=0.45pt] (4.62,1.85) -- (4.62,0.85);

    \fill[AccentBlue,draw=white,line width=0.5pt] (0.85,2.98) circle (2.0pt);
    \fill[AccentBlue,draw=white,line width=0.5pt] (2.85,1.85) circle (2.0pt);
    \fill[AccentTeal!90!black,draw=white,line width=0.5pt] (4.18,0.85) circle (2.0pt);
    \fill[white,draw=AccentBlue,line width=0.7pt] (4.62,0.85) circle (2.0pt);

    \node[left,font=\scriptsize,text=ContentsHeading] at (0,2.98) {$\phi(x)$};
    \node[left,font=\scriptsize,text=ContentsHeading] at (0,1.85) {$0$};
    \node[left,font=\scriptsize,text=ContentsHeading] at (0,0.85) {$\phi(z)$};

    \node[below,font=\scriptsize,text=ContentsHeading] at (0.85,1.82) {$x$};
    \node[below,font=\scriptsize,text=ContentsHeading] at (2.85,1.82) {$\lambda$};
    \node[above,font=\scriptsize,text=ContentsHeading] at (4.12,1.87) {$z$};
    \node[above,font=\scriptsize,text=AccentBrown!90!black] at (4.88,1.87) {$z{+}W_x(z)$};
  \end{scope}
\end{tikzpicture}
\caption{Slicing by threshold $t$ and the potential $W_x(z)$~\cite{CaiEtAl26b}. (a)~Every candidate in $S_{t+\Delta_v}$ precedes every candidate outside $S_t$ on voter $v$'s rank line, so $z_v(t + \Delta_v) \le q_v(t) \le s_v$. (b)~For $x < \lambda$ and $z > \lambda$, $W_x(z) = \lambda - z - \frac{\lambda - x}{\phi(x)}\phi(z)$ is the horizontal gap at ordinate $\phi(z)$ between $\phi(r)$ and the line through $(x, \phi(x))$ and $(\lambda, 0)$.}
\label{fig:slicing-and-potential}
\end{figure}

\paragraph{Step 1: Slicing and telescoping.}
For each threshold $t \ge 0$, let $S_t := \{b \in C : \rho(b) > t\}$, and let $p_t := D(S_t)$. On voter $v$'s rank line, a candidate $b$ satisfies $\lowrho_v(b) > t$ if and only if $b$ and every candidate ranked below $b$ belong to $S_t$. Thus, the candidates with $\lowrho_v(b) > t$ occupy an initial interval $[0, q_v(t))$, where $q_v(t) := \min_{c \notin S_t} \rank_v(c; D)$ is the rank of the first candidate outside $S_t$ (even if $D(c) = 0$). Let $z_v(t)$ be the rightmost endpoint of $I_v(b)$ over all candidates $b \in S_t \cap \supp(D)$ (with $z_v(t) = 0$ if $p_t = 0$), and let $s_v := \rank_v(a^\star; D)$ be the left endpoint of $a^\star$'s interval. Since $\rho(a^\star) = 0$, we have $a^\star \notin S_t$ and therefore $q_v(t) \le s_v$. Observe that $q_v(t)$ and $z_v(t)$ satisfy (\cref{fig:slicing-and-potential}(a))
\begin{equation}
  \int_0^\infty q_v(t)\,\mathrm{d}t \;=\; N_v(D),
  \qquad
  z_v(t + \Delta_v) \;\le\; q_v(t).
  \label{eq:slice-identities}
\end{equation}
The first relation integrates $\Pr_{b \sim D}[\lowrho_v(b) > t] = q_v(t)$ over $t \ge 0$. For the second, if some candidate $a \in S_{t + \Delta_v}$ were ranked by $v$ above some candidate $c \notin S_t$, then $\rho(a) > t + \Delta_v$ and $\lowrho_v(a) \le \rho(c) \le t$, contradicting $\Delta_v = \max_{b \in C}(\rho(b) - \lowrho_v(b))$.

To bound $\int_0^\infty q_v(t)\,\mathrm{d}t$, suppose we modify $q_v(t)$ into $L_v(t)$ satisfying $\E_v[q_v(t)] \le \E_v[L_v(t)]$ for every $t \ge 0$, and find a nondecreasing function $W \colon [0, 1] \to \R_{\ge 0}$ such that for every voter $v$ and threshold $t \ge 0$,
\begin{equation}
  \begin{aligned}
    L_v(t) &\;\le\; \lambda - W\bigl(q_v(t)\bigr), \\
    L_v(t) &\;\le\; W(z) - W\bigl(q_v(t)\bigr) \qquad \text{whenever } z_v(t) \le z \le s_v.
  \end{aligned}
  \label{eq:telescoping-target}
\end{equation}
Extend $V(t) := W(q_v(t))$ to $t < 0$ by setting $V(t) := \lambda$. For $0 \le t < \Delta_v$, the first inequality in \cref{eq:telescoping-target} gives $L_v(t) \le V(t - \Delta_v) - V(t)$. For $t \ge \Delta_v$, \cref{eq:slice-identities} implies $z_v(t) \le q_v(t - \Delta_v) \le s_v$, so applying the second inequality in \cref{eq:telescoping-target} with $z = q_v(t - \Delta_v)$ again gives $L_v(t) \le V(t - \Delta_v) - V(t)$. Integrating up to $X := \max_{b \in \supp(D)} \rho(b)$ (above which $p_t = 0$) telescopes $V(t)$ across thresholds separated by $\Delta_v$:
\[
  \int_0^X L_v(t)\,\mathrm{d}t
  \;\le\; \int_{-\Delta_v}^0 V(t)\,\mathrm{d}t - \int_{X - \Delta_v}^X V(t)\,\mathrm{d}t
  \;\le\; \lambda\Delta_v,
\]
since $V(t) = \lambda$ for $t < 0$ and $V(t) \ge 0$ for $t \ge 0$. Taking expectations over $v \sim V$ and using $\E_v[q_v(t)] \le \E_v[L_v(t)]$ gives $\E_v[N_v(D)] \le \lambda\,\E_v[\Delta_v]$.

\paragraph{Step 2: Using stability to offset each slice.}
To build the offset in $L_v(t)$, we use two consequences of the $g_\phi$-stability of $D$. First, by \cref{lem:stationarity}, every candidate in $\supp(D)$ has expected centered activation zero across voters, so the centered activation $k_v(t) := -\sum_{b \in S_t} \int_{I_v(b)} \phi(r)\,\mathrm{d}r$ of $S_t$ on voter $v$'s rank line satisfies
\begin{equation}
  \E_{v \sim V}\bigl[k_v(t)\bigr] \;=\; 0.
  \label{eq:kv-zero-mean}
\end{equation}
Second, since $\phi$ is nonincreasing, its value at the left endpoint $s_v$ of $a^\star$'s interval is at least its average over $I_v(a^\star)$; as $a^\star$ has expected centered activation at most zero, this gives
\begin{equation}
  \E_{v \sim V}\bigl[\phi(s_v)\bigr] \;\ge\; 0.
  \label{eq:sv-nonneg-mean}
\end{equation}
Let $A := \E_{v \sim V}\bigl[\1_{\{s_v < \lambda\}}\,\phi(s_v)\bigr]$, and assume first that $A > 0$ and $0 < p_t < 1$. Given $\beta(p) \ge 0$ and $c(p, x) \ge 0$ for $p \in (0, 1)$ and $x \in [0, \lambda)$, we define
\begin{equation}
  L_v(t) \;:=\; q_v(t) + \beta(p_t)\,k_v(t) +
  \begin{cases}
    c(p_t, s_v), & s_v < \lambda, \\
    \phi(s_v)\,\overline{c}(p_t), & s_v \ge \lambda,
  \end{cases}
  \qquad
  \overline{c}(p) \;:=\; \frac{\E_{v \sim V}\bigl[\1_{\{s_v < \lambda\}}\,c(p, s_v)\bigr]}{A}.
  \label{eq:Lv-def}
\end{equation}
The scaling in $\overline{c}(p_t)$ ensures that the last term in \cref{eq:Lv-def} has expectation $\overline{c}(p_t)\,A + \overline{c}(p_t)\,\E_v[\1_{\{s_v \ge \lambda\}}\,\phi(s_v)] = \overline{c}(p_t)\,\E_v[\phi(s_v)] \ge 0$ by \cref{eq:sv-nonneg-mean}, while $\beta(p_t)\,k_v(t)$ has expectation zero by \cref{eq:kv-zero-mean}. Thus $\E_v[q_v(t)] \le \E_v[L_v(t)]$.

\paragraph{Step 3: Choosing the potential.}
For each $x \in [0, \lambda)$, we define the potential (\cref{fig:slicing-and-potential}(b))
\begin{equation}
  W_x(z) \;:=\;
  \begin{cases}
    0, & z \le \lambda, \\
    \lambda - z - \dfrac{\lambda - x}{\phi(x)}\,\phi(z), & z > \lambda,
  \end{cases}
  \label{eq:potential-def}
\end{equation}
and average $W_{s_v}(z)$ over the voters with $s_v < \lambda$, weighted by $\phi(s_v)$ (just like $\overline{c}(p)$ in \cref{eq:Lv-def}):
\begin{equation}
  W(z) \;:=\; \frac{\E_{v \sim V}\bigl[\1_{\{s_v < \lambda\}}\,\phi(s_v)\,W_{s_v}(z)\bigr]}{A}.
  \label{eq:W-avg-def}
\end{equation}
Because $0 < \phi(x) \le \lambda - x$, the line through $(x, \phi(x))$ and $(\lambda, 0)$ has slope at least $-1$. Since $\phi'(z) \le -1$ on $[\lambda, 1]$, the curve $\phi$ falls at least as steeply as this line beyond $\lambda$, so the horizontal gap $W_x(z)$ between them (\cref{fig:slicing-and-potential}(b)) is nonnegative and nondecreasing in $z$, and satisfies the mass-one identity
\[
  \phi(x)\bigl[z + W_x(z) - \lambda\bigr] \;=\; \bigl[-\phi(z)\bigr](\lambda - x)
  \qquad (z \ge \lambda).
\]

\paragraph{Step 4: Checking each voter.}
To bound $L_v(t)$ for an individual voter $v$, we first bound $k_v(t)$ in terms of $p = p_t$, $q = q_v(t)$, and $z \ge z_v(t)$. On $v$'s rank line, $S_t$ completely fills $[0, q)$, which contributes $-J(q)$ to $k_v(t)$, while the other $p - q$ mass of $S_t$ lies inside $[q, z]$. Since $-\phi$ is nondecreasing, this mass attains its largest centered activation when pushed as far right as allowed---either to $[1 - p + q, 1]$ or against $z$ on $[z - p + q, z]$:
\begin{equation}
  \begin{aligned}
    k_v(t) &\;\le\; \delta_p(q) \;:=\; J(1 - p + q) - J(q), \\
    k_v(t) &\;\le\; K_p(q, z) \;:=\; J(z - p + q) - J(z) - J(q).
  \end{aligned}
  \label{eq:kv-packing-bounds}
\end{equation}
Notice that $\delta_p(q)$ is the centered activation of $\calH$ for Type~1 voters in Election~2, and $K_p(0, z)$ is the centered activation of $\calH$ for Type~2 voters in Election~3. We now check \cref{eq:telescoping-target} voter by voter when $0 < p < 1$ and $A > 0$, with $q = q_v(t)$ and $z \in [z_v(t), s_v]$:
\begin{itemize}
  \item \emph{For a voter with $s_v = x < \lambda$:} Here $q \le z \le x < \lambda$, so $W(q) = W(z) = 0$ and $K_p(q, z) \le K_p(q, x)$ since $K_p(q, \cdot)$ is nondecreasing. By \cref{eq:Lv-def,eq:kv-packing-bounds}, \cref{eq:telescoping-target} holds whenever
  \[
    q + \beta(p)\,\delta_p(q) + c(p, x) \;\le\; \lambda \;\; \bigl(0 \le q \le \min\{p, x\}\bigr),
    \quad
    q + \beta(p)\,K_p(q, x) + c(p, x) \;\le\; 0 \;\; (0 \le q \le p \le x).
  \]
  \item \emph{For a voter with $s_v \ge \lambda$:} Here $q \le z \le s_v$ and $\phi(s_v) \le 0$, so $[-\phi(q)]_+ \le -\phi(s_v)$ and $[-\phi(z)]_+ \le -\phi(s_v)$. By \cref{eq:Lv-def,eq:kv-packing-bounds}, \cref{eq:telescoping-target} holds for this voter whenever
  \[
    q + \beta(p)\,\delta_p(q) + W(q) - \lambda \;\le\; \bigl[-\phi(q)\bigr]_+\,\overline{c}(p),
    \qquad
    q + \beta(p)\,K_p(q, z) + W(q) - W(z) \;\le\; \bigl[-\phi(z)\bigr]_+\,\overline{c}(p).
  \]
  By \cref{eq:Lv-def,eq:W-avg-def}, these two inequalities follow by averaging over the voters with $s_v = x < \lambda$ (with weights $\phi(x)/A$) whenever \cref{eq:ver-3,eq:ver-4} below hold for every $x \in [0, \lambda)$.
\end{itemize}
The boundary cases $p_t \in \{0, 1\}$ and $A = 0$ follow from the same inequalities and are checked in \cref{app:verifier-proof}. Summarizing Steps~1--4 proves the following lemma~\cite{CaiEtAl26b}.

\begin{lemma}[Verifier Lemma]
\label{lem:verifier}
Let $\phi \colon [0, 1] \to \R$ be continuous and strictly decreasing with $\int_0^1 \phi(r)\,\mathrm{d}r = 0$, $\phi(\lambda) = 0$, $\phi(1) = -1$, and $\phi(r) \le \lambda - r$ on $[0, 1]$, and suppose $\phi$ is absolutely continuous on $[\lambda, 1]$ with $\phi' \le -1$ almost everywhere there. Suppose that for every $p \in (0, 1)$ there is a finite coefficient $\beta(p) \ge 0$ and, for each $x \in [0, \lambda)$, a finite correction $c(p, x) \ge 0$ such that
\begin{align}
  q + \beta(p)\,\delta_p(q) + c(p, x) &\;\le\; \lambda
    && \bigl(0 \le q \le \min\{p, x\}\bigr), \label{eq:ver-1} \\
  q + \beta(p)\,K_p(q, x) + c(p, x) &\;\le\; 0
    && (0 \le q \le p \le x < \lambda), \label{eq:ver-2} \\
  \phi(x)\bigl[q + \beta(p)\,\delta_p(q) + W_x(q) - \lambda\bigr] &\;\le\; \bigl[-\phi(q)\bigr]_+\,c(p, x)
    && (0 \le q \le p), \label{eq:ver-3} \\
  \phi(x)\bigl[q + \beta(p)\,K_p(q, z) + W_x(q) - W_x(z)\bigr] &\;\le\; \bigl[-\phi(z)\bigr]_+\,c(p, x)
    && (0 \le q \le p \le z \le 1). \label{eq:ver-4}
\end{align}
Then every $g_\phi$-stable lottery has distortion at most $1 + 2\lambda$.
\end{lemma}

\subsection{Verifier Coefficients from Supporting Lines}
\label{sec:duality-coefficients}

We now show why a minimizing function $\phi$ of the lower-bound variational problem satisfies the upper-bound inequalities of \cref{lem:verifier} at the same parameter $\lambda = \lambda^\star$.

Fix a mass $p \in [p_-, p_+]$ used in Election~3, and write $R = J(p)$ and $P = J(1 - p)$. Look first at the high-rank verifier inequality~\eqref{eq:ver-4} at $q = 0$ and $z > \lambda$: dividing by $\phi(x) > 0$ (where $W_x(0) = 0$, $K_p(0, z) = J(z - p) - J(z)$, and $W_x(z) = \lambda - z - \frac{\lambda - x}{\phi(x)}\phi(z)$) gives
\[
  z - \lambda + \beta\bigl[J(z - p) - J(z)\bigr] + \frac{\lambda - x + c(p, x)}{\phi(x)}\,\phi(z) \;\le\; 0.
\]
Since $\phi(z) < 0$, making the coefficient $\frac{\lambda - x + c(p, x)}{\phi(x)}$ at least a constant $\eta > 0$ across all $x \in [0, \lambda)$ requires $c(p, x) \ge \eta\phi(x) - (\lambda - x)$. Choosing the smallest nonnegative correction
\[
  c(p, x) \;:=\; \bigl[\eta\phi(x) - (\lambda - x)\bigr]_+
  \qquad (0 \le x < \lambda)
\]
gives $\frac{\lambda - x + c(p, x)}{\phi(x)} = \max\bigl\{\frac{\lambda - x}{\phi(x)}, \eta\bigr\} \ge \eta$, reducing the high-rank inequality at $q = 0$ to the $x$-independent bound
\begin{equation}
  z - \lambda + \beta\bigl[J(z - p) - J(z)\bigr] + \eta\phi(z) \;\le\; 0.
  \label{eq:high-q0-target}
\end{equation}

Next, consider the two low-rank verifier inequalities \cref{eq:ver-1,eq:ver-2} at their tightest split ranks $q = 0$ and $q = p$:
\begin{itemize}
  \item At $q = 0$, we have $\delta_p(0) = J(1 - p) = P$, where the entire mass $p$ of $S_t$ lies at the top of the rank line on $[1 - p, 1)$. Thus \cref{eq:ver-1} requires $\beta P + c(p, x) \le \lambda$. When $c(p, x) > 0$, substituting $c(p, x) = \eta\phi(x) - (\lambda - x)$ turns this condition into
  \[
    x + \beta P + \eta\phi(x) \;\le\; 2\lambda.
  \]
  \item At $q = p \le y < \lambda$, we have $K_p(p, y) = -J(p) = -R$, where the entire mass $p$ of $S_t$ lies at the bottom of the rank line on $[0, p)$ below $a^\star$ at rank $y$. Thus \cref{eq:ver-2} requires $p - \beta R + c(p, y) \le 0$. When $c(p, y) > 0$, substituting $c(p, y) = \eta\phi(y) - (\lambda - y)$ turns this condition into
  \[
    y + p - \beta R + \eta\phi(y) \;\le\; \lambda.
  \]
\end{itemize}
Evaluating these two displayed inequalities at $x = \lambda$ and $y = \lambda$ (where $\phi(\lambda) = 0$) also gives $\beta P \le \lambda$ and $\beta R \ge p$, which cover the cases where $c(p, x) = 0$ or $c(p, y) = 0$.

To make both low-rank inequalities hold simultaneously across all $x$ and $y$ with a single coefficient $\beta$, we eliminate $\beta$ by multiplying the first inequality by $R$ and the second by $P$. We then choose $\eta > 0$ so that the weighted sum of their left-hand maxima matches the right-hand side:
\begin{equation}
  R \max_{x}\bigl\{x + \eta\phi(x)\bigr\} + P \max_{y \ge p}\bigl\{y + \eta\phi(y)\bigr\}
  \;=\; 2\lambda R + (\lambda - p)P.
  \label{eq:eta-balance}
\end{equation}
Initially, let the maxima in \cref{eq:eta-balance} run over the compact low-rank intervals $x \in [x_-, x_+]$ and $y \in [\max\{x_-, p\}, x_+]$ used in the variational problem. Because $\phi > 0$ on $[x_-, x_+]$, the left-hand side of \cref{eq:eta-balance} is continuous and strictly increasing in $\eta$. At $\eta = 0$, it is strictly smaller than the right-hand side because $\Delta = (2\lambda - x)R - (p + y - \lambda)P > 0$ for all $x, y \in [x_-, x_+]$; as $\eta \to \infty$, it tends to $\infty$. Hence \cref{eq:eta-balance} has a unique solution $\eta > 0$. Setting
\begin{equation}
  \beta \;:=\; \frac{\max_{y \ge p}\{y + \eta\phi(y)\} - \lambda + p}{R},
  \label{eq:beta-def}
\end{equation}
both low-rank inequalities hold for all $x$ and $y$ in the tested intervals, with equality at the maximizing positions. Furthermore, the slope bounds of \cref{prop:shape} ensure that $\beta > 0$ and that extending the maxima to all $x \in [0, \lambda)$ and $y \in [p, \lambda)$ does not increase either maximum, because $r + \eta\phi(r)$ increases before $[x_-, x_+]$ and decreases after it.

Now let $x$ and $y$ be ranks attaining the two maxima in \cref{eq:eta-balance,eq:beta-def}, so that
\[
  \phi(x) \;=\; \frac{2\lambda - x - \beta P}{\eta}
  \qquad\text{and}\qquad
  \phi(y) \;=\; \frac{\lambda - p - y + \beta R}{\eta}.
\]
Substituting these two expressions into the combinations of $\phi(x)$ and $\phi(y)$ that appear in the three-class inequality~\eqref{eq:three-class} of Election~3 cancels the $-\beta R P$ and $+\beta P R$ terms in the first sum and the $(2\lambda - x)(p + y - \lambda)$ terms in the second:
\[
  R\phi(x) + P\phi(y) \;=\; \frac{(2\lambda - x)R - (p + y - \lambda)P}{\eta} \;=\; \frac{\Delta}{\eta},
\]
\[
  (p + y - \lambda)\phi(x) + (2\lambda - x)\phi(y) \;=\; \frac{\beta\bigl[(2\lambda - x)R - (p + y - \lambda)P\bigr]}{\eta} \;=\; \frac{\beta\Delta}{\eta}.
\]
With $K = J(z - p) - J(z)$, the three-class inequality~\eqref{eq:three-class} of Election~3 is
\[
  (z - \lambda)\bigl[R\phi(x) + P\phi(y)\bigr] + K\bigl[(p + y - \lambda)\phi(x) + (2\lambda - x)\phi(y)\bigr] + \Delta\,\phi(z) \;\le\; 0.
\]
Substituting the two identities above factors out $\Delta / \eta > 0$:
\[
  \frac{\Delta}{\eta}\,\bigl[z - \lambda + \beta K + \eta\phi(z)\bigr] \;\le\; 0.
\]
Dividing by $\Delta / \eta > 0$ gives \cref{eq:high-q0-target}. Thus, the two independent low-rank positions $x, y$ in Election~3 determine a single pair $(\beta(p), \eta)$ that satisfies the low-rank and high-rank verifier inequalities simultaneously. For masses $p$ outside $[p_-, p_+]$ and for positive initial block lengths $q > 0$, monotonicity and the slope bounds of \cref{prop:shape} reduce \cref{eq:ver-1,eq:ver-2,eq:ver-3,eq:ver-4} to the finite checks in \cref{app:upper-all}.

\subsection{Completing \texorpdfstring{$\phi^\star$}{phi*} and Proving Theorem 1.1}
\label{sec:completing-score}

The variational problem determines the minimizing function $\phi^\star$ and its primitive $J^\star$ on $[p_-, z_1]$. To extend $\phi^\star$ to $[0, 1]$ with $J^\star(0) = J^\star(1) = 0$ and $\phi^\star(1) = -1$, we join it to $r = 0$ and $r = 1$ with two simple polynomial pieces (\cref{fig:supporting-lines} and \cref{app:joining-curve}):
\begin{itemize}
  \item \textbf{Near $r = 0$:} Because $C_{\lambda^\star, S}(0) = \lambda^\star(\lambda^\star - \frac{1}{2} - S) > 0$, the rational curve $C_{\lambda^\star, S}$ does not pass through the origin. At $x_h = 0.2635 \in (p_-, \ell)$, we replace $C'_{\lambda^\star, S}$ on $[0, x_h]$ by the unique affine function $\phi^\star(r) = \phi_0 + s_h r$ whose primitive starts at $J^\star(0) = 0$ and matches both $J^\star(x_h) = C_{\lambda^\star, S}(x_h)$ and $\phi^\star(x_h) = C'_{\lambda^\star, S}(x_h)$.
  \item \textbf{Near $r = 1$:} By \cref{prop:attainment} and the boundary exclusions in \cref{app:checking-binary}, the continuation $Y(z)$ touches $T_{\lambda^\star, S}(z, q)$ at an interior point $(z_\star, q_\star) \in (z_B, z_1) \times \calQ$, where both first-order partial derivatives of $Y(z) - T_{\lambda^\star, S}(z, q)$ vanish. On $[z_\star, 1]$, we complete $\phi^\star$ by the unique quadratic polynomial that matches $\phi^\star(z_\star) = Y'(z_\star)$, ends at $\phi^\star(1) = -1$, and has integral $-Y(z_\star)$ so that $J^\star(1) = 0$.
\end{itemize}

\begin{lemma}[Verification of the completed function $\phi^\star$]
\label{lem:completed-score}
The completed function $\phi^\star \colon [0, 1] \to \R$ is continuous and strictly decreasing, with $\int_0^1 \phi^\star(r)\,\mathrm{d}r = 0$, $\phi^\star(\lambda^\star) = 0$, $\phi^\star(1) = -1$, and $\phi^\star(r) \le \lambda^\star - r$ on $[0, 1]$, and is concave on $[\lambda^\star, 1]$ with $(\phi^\star)' \le -1$ almost everywhere there. Furthermore, for every $p \in (0, 1)$, there exist $\beta(p) \ge 0$ and $c(p, x) \ge 0$ ($0 \le x < \lambda^\star$) satisfying \cref{eq:ver-1,eq:ver-2,eq:ver-3,eq:ver-4} of \cref{lem:verifier}.
\end{lemma}

\cref{app:upper-all} proves \cref{lem:completed-score} by extending the coefficient construction of \cref{sec:duality-coefficients} across all $p \in (0, 1)$, and \cref{app:numerical} collects the accompanying interval-arithmetic verifications.

\begin{proof}[Proof of {\cref{thm:main}}]
By \cref{lem:completed-score} and \cref{lem:verifier}, the activation $g^\star = g_{\phi^\star} \in \calG$ is continuous and strictly increasing and satisfies $M(g^\star) \le 1 + 2\lambda^\star$.

Conversely, suppose some activation $g \in \calG$ had $M(g) < 1 + 2\lambda^\star$. Choose $\lambda \in [\lambda_0, \lambda^\star)$ with $M(g) \le 1 + 2\lambda$ (taking $\lambda = \lambda_0$ if $M(g) \le 1 + 2\lambda_0$). Then \cref{prop:reduction} produces a pair $(S, \phi)$ satisfying all conditions in the definition~\eqref{eq:lambda-star-def} of $\lambda^\star$ at parameter $\lambda < \lambda^\star$, contradicting the definition of $\lambda^\star$. Therefore,
\[
  M(g^\star) \;=\; \min_{g \in \calG} M(g) \;=\; 1 + 2\lambda^\star.
\]
Finally, \cref{lem:numerical-calibration} in \cref{app:calibration} certifies $0.56232844 < \lambda^\star < 0.56232847$, so $1 + 2\lambda^\star \approx 2.1246569$.
\end{proof}

\appendix
\crefalias{section}{appendix}

\section{Proofs for \texorpdfstring{$g$}{g}-Stable Lotteries and the Three Elections}
\label{app:elections}

\subsection{Existence of \texorpdfstring{$g$}{g}-Stable Lotteries and the Biased Metric Construction}
\label{app:existence}

\begin{proof}[Proof of existence in {\cref{lem:stationarity}}]
Let $g \colon [0, 1] \to [0, 1]$ be any continuous function with mean $\tau = \int_0^1 g(r)\,\mathrm{d}r$. To find a $g$-stable lottery, we increase the probabilities of candidates whose expected activations exceed $\tau$. For a lottery $D$ on the simplex $\Delta(C)$, write
\[
  e_a(D) \;:=\; \bigl[\E_{v \sim V}\psi_v(a; D) - \tau\bigr]_+,
  \qquad
  e(D) \;:=\; \sum_{a \in C} e_a(D),
\]
and consider the map $T \colon \Delta(C) \to \Delta(C)$ defined by
\[
  T_a(D) \;:=\; \frac{D(a) + e_a(D)}{1 + e(D)}
  \qquad (a \in C).
\]
Continuity of $g$ makes each function $D \mapsto \psi_v(a; D)$, and hence $T$, continuous on $\Delta(C)$. By Brouwer's fixed-point theorem, $T$ has a fixed point $D$. If $e(D) > 0$ at this fixed point, then $D(a)\,e(D) = e_a(D)$ implies $e_a(D) > 0$ for every $a \in \supp(D)$, so every supported candidate has $\E_v[\psi_v(a; D)] > \tau$, contradicting $\sum_{a \in C} D(a)\,\E_v[\psi_v(a; D)] = \tau$. Thus $e(D) = 0$, which means $e_a(D) = 0$ for all $a \in C$ and $D$ is $g$-stable.
\end{proof}

\begin{proof}[Proof of the converse in {\cref{lem:biased-reduction}}]
Fix $a^\star \in C$ and any distance function $\rho \colon C \to \R_{\ge 0}$ with $\rho(a^\star) = 0$. Define the distance between any two distinct candidates $a, b \in C$ to be $\max\{\rho(a), \rho(b)\}$ (and $0$ if $a = b$), and define the distance between voter $v \in V$ and candidate $b \in C$ to be
\[
  d(v, b) \;:=\; \frac{1}{2}\Delta_v + \lowrho_v(b).
\]
To check that these distances extend to a pseudometric on $V \cup C$, write $u_b := \lowrho_v(b)$ for a fixed voter $v$. By definition of $\lowrho_v$ and $\Delta_v = \max_{c \in C}(\rho(c) - \lowrho_v(c))$, we have $0 \le u_b \le \rho(b) \le \Delta_v + u_b$ for every $b \in C$, and hence
\[
  |u_a - u_b| \;\le\; \max\{\rho(a), \rho(b)\} \;\le\; \Delta_v + u_a + u_b
  \;=\; d(v, a) + d(v, b).
\]
Thus, no two-step path through a voter $v$ is shorter than the direct distance $\max\{\rho(a), \rho(b)\}$ between candidates $a$ and $b$, and no path through candidates is shorter than $d(v, b)$. Setting the distance between any two voters to be their shortest-path distance in the graph of these prescribed edges therefore gives a pseudometric $d$ on $V \cup C$ that preserves all prescribed voter--candidate and candidate--candidate distances.

Because $\lowrho_v(b)$ is nonincreasing as $b$ moves up voter $v$'s ranking, $d$ is consistent with the election. Since $\lowrho_v(a^\star) = \rho(a^\star) = 0$, candidate $a^\star$ has social cost $\cost_d(a^\star) = \frac{1}{2}\E_v[\Delta_v]$, every candidate $b \in C$ has $\cost_d(b) \ge \cost_d(a^\star)$, and
\[
  \cost_d(D) \;=\; \frac{1}{2}\E_{v \sim V}\bigl[\Delta_v\bigr] + \E_{v \sim V}\bigl[N_v(D)\bigr].
\]
Hence, any strict violation of \cref{eq:height-inequality} gives $\cost_d(D) > (1 + 2\lambda)\cost_d(a^\star)$. Adding $\varepsilon > 0$ to every distance between distinct points produces a genuine metric that remains consistent with the reported rankings; as $\varepsilon \downarrow 0$, its cost ratio converges to $1 + 2\E_v[N_v(D)] / \E_v[\Delta_v]$ when $\E_v[\Delta_v] > 0$, and to $\infty$ when $\E_v[\Delta_v] = 0 < \E_v[N_v(D)]$.
\end{proof}

\subsection{Normalizing \texorpdfstring{$\phi$}{phi} and \texorpdfstring{$J$}{J}}
\label{app:normalization}

\begin{proof}[Proof of {\cref{lem:election-1}}]
In Election~1 (\cref{sec:election-1}), if all voters rank $a^\star \notin \supp(D)$ at the top ($r = 1$) and permute $m$ equal-mass candidates in $\supp(D)$ uniformly, every supported candidate has expected activation $\tau$, while $a^\star$ has activation $g(1)$, $N_v = 1$, and $\Delta_v = 0$. Since $M(g) \le 1 + 2\lambda < \infty$, this lottery cannot be $g$-stable, so $g(1) > \tau = \int_0^1 g(r)\,\mathrm{d}r$ (and in particular $g$ is nonconstant). Set $\widetilde{g}(r) := \frac{g(r) - \tau}{g(1) - \tau}$, so that $\int_0^1 \widetilde{g}(r)\,\mathrm{d}r = 0$ and $\widetilde{g}(1) = 1$.

In the construction of Election~1 (\cref{sec:election-1}), any finitely supported random rank $R \in [0, 1]$ for $a^\star$ must satisfy the implication
\begin{equation}
  \E\bigl[\widetilde{g}(R)\bigr] \;\le\; 0 \quad\Longrightarrow\quad \E\bigl[f_\lambda(R)\bigr] \;\le\; 0.
  \label{eq:f-lambda-implication}
\end{equation}
Indeed, if a finitely supported distribution on $[0, 1]$ violated \cref{eq:f-lambda-implication}, we could first mix in a small probability mass at a rank $r_0$ with $\widetilde{g}(r_0) < 0$ so that $\E[\widetilde{g}(R)] < 0$ and $\E[f_\lambda(R)] > 0$ both hold strictly. Approximating all ranks in $\supp(R) \setminus \{1\}$ by rational numbers with common denominator $m$ (while keeping any atom at $r = 1$ fixed) preserves both strict inequalities by continuity of $g$ and of $f_\lambda$ on $[0, 1)$. Approximating the voter-type weights by rational numbers and taking $m$ equal-mass candidates $b \in \supp(D)$ uniformly permuted within each voter type keeps every supported candidate's expected activation equal to $\tau$ while $a^\star$ has expected activation strictly below $\tau$, giving a finite election that violates \cref{lem:biased-reduction}. Thus \cref{eq:f-lambda-implication} holds for all finitely supported random variables $R$ on $[0, 1]$.

Now take any two ranks $s, r \in [0, 1]$ with $\widetilde{g}(s) < 0 < \widetilde{g}(r)$, and mix them with probabilities proportional to $\widetilde{g}(r)$ and $-\widetilde{g}(s)$ so that $\E[\widetilde{g}(R)] = 0$. Then \cref{eq:f-lambda-implication} gives
\[
  \frac{f_\lambda(r)}{\widetilde{g}(r)} \;\le\; \frac{f_\lambda(s)}{\widetilde{g}(s)}.
\]
Hence the supremum
\[
  c \;:=\; \sup_{\substack{r \in [0, 1] \\ \widetilde{g}(r) > 0}} \frac{f_\lambda(r)}{\widetilde{g}(r)}
\]
is finite (bounded above by $f_\lambda(s)/\widetilde{g}(s)$ for any fixed $s$ with $\widetilde{g}(s) < 0$), and evaluating at $r = 1$ gives $c \ge f_\lambda(1)/\widetilde{g}(1) = 1$. The same inequality shows $c\,\widetilde{g}(s) \ge f_\lambda(s)$ whenever $\widetilde{g}(s) < 0$; the definition of $c$ gives $c\,\widetilde{g}(r) \ge f_\lambda(r)$ whenever $\widetilde{g}(r) > 0$; and taking a point mass in \cref{eq:f-lambda-implication} gives $0 = c\,\widetilde{g}(r) \ge f_\lambda(r)$ whenever $\widetilde{g}(r) = 0$.

Define $\phi(r) := -c\,\widetilde{g}(r) = \frac{c}{g(1) - \tau}(\tau - g(r))$. Then $\phi$ is continuous with $\int_0^1 \phi(r)\,\mathrm{d}r = 0$, $\phi(1) = -c \le -1$, and $\phi(r) \le -f_\lambda(r) = \lambda - r$ for all $r \in [0, 1)$ (and at $r = 1$ by continuity). Since $J(1) = 0$ and $J'(r) = \phi(r) \le \lambda - r < 0$ on $(\lambda, 1]$, we have $J(r) > 0$ for all $r \in [\lambda, 1)$ (and $\lambda > 1/2$ since $\int_0^1 (\lambda - r)\,\mathrm{d}r > \int_0^1 \phi(r)\,\mathrm{d}r = 0$). If $\{r \in (0, 1) : J(r) \le 0\}$ were nonempty, let $r_0 \in (0, \lambda)$ be its maximum, so $J(r_0) = 0$ and $J(p) > 0$ on $(r_0, 1)$. For every $p \in (r_0, 1)$ and $q \in [0, r_0]$, the positive-$J(p)$ case of the Election~2 construction (\cref{sec:election-2}) applies since $J(p) > 0$ and $q < \min\{p, \lambda\}$, giving $p[J(1 - p + q) - J(q)] \le (\lambda - q)J(p)$. If $\min_{[0, 1]} J < 0$, taking $q = p_\star := \max(\arg\min_{[0, 1]} J) \in (0, r_0)$ and letting $p \downarrow r_0$ gives $J(1 - r_0 + p_\star) \le J(p_\star)$, contradicting $1 - r_0 + p_\star \in (p_\star, 1)$; if $\min_{[0, 1]} J = 0$, then $J(r_0) = J'(r_0) = 0$, and dividing $p[J(1 - p + r_0) - J(r_0)] \le (\lambda - r_0)J(p)$ by $p - r_0 > 0$ as $p \downarrow r_0$ gives $r_0(-\phi(1)) \le (\lambda - r_0)J'(r_0) = 0$, contradicting $-\phi(1) \ge 1$. Thus $J(r) > 0$ for all $r \in (0, 1)$.
\end{proof}

\subsection{Finite Elections for Elections 2 and 3}
\label{app:finite-two-three}

For Election~2 (\cref{sec:election-2}), the endpoint $q=p$ is immediate from $J(1)-J(p)=-J(p)$; equivalently, it follows by continuity from $q<p$. Any strict violation of \cref{eq:two-class} persists after approximating the class mass $p$ and split rank $q$ by rational numbers with a large odd common denominator $m$. Realizing the two classes $\calH$ and $\calZ$ with $m$ equal-mass candidates in total and recomputing the balancing voter-type weights from the values of $J$ at the rational ranks gives an exactly $g$-stable weighted election with the same strict violation of \cref{eq:height-inequality}. Applying \cref{lem:equal-weights} below converts this weighted election into a finite election with equally weighted voters.

\begin{proof}[Proof of {\cref{lem:election-3}}]
Integrating $\phi(r) \le \lambda - r$ over $[z - p, z]$ gives
\[
  K \;=\; J(z - p) - J(z) \;=\; -\int_{z-p}^z \phi(r)\,\mathrm{d}r \;\ge\; p\Bigl(z - \lambda - \frac{p}{2}\Bigr) \;>\; 0,
\]
since $z > \lambda + p/2$. For $\Delta = (2\lambda - x)R - (p + y - \lambda)P$, if $p + y - \lambda \le 0$, then $\Delta > 0$ follows from $R, P > 0$ and $x < \lambda$. Otherwise $p + y - \lambda > 0$, and the $q = 0$ two-class inequality $p P \le \lambda R$ from \cref{eq:two-class-zero} gives
\[
  \frac{\Delta}{R} \;\ge\; 2\lambda - x - (p + y - \lambda)\frac{\lambda}{p}
  \;=\; \lambda - x + \frac{\lambda(\lambda - y)}{p} \;>\; 0,
\]
since $x < \lambda$ and $y < \lambda$.

Write $E_3$ for the left-hand side of \cref{eq:three-class}, and recall the voter-type weights $w_1, w_2, w_3$ from \cref{eq:three-class-weights}, which satisfy $R w_3 = K w_2 + P w_1$ and $E_3 = -(2\lambda - x)w_1 + (z - \lambda)w_2 + (p + y - \lambda)w_3$:
\begin{itemize}
  \item \emph{Case 1: $w_1 \le 0$.} Using only Type~2 and Type~3 voters from \cref{sec:election-3} with positive weights $R$ and $K$ gives $\calH$ centered activation $R K - K R = 0$ and gives $a^\star$ centered activation $-R\phi(z) - K\phi(y) = w_1 \le 0$. Their weighted contribution to $N_v(D) - \lambda\Delta_v$ is $(z - \lambda)R + (p + y - \lambda)K > 0$, which contradicts \cref{lem:biased-reduction}.
  \item \emph{Case 2: $w_1 > 0$ and $w_2 \le 0$.} Substituting $w_3 = (K w_2 + P w_1)/R$ into $E_3 = -(2\lambda - x)w_1 + (z - \lambda)w_2 + (p + y - \lambda)w_3$ gives
  \[
    E_3 \;=\; -\frac{\Delta}{R}\,w_1 + \Bigl(z - \lambda + \frac{(p + y - \lambda)K}{R}\Bigr)w_2 \;<\; 0,
  \]
  so \cref{eq:three-class} holds strictly.
  \item \emph{Case 3: $w_1 > 0$ and $w_2 > 0$.} Then $w_3 = (K w_2 + P w_1)/R > 0$. All three voter-type weights are positive, and their weighted contribution to $N_v(D) - \lambda\Delta_v$ is $E_3$, so $E_3 > 0$ would again contradict \cref{lem:biased-reduction}.
\end{itemize}

To turn either strict contradiction (in Case~1 or Case~3) into a finite election with equally weighted voters, we first make the unsupported candidate $a^\star$ strictly stable. Choose $0<r_0<\min\{p,1-p\}$ with $\phi(r_0)>0$. Such a rank exists because $J(0)=0$ and $J(t)>0$ for $0<t<1$. Use two auxiliary voter types with $\calH$ at the top for the first type and at the bottom for the second. Insert the zero-mass candidate $a^\star$ at rank $r_0$ in both types, splitting the block that contains it. This leaves the centered activations of $\calH$ equal to $P$ and $-R$. Weights proportional to $R$ and $P$ balance $\calH$ and give $a^\star$ total centered activation $-(R+P)\phi(r_0)<0$. Adding a small enough positive multiple of this pair keeps $\E_v[N_v(D) - \lambda\Delta_v] > 0$ and makes the expected centered activation of $a^\star$ strictly negative in both Case~1 and Case~3.

Next, approximate $p,x,y,z,r_0$ on a grid with a large odd denominator $m$, keeping $\phi(r_0)>0$ at the rounded auxiliary rank and preserving the boundary equalities among the original ranks. If the rounded ranks leave a small residual $\epsilon$ in the centered activation of $\calH$, shift the two auxiliary weights by
\[
  \Bigl(-\frac{\epsilon}{P + R},\, \frac{\epsilon}{P + R}\Bigr),
\]
evaluated at the rounded values of $P$ and $R$. This restores the centered activation of $\calH$ to zero while preserving the total weight. When the rational approximation is close enough, all voter-type weights remain positive, the expected centered activation of $a^\star$ remains strictly negative, and $\E_v[N_v(D) - \lambda\Delta_v] > 0$ continues to hold. Realize $\calH$ and $\calM$ with $m p$ and $m(1 - p)$ candidates of equal lottery mass $1/m$, and average each voter type uniformly over all permutations within $\calH$ and within $\calM$. By symmetry within each class, every supported candidate has expected activation $\tau$, while $a^\star$ has expected activation strictly below $\tau$. Applying \cref{lem:equal-weights} then replaces the real voter weights with a finite election of equally weighted voters.
\end{proof}

\subsection{Equal-Weight Elections}
\label{app:equal-weights}

The constructions above use real voter-type weights to balance the centered activation of $\calH$ to zero. The following lemma shows that any strict violation on a uniform lottery of odd support size $m$ can be realized with finitely many equally weighted voters.

\begin{lemma}[Realizing real voter weights with equally weighted voters]
\label{lem:equal-weights}
Fix a nonconstant activation $g \in \calG$. Suppose a finite weighted election has a $g$-stable lottery $D_{\mathrm{unif}}$ that is uniform on an odd number $m$ of supported candidates, every unsupported candidate has expected activation strictly below $\tau$, and some candidate distances $\rho$ satisfy $\E_v[N_v(D_{\mathrm{unif}})] > \lambda\,\E_v[\Delta_v]$. If $m$ is large enough, then there is a finite election with equally weighted voters that has a $g$-stable lottery violating the same inequality \cref{eq:height-inequality}.
\end{lemma}

\begin{proof}
For each supported candidate $i \in \{1, \dots, m\}$, write $S_i(D) := \E_v[\psi_v(i; D)]$, and let $T := \{h \in \R^m : \sum_{i=1}^m h_i = 0\}$ be the tangent space to the simplex at $D_{\mathrm{unif}}$. Because $G(t) = \int_0^t g(r)\,\mathrm{d}r$ is $C^1$ and each supported candidate has positive mass $D_i > 0$ near $D_{\mathrm{unif}}$, the interval averages $[G(r + D_i) - G(r)]/D_i$ are $C^1$ functions of $D$ on a neighborhood of $D_{\mathrm{unif}}$. Differentiating the identity $\sum_{i=1}^m D_i S_i(D) = \tau$ at the stationary lottery $D_{\mathrm{unif}}$ (where $S_i(D_{\mathrm{unif}}) = \tau$ for all $i$) shows that the Jacobian $\mathrm{D}S(D_{\mathrm{unif}})$ maps $T$ into $T$.

\paragraph{An auxiliary profile with an invertible Jacobian.}
Start from the profile that assigns equal weight to all $m!$ rankings of the $m$ supported candidates and places all unsupported candidates at the bottom of every ranking. By symmetry among the $m$ supported candidates, $D_{\mathrm{unif}}$ is stationary and $\mathrm{D}S(D_{\mathrm{unif}})\big|_T = a I_T$ for some scalar $a \in \R$. Since $g \in \calG$ is continuous and nonconstant, for all large enough $m$ there is an index $r \in \{0, \dots, m - 3\}$ with
\[
  d \;:=\; \frac{m}{3}\Bigl[g\Bigl(\frac{r + 2}{m}\Bigr) - g\Bigl(\frac{r + 1}{m}\Bigr)\Bigr] \;\ne\; 0.
\]
For any ordered triple $(i, j, k)$ of distinct supported candidates, consider the rankings that place $\{i, j, k\}$ in consecutive positions $r + 1, r + 2, r + 3$ and permute the other $m - 3$ supported candidates uniformly. Add the uniform mixture of the three cyclic orders $(i, j, k), (j, k, i), (k, i, j)$ and subtract the uniform mixture of their reversals $(k, j, i), (i, k, j), (j, i, k)$. Under this signed perturbation, each of $i, j, k$ occupies each of the positions $r + 1, r + 2, r + 3$ with net weight zero, so every candidate's marginal position distribution is unchanged and $D_{\mathrm{unif}}$ remains stationary. Differentiating the interval averages at $D_{\mathrm{unif}}$ shows that on the coordinates $(i, j, k)$, this signed perturbation changes $\mathrm{D}S(D_{\mathrm{unif}})$ by
\[
  d\begin{pmatrix}
    0 & -1 & 1 \\
    1 & 0 & -1 \\
    -1 & 1 & 0
  \end{pmatrix},
\]
with all other entries equal to zero: in the $(i, j)$ entry, the terms where $j$ is third in the triple cancel between $(k, i, j)$ and $(i, k, j)$, leaving $-d$ from the middle grid point $(r + 2)/m$ minus $(r + 1)/m$, and the other entries follow by cyclic symmetry. Own-mass derivatives and derivatives with respect to candidates outside $\{i, j, k\}$ cancel because $\{i, j, k\}$ occupies consecutive positions and each candidate's marginal position distribution is unchanged.

Since $m$ is odd, we can partition $\{2, \dots, m\}$ into pairs and sum these signed perturbations over the triples $(1, 2, 3), (1, 4, 5), \dots, (1, m - 1, m)$. The resulting skew-symmetric matrix $B$ has kernel $\R\1$: for each pair $(2s, 2s + 1)$, the two equations $(B h)_{2s} = d(h_1 - h_{2s+1}) = 0$ and $(B h)_{2s+1} = d(h_{2s} - h_1) = 0$ imply $h_{2s} = h_{2s+1} = h_1$. Thus $B|_T$ is invertible, and since $B$ is skew-symmetric, $a I_T + t B|_T$ is invertible for every $t \ne 0$ (if $(a I_T + t B)h = 0$ for $h \in T$, taking the inner product with $h$ gives $a\|h\|^2 = 0$, so either $a = 0$ and $B h = 0$, or $h = 0$). Choosing $t > 0$ small enough keeps the weight of every ranking in the perturbed profile positive. This gives an auxiliary profile at which $D_{\mathrm{unif}}$ is stationary and $\mathrm{D}S(D_{\mathrm{unif}})\big|_T$ is invertible.

\paragraph{Rational voter weights and an exactly stable lottery.}
Mix a fraction $t \in (0, 1)$ of this auxiliary profile into the original weighted election. At every $t$, $D_{\mathrm{unif}}$ remains stationary. Because $\det\bigl(\mathrm{D}S_t(D_{\mathrm{unif}})\big|_T\bigr)$ is a polynomial in $t$ that is nonzero at $t = 1$, it is nonzero for arbitrarily small $t > 0$. By the strict margins in the original profile and continuity in $t$, we may fix such a $t>0$ small enough that every unsupported candidate still has expected activation strictly below $\tau$ and $\E_v[N_v(D_{\mathrm{unif}})] > \lambda\,\E_v[\Delta_v]$ still holds.

Applying the Implicit Function Theorem to the map $(S_1 - S_m, \dots, S_{m-1} - S_m)$ on $T$ at $(w_t, D_{\mathrm{unif}})$ shows that for rational voter weights $w'$ close enough to the mixed weights $w_t$, there is a nearby interior lottery $D'$ on the same $m$ supported candidates at which $S_1(D') = \dots = S_m(D')$. Since $\sum_{i=1}^m D'_i S_i(D') = \tau$, their common value is $\tau$. By continuity, every unsupported candidate still has expected activation strictly below $\tau$ at $D'$, and $\E_v[N_v(D')] > \lambda\,\E_v[\Delta_v]$. Clearing the common denominator of the rational weights $w'$ gives a finite election with equally weighted voters in which $D'$ is $g$-stable and violates \cref{eq:height-inequality}.
\end{proof}

In our constructions of Elections~2 and~3, the rational approximations to the class masses and split ranks can be chosen with an arbitrarily large odd common denominator $m$; uniform permutation within each class then produces the uniform lottery $D_{\mathrm{unif}}$ required by \cref{lem:equal-weights}.

\section{Proofs for the Variational Problem}
\label{app:variational}

All intervals below have the exact rational endpoints listed in \cref{app:rank-intervals}. Fix $\lambda \in [\lambda_0, \lambda_1]$, and write $\calK_\lambda$ for the set of pairs $(S, \phi) \in \R \times L^\infty([\ell, t_e])$ satisfying $S \le S_+(\lambda)$ and \cref{eq:low-constraints}, with the primitive $J(r)$ defined on $[p_-, z_0]$ by
\[
  J(r) \;=\;
  \begin{cases}
    C_{\lambda, S}(r), & p_- \le r \le \ell, \\
    C_{\lambda, S}(\ell) + \displaystyle\int_\ell^r \phi(u)\,\mathrm{d}u, & \ell \le r \le t_e, \\
    Q_\lambda(r) - S, & t_e \le r \le z_0.
  \end{cases}
\]
The derivative bound $\phi(r) \le \lambda - r$ and the endpoint condition $J(t_e) = Q_\lambda(t_e) - S$ imply $J(r) \ge Q_\lambda(r) - S$ on $[p_-, z_0]$. Combining this lower bound with $J(r) \le C_{\lambda, S}(r)$ at an interior point $r \in (\ell, t_e)$ bounds $S$ from below by $S_- > 0$ (\cref{app:rank-intervals}). Thus $S$ and $\phi$ range over bounded subsets of $\R$ and $L^\infty([\ell, t_e])$.

In the ODE~\eqref{eq:high-ode}, the essential infimum is taken with respect to three-dimensional Lebesgue measure on the control domain
\[
  \calU \;:=\; \bigl\{(p, x, y) : p \in [p_-, p_+],\ x \in [x_-, x_+],\ \max\{x_-, p\} \le y \le x_+\bigr\},
\]
so it does not depend on the representative of $\phi$. Because the denominator $\Delta > 0.0519$ is bounded away from zero on $\calK_\lambda \times \calU$ (\cref{app:rank-intervals,app:variational-estimates}) and $H(z, Y; p, x, y)$ is polynomial in $z$ and affine in $Y$ with uniformly bounded coefficients, the right-hand side of \cref{eq:high-ode} is Lipschitz in $(z, Y)$ uniformly over $(p, x, y) \in \calU$. Taking the essential infimum over $\calU$ and the minimum with $\lambda - z$ preserves this Lipschitz bound. Thus every pair $(S, \phi) \in \calK_\lambda$ has a unique continuation $Y(z)$ on $[z_0, z_1]$ with Lipschitz derivative $Y'(z)$.

In \cref{eq:high-binary-compatibility}, the tested pairs $(z, q)$ range over $[z_B, z_1] \times \calQ$, where $\calQ \subset [p_-, \ell]$. Define the maximum violation functional
\begin{equation}
  j_\lambda(S, \phi) \;:=\; \max_{(z, q) \in [z_B, z_1] \times \calQ}\bigl[T_{\lambda, S}(z, q) - Y(z)\bigr].
  \label{eq:j-lambda-def}
\end{equation}
Then $(S, \phi) \in \calK_\lambda$ satisfies the compatibility condition~\eqref{eq:high-binary-compatibility} if and only if $j_\lambda(S, \phi) \le 0$.

Notice that whenever $j_\lambda(S, \phi) \le 0$, the lower bound $Y(z) \ge Q_\lambda(z) - Q_\lambda(1)$ holds automatically across $[z_0, z_1]$. Indeed, $Y'(z) \le \lambda - z = Q'_\lambda(z)$, so $Y(z) - Q_\lambda(z) + Q_\lambda(1)$ is nonincreasing in $z$. At the right endpoint $z_1$ and $q_0 \in \calQ$, the numerical check in \cref{app:variational-estimates} shows $T_{\lambda, S}(z_1, q_0) > Q_\lambda(z_1) - Q_\lambda(1)$ for all $\lambda \in [\lambda_0, \lambda_1]$ and $S \in [S_-, S_+(\lambda)]$, so $Y(z_1) \ge T_{\lambda, S}(z_1, q_0)$ implies $Y(z) \ge Q_\lambda(z) - Q_\lambda(1)$ for all $z \in [z_0, z_1]$.

\subsection{Reducing to the Low-Rank Interval}
\label{app:reduction}

\begin{proof}[Proof of {\cref{prop:reduction}}]
Starting from a primitive $J$ and derivative $\phi = J'$ given by \cref{lem:election-1}, \cref{lem:election-2}, and \cref{lem:election-3}, we first remove the slack $\sigma(r) = \lambda - r - \phi(r)$ from the middle ranks so that $\phi(r) = \lambda - r$ there, and then raise $J$ to $C_{\lambda, S}$ on $[p_-, \ell]$.

For a triple $(p, x, y) \in \calU$ and high rank $z \in [z_0, z_1]$, write
\[
  R = J(p), \qquad P = J(1 - p), \qquad K = J(z - p) - J(z),
\]
\[
  h = 2\lambda - x, \qquad k = p + y - \lambda, \qquad U = z - \lambda, \qquad \Delta = h R - k P.
\]
The necessary inequalities from Elections~2 and~3 include $(1 - p)R \le \lambda P$, $p P \le \lambda R$, and
\[
  E_3 \;:=\; R\bigl(U\phi(x) + h\phi(z)\bigr) + P\bigl(U\phi(y) - k\phi(z)\bigr) + K\bigl(k\phi(x) + h\phi(y)\bigr) \;\le\; 0,
\]
which solves to $\phi(z) \le H(z, J(z); p, x, y)$ as in \cref{eq:H-def}. All auxiliary ranks and numerical constants used below are recorded in \cref{app:variational-estimates}.

\paragraph{Step 1: Making $\phi(r) = \lambda - r$ on the middle ranks.}
Let $\sigma(r) = \lambda - r - \phi(r) \ge 0$. \cref{lem:variational-estimates} bounds the total slack $\mu = \int_{x_+}^{z_0} \sigma(r)\,\mathrm{d}r < \mu_0$ on an interval containing $[x_+, z_0]$ and bounds $\int_0^{0.5} \sigma(r)\,\mathrm{d}r \le S_+(\lambda)$. We now reduce the gap $\sigma(r) = \lambda - r - \phi(r)$ to zero on $(x_+, z_0)$ by raising $\phi(r)$ to $\lambda - r$ there. Let $s := \max_{(z, q) \in [z_B, z_1] \times \calQ}(1 - z + q) = 0.44$, and split $\int_{x_+}^{z_0} \sigma(r)\,\mathrm{d}r = \mu_1 + \mu_2$ into mass $\mu_1$ on $(x_+, s)$ and mass $\mu_2$ on $[s, z_0)$, so that modifications on $[s, z_0)$ do not affect the two-class lower bound $J(1 - z + q) - J(q)$.
\begin{itemize}
  \item In the first phase, we reduce $\sigma$ to $0$ on $(x_+, s)$ while simultaneously reducing $\int_0^{p_-} \sigma(r)\,\mathrm{d}r$ by $\eta = \kappa_f \mu_1$ below $p_-$ (\cref{app:variational-estimates} certifies enough available gap $\sigma$ below $p_-$). Because the low positions $x, y \in [x_-, x_+]$ lie outside the modified intervals, $\phi(x)$ and $\phi(y)$ remain unchanged. Decreasing $\sigma$ by $\eta = \kappa_f\mu_1$ below $p_-$ and by $\mu_1$ on $(x_+, s)$ increases $R = J(p)$ by $\delta R = \kappa_f\mu_1$ and $P = J(1 - p)$ by $\delta P = (\kappa_f + 1)\mu_1$. Thus
  \[
    \lambda\,\delta R - p\,\delta P \;=\; \bigl[(\lambda - p)\kappa_f - p\bigr]\mu_1 \;>\; 0,
    \qquad
    \lambda\,\delta P - (1 - p)\,\delta R \;=\; \bigl[\lambda - (1 - p - \lambda)\kappa_f\bigr]\mu_1 \;>\; 0,
  \]
  for every $p \in [p_-, p_+]$ because
  \[
    \frac{p_+}{\lambda - p_+} \;<\; \kappa_f \;<\; \frac{\lambda}{1 - p_- - \lambda},
  \]
  so both $q = 0$ two-class inequalities in \cref{eq:two-class-zero} are preserved.
  \item In the second phase, reducing $\sigma$ to $0$ on $[s, z_0)$ increases $P$ without changing $R$, which strengthens $(1 - p)R \le \lambda P$. For the opposite inequality $p P \le \lambda R$, the universal slack bounds in \cref{lem:variational-estimates} initially give
  \[
    R \;\ge\; Q_{\lambda_0}(p) - 0.0267,
    \qquad
    P \;\le\; Q_{\lambda_1}(1 - p) - 0.02196.
  \]
  The first phase increases $\lambda R - p P$, and the second decreases it by at most $p\mu_2 \le p\mu_0$. By the Bernstein bound in \cref{lem:variational-estimates}, throughout both phases we have
  \[
    \lambda R - p P \;\ge\; \lambda_0\bigl[Q_{\lambda_0}(p) - 0.0267\bigr] - p\bigl[Q_{\lambda_1}(1 - p) - 0.02196 + \mu_0\bigr] \;>\; 0.0143 \;>\; 0.
  \]
\end{itemize}

Before comparing the high-rank continuations, we record the domain bounds along the linear interpolation of each phase. Let $Y$ be the maximal continuation at the start of the phase, with $Y(z_0) = J_{\mathrm{old}}(z_0)$. Initially $Y \ge J_{\mathrm{old}} \ge Q_\lambda - Q_\lambda(1)$ by ODE comparison and $J' \le \lambda - r$. For a phase that removes slack $\mu \in \{\mu_1, \mu_2\}$ on the middle ranks and $\eta \in \{\kappa_f\mu_1, 0\}$ below $p_-$, let $e(z)$ solve the linear ODE
\[
  e'(z) \;=\; \overline{\alpha}\,e(z) + \overline{p}\,\mu, \qquad e(z_0) \;=\; 0,
\]
with $\overline{\alpha} := 4.53$ and $\overline{p} := 6.19$, and define the candidate subsolution
\[
  \widetilde{Y}(z) \;:=\; Y(z) + \eta + \mu - e(z).
\]
Setting
\[
  \omega \;:=\; \frac{\overline{p}}{\overline{\alpha}}\bigl(\mathrm{e}^{\overline{\alpha}(z_1 - z_0)} - 1\bigr) \;<\; 0.784 \;<\; 1,
\]
we have $0 \le e(z) \le \omega\mu < \mu$ on $[z_0, z_1]$. Thus $\widetilde{Y}(z) \ge Y(z) + \eta + (1 - \omega)\mu \ge Y(z) \ge Q_\lambda(z) - Q_\lambda(1)$, so every interpolated high-rank value between $Y(z)$ and $\widetilde{Y}(z)$ stays above $Q_\lambda(z) - Q_\lambda(1)$. Because the low-rank modifications preserve $J \le Q_\lambda$, along the interpolation we have
\[
  K \;\le\; p U - \frac{p^2}{2} + Q_\lambda(1) \;\le\; \frac{U^2}{2} + Q_\lambda(1) \;<\; 0.111315.
\]
Meanwhile, $J'_{\mathrm{old}} \le \lambda - r$ gives $K_{\mathrm{old}} \ge p(U - p/2)$, and the phase increments satisfy $K \ge K_{\mathrm{old}} - \mu$, so \cref{lem:variational-estimates} gives
\[
  K \;\ge\; p\Bigl(U - \frac{p}{2}\Bigr) - \mu \;>\; 0,
  \qquad
  c_x \;:=\; U R + k K \;\ge\; (0.775 - \lambda_1)\cdot 0.0826 - 0.02733 \cdot 0.111315 \;>\; 0.0145,
\]
while $c_y := U P + h K > 0$ since $U, P, h, K > 0$. Together with $\Delta > 0.0519$, these bounds hold across the entire interpolation before comparing $\widetilde{Y}$ with the new continuation; once that comparison is proved at the end of Phase~1, it preserves $Y \ge Q_\lambda - Q_\lambda(1)$ at the start of Phase~2.

To track how the three-class differential bound $\widehat{H} := \min\{\lambda - z, H\}$ changes during each phase, write $M_{xy} := k\phi(x) + h\phi(y)$. If $M_{xy} \le 0$, evaluating $E_3$ at $\phi(z) = \lambda - z = -U$ gives
\[
  E_3\big|_{\phi(z) = \lambda - z} \;\le\; U(-\lambda R + p P) + K M_{xy} \;\le\; 0,
\]
so $H \ge \lambda - z$ and this triple imposes no restriction beyond $\phi(z) \le \lambda - z$. Whenever $H < \lambda - z$, we have $\alpha := M_{xy}/\Delta > 0$. By definition of $H$, we have $c_x\phi(x) + c_y\phi(y) + \Delta\,H = 0$; combining this with $c_x h - c_y k = U(h R - k P) = U\Delta$ gives $U\phi(x) + h H = -\alpha c_y$ and $U\phi(y) - k H = \alpha c_x$. Differentiating $c_x\phi(x) + c_y\phi(y) + \Delta\,H = 0$ with respect to $R, P, K$ then gives
\[
  H_R \;=\; \frac{\alpha c_y}{\Delta} \;>\; 0,
  \quad
  H_P \;=\; -\frac{\alpha c_x}{\Delta} \;<\; 0,
  \quad
  H_K \;=\; -\alpha,
  \quad
  H_R + H_P \;=\; \frac{\alpha\bigl[(h - k)K + U(P - R)\bigr]}{\Delta} \;\ge\; 0,
\]
since $h - k = 3\lambda - x - p - y > 0$ and $P \ge \frac{1 - p}{\lambda}R > R$ for $p \le p_+$. We bound $\alpha$ and $|H_P|$ according to the sign of $k$:
\begin{itemize}
  \item When $k < 0$, the condition $H < \lambda - z = -U$ implies $\phi(x) > (U\Delta - c_y\phi(y))/c_x$. Multiplying by $k < 0$ and using $c_x h - c_y k = U\Delta$ gives
  \[
    M_{xy} \;<\; \frac{U\Delta\bigl[k + \phi(y)\bigr]}{c_x} \;\le\; \frac{U\Delta p}{c_x},
  \]
  where the second inequality uses $\phi(y) \le \lambda - y$ and $k + \lambda - y = p$. Here $p < \lambda - y \le 0.27733$ and $-k \le 0.02733$, so the bounds in \cref{lem:variational-estimates} give
  \[
    \alpha \;<\; \frac{0.27733}{0.0826 - 0.02733 \cdot 0.159/(0.775 - \lambda_1)} \;<\; 4.47,
    \qquad
    |H_P| \;=\; \frac{\alpha c_x}{\Delta} \;<\; \frac{0.313 \cdot 0.27733}{0.0519} \;<\; 1.68.
  \]
  \item When $k \ge 0$, bounding $\phi(x) \le \lambda - x$ and $\phi(y) \le \lambda - y$ directly in $M_{xy} = k\phi(x) + h\phi(y)$ and using $p \le y$ gives
  \[
    M_{xy} \;\le\; k(\lambda - x) + h(\lambda - y) \;=\; p(\lambda - x) + \lambda(\lambda - y) \;\le\; \lambda^2 - x y \;\le\; M_+ \;:=\; \lambda_1^2 - x_-^2,
  \]
  and hence
  \[
    \alpha \;<\; \frac{M_+}{0.0519} \;<\; 4.53,
    \qquad
    |H_P| \;<\; \frac{M_+(0.313 \cdot 0.136 + 0.178 \cdot 0.159)}{0.0519^2} \;<\; 6.19.
  \]
\end{itemize}
In both cases, $\alpha < \overline{\alpha} = 4.53$ and $|H_P| < \overline{p} = 6.19$.

Along the linear interpolation of the data, decreasing $\sigma$ raises $J$ by $\eta$ on $[p_-, x_+]$ and by $\eta + \mu$ at $z_0$, so the increments of $R$, $P$, and $K = J(z - p) - \widetilde{Y}(z)$ satisfy $\delta R = \eta$, $\delta P = \eta + \delta_P$, and $\delta K = \delta_K - \mu + e(z)$ with $0 \le \delta_P, \delta_K \le \mu$. Using $H_R + H_P \ge 0$, $H_P \in [-\overline{p}, 0)$, and $H_K = -\alpha \in [-\overline{\alpha}, 0)$, the resulting change in $H$ along the interpolation is bounded below by
\[
  \delta H \;=\; (H_R + H_P)\eta + H_P\,\delta_P - \alpha\bigl(\delta_K - \mu + e(z)\bigr)
  \;\ge\; -\overline{p}\,\mu - \overline{\alpha}\,e(z) \;=\; -e'(z).
\]
Thus $\widehat{H}_{\mathrm{new}}\bigl(z, \widetilde{Y}(z); p, x, y\bigr) \ge \widehat{H}_{\mathrm{old}}\bigl(z, Y(z); p, x, y\bigr) - e'(z) \ge Y'(z) - e'(z) = \widetilde{Y}'(z)$ (including at points where the active branch of $\min\{\lambda - z, H\}$ switches), so the new maximal continuation lies at or above $\widetilde{Y}(z) \ge Y(z) + \eta + (1 - \omega)\mu$.

In the high-rank two-class bound at $(z, q) \in [z_B, z_1] \times \calQ$, both $1 - z + q$ and $q$ lie below $s$. Thus the second phase (above $s$) leaves $J(1 - z + q) - J(q)$ unchanged, while the first phase (below $s$) increases $J(q)$ by $\eta$ and $J(1 - z + q)$ by at most $\eta + \mu_1$, increasing their difference by at most $\mu_1$. With $\Gamma := \sup_{(z, q) \in [z_B, z_1] \times \calQ} \frac{z}{\lambda - q} \le 3.100$, the margin $Y(z) - \frac{z}{\lambda - q}\bigl[J(1 - z + q) - J(q)\bigr]$ increases by at least
\[
  (\kappa_f + 1 - \omega - \Gamma)\mu_1 + (1 - \omega)\mu_2 \;\ge\; 0.096\,\mu_1 + 0.216\,\mu_2 \;\ge\; 0,
\]
and strictly increases whenever $\mu_1 + \mu_2 > 0$. After both phases, $J(r) = Q_\lambda(r) - S$ holds identically on $[x_+, z_0]$ (with the new cumulative slack $S \le S_+(\lambda)$), and $J(r) \le C_{\lambda, S}(r)$ holds through $t_e$ because $C_{\lambda, S}(r) - (Q_\lambda(r) - S)$ is positive at $x_+$ and increasing on $[x_+, t_e]$.

\paragraph{Step 2: Raising $J$ to $C_{\lambda, S}$ on $[p_-, \ell]$ and taking the nondecreasing envelope.}
Define the modified low-rank primitive
\[
  \widetilde{J}(r) \;:=\;
  \begin{cases}
    C_{\lambda, S}(r), & p_- \le r \le \ell, \\
    \max\Bigl\{C_{\lambda, S}(\ell),\, \max_{u \in [\ell, r]} J(u)\Bigr\}, & \ell \le r \le t_e.
  \end{cases}
\]
Because the relocated $J$ is Lipschitz with $J'(r)=\phi(r)\le\lambda-r$ almost everywhere and $J \le C_{\lambda, S}$ on $[p_-, t_e]$, the function $\widetilde{J}$ is Lipschitz and nondecreasing, lies between $J$ and $C_{\lambda, S}$, and satisfies $0 \le \widetilde{\phi}(r) := \widetilde{J}'(r) \le \lambda - r$ almost everywhere on $[\ell, t_e]$. To evaluate the envelope at $t_e$, first note that, on $[\ell,0.345]$, we have $J(u)\le C_{\lambda,S}(0.36)<Q_\lambda(x_+)-S$ by monotonicity of the cap and \cref{app:variational-estimates}. For $0.345\le u\le x_+$, integrating backwards from $J(x_+)=Q_\lambda(x_+)-S$ gives
\[
  J(u)=Q_\lambda(u)-S+\int_u^{x_+}\sigma(r)\,\mathrm{d}r
       \le Q_\lambda(x_+)-S+\mu_0 < Q_\lambda(t_e)-S.
\]
Indeed, the relocation leaves the slack on $[0.345,x_+]$ unchanged, and $Q_\lambda(0.385)-Q_\lambda(0.38)\ge 0.000899125>\mu_0$. On $[x_+,t_e]$ the primitive equals the increasing function $Q_\lambda-S$. These bounds also dominate the floor $C_{\lambda,S}(\ell)$, so $\widetilde J(t_e)=Q_\lambda(t_e)-S$. At almost every $r \in [\ell, t_e]$, either $\widetilde{\phi}(r) = 0$ or $\widetilde{\phi}(r) = \phi(r) > 0$. By \cref{lem:variational-estimates}, any triple with $\widetilde{H} \le \lambda - z$ has $\widetilde{\phi}(y) > 0.0332 > 0$ and $0 < \widetilde{\theta} < 1$, and comparing $x$ with $[0.34, 0.35]$ using \cref{eq:control-change-identities} shows that any near-minimizing first position has $\widetilde{\phi}(x) > 0.0618 - 8\varepsilon > 0$. Thus the essential infimum in \cref{eq:high-ode} is determined by triples where $\widetilde{\phi}(x) = \phi(x) > 0$ and $\widetilde{\phi}(y) = \phi(y) > 0$. Since $\widetilde{R} = \widetilde{J}(p) \ge R$ and $\partial E_3/\partial R = U\phi(x) + h\phi(z) \le -\lambda U < 0$, the maximal continuation $Y(z)$ can only increase, while raising $J(q)$ to $C_{\lambda, S}(q)$ for $q \in \calQ \subset [p_-, \ell]$ decreases $T_{\lambda, S}(z, q)$. Hence $(S, \widetilde{\phi})$ satisfies \cref{eq:low-constraints} and its continuation satisfies \cref{eq:high-binary-compatibility}.
\end{proof}

\subsection{Taking a Limit of Feasible Functions}
\label{app:attainment}

\begin{proof}[Proof of {\cref{prop:attainment}}]
Take a sequence of feasible parameters $\lambda_n \downarrow \lambda^\star$ with corresponding feasible pairs $(S_n, \phi_n) \in \calK_{\lambda_n}$ satisfying $j_{\lambda_n}(S_n, \phi_n) \le 0$.

\paragraph{Step 1: Passing to a weak-* limit.}
Because $S_n \in [S_-, S_+(\lambda_n)]$ and $0 \le \phi_n(r) \le \lambda_n - r$ are uniformly bounded, we may pass to a subsequence along which $S_n \to S \in [S_-, S_+(\lambda^\star)]$ and $\phi_n \stackrel{*}{\rightharpoonup} \phi$ in $L^\infty([\ell, t_e])$. Their primitives $J_n$ converge uniformly to $J$ on $[p_-, z_0]$, so $(S, \phi) \in \calK_{\lambda^\star}$.

To show that $j_\lambda(S, \phi)$ is jointly lower semicontinuous with respect to $(\lambda, S)$ and the weak-* topology on $\phi$, represent the branch $\lambda - z$ in \cref{eq:high-ode} as an extra control $0$ with right-hand side $\lambda - z$. Define an \emph{averaged control} $\nu = (\nu_z)_{z \in [z_0, z_1]}$ to be a measurable family of probability measures on $\{0\} \cup \calU$ whose restriction to $\calU$ at each $z$ has absolutely continuous $x$- and $y$-marginals on $[x_-, x_+]$. The Lebesgue-point approximation in \cref{sec:reference-comparison} extends the essential-infimum bound $Y'(z) \le H(z, Y(z); p, x, y)$ to all triples $(p, x, y) \in \calU$ whose low ranks $x$ and $y$ are Lebesgue points of $\phi$ (including boundary triples with $y = p$). Because the $x$- and $y$-marginals of $\nu_z$ are absolutely continuous, each averaged control assigns full measure to such triples without needing a joint density on $\calU$.

For each fixed averaged control $\nu$, the corresponding trajectory $Y_\nu(z)$ solves a linear ODE on $[z_0, z_1]$ with $Y_\nu(z_0) = Q_\lambda(z_0) - S$. As $(\lambda_n, S_n) \to (\lambda^\star, S)$ and $\phi_n \stackrel{*}{\rightharpoonup} \phi$, the primitives $J_n$ converge uniformly, and at each fixed $z \in [z_0, z_1]$ the $\nu_z$-averaged coefficients converge because the evaluations of $\phi$ are integrals against fixed $L^1$ marginal densities in $x$ and $y$. Since these coefficients are uniformly bounded, dominated convergence gives $L^1([z_0, z_1])$ convergence of the ODE coefficients, and Gr\"onwall's inequality gives uniform convergence $Y_{\nu, n} \to Y_\nu$ on $[z_0, z_1]$.

For any $(S, \phi) \in \calK_\lambda$ and any averaged control $\nu$, the $\nu_z$-averaged right-hand side is at least the essential infimum in \cref{eq:high-ode}, so ODE comparison gives $Y(z) \le Y_\nu(z)$. Conversely, for any $\varepsilon > 0$, at each $z \in [z_0, z_1]$ we may select control $0$ if $\lambda - z$ is within $\varepsilon$ of $Y'(z)$, and otherwise take the normalized Lebesgue measure on the positive-measure subset of triples in $\calU$ where $H(z, Y(z); p, x, y) \le Y'(z) + \varepsilon$. By Gr\"onwall's inequality, the resulting averaged control satisfies $0 \le Y_\nu(z) - Y(z) \le C\varepsilon$ with $C$ independent of $\varepsilon$. Therefore,
\[
  Y(z) \;=\; \inf_\nu Y_\nu(z)
  \qquad\text{and}\qquad
  j_\lambda(S, \phi) \;=\; \sup_{\nu,\, (z, q) \in [z_B, z_1] \times \calQ}\bigl[T_{\lambda, S}(z, q) - Y_\nu(z)\bigr].
\]
As a supremum of functionals that are continuous under $\lambda_n \to \lambda^\star$, $S_n \to S$, and $\phi_n \stackrel{*}{\rightharpoonup} \phi$, the functional $j_\lambda(S, \phi)$ is lower semicontinuous, so $j_{\lambda^\star}(S, \phi) \le \liminf_{n \to \infty} j_{\lambda_n}(S_n, \phi_n) \le 0$.

\paragraph{Step 2: Excluding $\lambda \le \underline{\lambda}$ and proving $j_{\lambda^\star}(S, \phi) = 0$.}
In \cref{sec:reference-comparison,app:calibration}, we exclude feasibility at $\underline{\lambda} = 0.56232844$ and bound the distance of any feasible function at $\lambda_{\mathrm{ref}} = 0.5623285$ from an explicit reference function. To transfer those bounds from a parameter $\lambda$ up to $\lambda + \delta$ ($\delta > 0$), note that increasing $\lambda$ weakens the three-class inequality $E_3 \le 0$ because
\[
  \frac{\partial E_3}{\partial \lambda}
  \;=\; -(R + K)\phi(x) + (2K - P)\phi(y) + (2R + P)\phi(z) \;<\; 0.
\]
Removing the additional constant slack $\delta$ on $(x_+, z_0)$ by the same relocation as in \cref{app:reduction}---with $\eta_\delta = \kappa_f(s - x_+)$ and $\zeta_\delta = x_+ - \eta_\delta$---shifts $J$ and $S$ to
\[
  \widetilde{J}(r) \;=\; J(r) + \eta_\delta\delta + \delta(r - x_+)_+,
  \qquad
  \widetilde{S} \;=\; S + \zeta_\delta\delta.
\]
As verified in \cref{app:variational-estimates}, the gap $C_{\lambda+\delta, \widetilde{S}}(r) - \widetilde{J}(r)$ increases on $[p_-, t_e]$, $\widetilde{S}$ remains inside the calibrated interval, and the high-rank subsolution is $\widetilde Y(z)=Y(z)+(z_0-\zeta_\delta)\delta-e(z)$ with $e(z)\ge0$. Consequently
\[
 Q_{\lambda+\delta}(z)-\widetilde S-\widetilde Y(z)
 =Q_\lambda(z)-S-Y(z)+(z-z_0)\delta+e(z)
 \ge Q_\lambda(z)-S-Y(z).
\] Because the reference comparison in \cref{sec:reference-comparison} requires only $J \le C_{\lambda, S}$ on $[p_-, \ell]$ (and not equality $J = C_{\lambda, S}$), any feasible pair at $\lambda$ produces a valid competitor at $\lambda + \delta$. Together with the strictly feasible trial at $\overline{\lambda} = 0.56232847$ from \cref{lem:numerical-calibration}, this proves $\underline{\lambda} < \lambda^\star \le \overline{\lambda}$.

Finally, suppose toward a contradiction that $j_{\lambda^\star}(S, \phi) < 0$. Choose $\overline{S} = 0.026 \in (S_-, S_+(\lambda^\star))$, set $h = t_e - \ell$ and $M = \overline{S} - Q_{\lambda^\star}(\ell) + C_{\lambda^\star, \overline{S}}(\ell)$, and consider the quadratic function
\[
  \overline{\phi}(r) \;:=\; \lambda^\star - r - \epsilon - 3\Bigl(\frac{M}{h} - \epsilon\Bigr)\Bigl(1 - \frac{r - \ell}{h}\Bigr)^2
  \qquad (\epsilon = 0.001).
\]
Integrating gives $\overline{J}(t_e) = Q_{\lambda^\star}(t_e) - \overline{S}$, and exact rational evaluation (\cref{app:variational-estimates}) shows that $\overline{\phi} > 0$, $\lambda^\star - r - \overline{\phi} \ge \epsilon$, and $\bigl[C_{\lambda^\star, \overline{S}}(r) - \overline{J}(r)\bigr]/(r - \ell)$ are bounded away from zero on $(\ell, t_e]$. Mixing a small fraction of $(\overline{S}, \overline{\phi})$ into $(S, \phi)$ preserves $j_{\lambda^\star} < 0$ by norm continuity while creating uniform strict slack in all low-rank constraints \cref{eq:low-constraints}. We can then decrease $\lambda$ slightly below $\lambda^\star$ and add a tiny constant to $\phi$ to maintain $J(t_e) = Q_\lambda(t_e) - S$, giving a feasible pair at $\lambda < \lambda^\star$, a contradiction. Thus $j_{\lambda^\star}(S, \phi) = 0$ (and applying the same perturbation to the strictly feasible trial at $\overline{\lambda}$ gives $\lambda^\star < \overline{\lambda}$).
\end{proof}

\subsection{How Active Inequalities Determine \texorpdfstring{$\phi$}{phi}}
\label{app:shape}

Because the strict-trial perturbation in the proof of \cref{prop:attainment} rules out $j_{\lambda^\star} < 0$ anywhere on $\calK_{\lambda^\star}$, every feasible pair $(S, \phi) \in \calK_{\lambda^\star}$ minimizes $j_{\lambda^\star}$ on the convex set $\calK_{\lambda^\star}$, with minimum $j_{\lambda^\star}(S, \phi) = 0$.

\begin{lemma}[Active-rank condition for a minimizing pair]
\label{lem:first-order}
Let $(S, \phi) \in \calK_{\lambda^\star}$ be a minimizing pair, write $I := [x_-, x_+]$, and let $c := \max\{r \in [\ell, t_e] : J(r) = C_{\lambda^\star, S}(r)\}$. For almost every $r \in (c, t_e]$, either $\phi(r) = \lambda^\star - r$, or $r \in I$ and there is a sequence of admissible triples in $\calU$ with $x = r$ or $y = r$ along which $H(z, Y(z); p, x, y) - Y'(z) \to 0$. In particular, $\phi(r) \ge \varepsilon_0 := 0.02$ almost everywhere on $(c, t_e]$, and $\phi(r) = \lambda^\star - r$ almost everywhere on $(c, t_e] \setminus I$.
\end{lemma}

\begin{proof}
Write $\lambda = \lambda^\star$, fix a Borel representative of $\phi$, and let $L \subset (x_-, x_+)$ be its full-measure Borel set of Lebesgue points. For $x, y \in L$, define
\[
  \Gamma(z, p, x, y) \;:=\; H\bigl(z, Y(z); p, x, y\bigr) - Y'(z) \;\ge\; 0,
\]
where nonnegativity holds for all $z \in [z_0, z_1]$ and $p \in [p_-, \min\{p_+, y\}]$ (including $y = p$, since $y \ge x_- > p_-$ allows interior approximations with mass strictly below the second rank) by approximating $(p, x, y)$ with interior triples at the Lebesgue points $x, y$ as in \cref{sec:reference-comparison}. Every averaged control assigns zero measure to non-Lebesgue points and endpoint ranks in $x$ and $y$ because its $x$- and $y$-marginals are absolutely continuous.

For $r \in L$, define the slice infima
\[
  m_x(r) \;:=\; \inf_{\substack{z \in [z_0, z_1],\, y \in L \\ p_- \le p \le \min\{p_+, y\}}} \Gamma(z, p, r, y),
  \qquad
  m_y(r) \;:=\; \inf_{\substack{z \in [z_0, z_1],\, x \in L \\ p_- \le p \le \min\{p_+, r\}}} \Gamma(z, p, x, r),
\]
and set $m(r) := \min\{m_x(r), m_y(r)\}$ for $r \in L$ and $m(r) := +\infty$ for $r \notin I$ (its values on the null set $I \setminus L$ do not matter). To see that $m_x$ and $m_y$ are Borel measurable on $L$, choose a countable set $\{d_j\} \subset L$ such that $\{(d_j, \phi(d_j))\}$ is dense in the graph $\{(r, \phi(r)) : r \in L\}$, a countable dense subset of $z \in [z_0, z_1]$, and rational $\vartheta \in [0, 1]$ parameterizing $p = p_- + \vartheta(\min\{p_+, d_j\} - p_-)$ in $m_x(r)$ and $p = p_- + \vartheta(\min\{p_+, r\} - p_-)$ in $m_y(r)$. By continuity of $J(p)$ and density of $\{(d_j, \phi(d_j))\}$, the infima over these countable sets equal $m_x(r)$ and $m_y(r)$ (including along the boundary $y = p$), and each term is a Borel function of $r$.

\paragraph{Sensitivity of near-active terminal tests.}
For an averaged control $\nu$ and terminal test $(t, q) \in [z_B, z_1] \times \calQ$, write $F_{\nu, t, q}(S, \phi) := T_{\lambda, S}(t, q) - Y_\nu(t)$. Whenever $F_{\nu, t, q}(S, \phi) \ge -\varepsilon$, we have $0 \le Y_\nu(t) - Y(t) \le \varepsilon$ since $Y(t) \ge T_{\lambda, S}(t, q)$. Define the averaged state coefficient $A_\nu(z)$ and adjoint $\rho_\nu(z)$ by
\[
  A_\nu(z) \;:=\; \int_{\calU} \alpha(p, x, y)\,\mathrm{d}\nu_z(p, x, y),
  \qquad
  \rho_\nu(z) \;:=\; \exp\!\Bigl(\int_z^t A_\nu(s)\,\mathrm{d}s\Bigr),
\]
where $\alpha(p, x, y) = \bigl[(p + y - \lambda)\phi(x) + (2\lambda - x)\phi(y)\bigr]/\Delta$ and control $0$ (the branch $\lambda - z$) contributes zero to $A_\nu(z)$. Then $w(z) := Y_\nu(z) - Y(z) \ge 0$ satisfies $w(z_0) = 0$ and
\[
  w'(z) \;=\; A_\nu(z)\,w(z) + d_\nu(z),
  \qquad
  d_\nu(z) \;:=\; \nu_z(\{0\})\bigl[\lambda - z - Y'(z)\bigr] + \int_{\calU} \Gamma(z, p, x, y)\,\mathrm{d}\nu_z(p, x, y) \;\ge\; 0.
\]
Since $A_\nu$ is uniformly bounded and $w(t) \le \varepsilon$, variation of constants gives, uniformly over all $(t, q)$ and $\nu$ with $F_{\nu, t, q}(S, \phi) \ge -\varepsilon$,
\begin{equation}
  \|w\|_{L^\infty([z_0, t])} + \int_{z_0}^t d_\nu(z)\,\mathrm{d}z \;\le\; C\varepsilon.
  \label{eq:needle-gap-budget}
\end{equation}
In particular, the $\nu_z\,\mathrm{d}z$-measure of triples in $\calU$ with $\Gamma(z, p, x, y) > 0.005$ is $O(\varepsilon)$. On all triples with $\Gamma(z, p, x, y) \le 0.005$, the near-minimizer bounds in \cref{lem:variational-estimates} give $\phi(x) > 0.0218 > \varepsilon_0$, $\phi(y) > 0.0289 > \varepsilon_0$, and $c_x/c_y < 0.74 < v_p := \frac{\lambda}{1 - \min\{p, \ell\}}$. For $k = p + y - \lambda \ge 0$, these lower bounds on $\phi(x)$ and $\phi(y)$ immediately give $\alpha > 0$; for $k < 0$, the bounds in \cref{lem:variational-estimates} give
\[
  k\phi(x) + h\phi(y) \;>\; 0.74465 \cdot 0.0289 - 0.02733 \cdot 0.27733 \;=\; 0.0139409561 \;>\; 0.0139,
\]
so $\alpha > 0$ and $H_R = \alpha c_y/\Delta > 0$ in both cases, and hence
\[
  H_R \;\ge\; 0,
  \qquad
  f \;:=\; v_p H_R + H_P \;=\; H_R\Bigl(v_p - \frac{c_x}{c_y}\Bigr) \;\ge\; 0,
\]
while control $0$ has zero derivatives with respect to $S$, $\phi$, and $R$.

Now differentiate $F_{\nu, t, q}$ at fixed $(\nu, t, q)$. Writing $y_S(z) := \partial_S Y_\nu(z) + 1$, we have $y_S(z_0) = 0$ and
\[
  y'_S(z) \;=\; A_\nu(z)\,y_S(z) - \int_{\calU} f_\nu(z, p, x, y)\,\mathrm{d}\nu_z(p, x, y),
\]
where the subscript $\nu$ indicates evaluation at $Y_\nu(z)$. Replacing $Y_\nu$ by $Y$ changes the integrand by $O(\varepsilon)$ by \cref{eq:needle-gap-budget}, and the measure of triples with $\Gamma > 0.005$ is also $O(\varepsilon)$. Since $\rho_\nu > 0$ is bounded and $|\partial_S T_{\lambda, S}(t, q)| \le b_S < 1$ (\cref{lem:variational-estimates}), we get
\begin{equation}
  \frac{\partial F_{\nu, t, q}}{\partial S}(S, \phi) \;=\; 1 + \frac{\partial T_{\lambda, S}}{\partial S}(t, q) - y_S(t) \;\ge\; d - C\varepsilon,
  \qquad
  d \;:=\; 1 - b_S \;>\; 0.
  \label{eq:needle-scalar-sign}
\end{equation}
Similarly, for any perturbation $h \in L^\infty([\ell, t_e])$ of $\phi$ at fixed $S$, write $B_h(p) := \int_\ell^p h(u)\,\mathrm{d}u$ for $p \ge \ell$ and $B_h(p) := 0$ for $p < \ell$. Variation of constants gives
\begin{equation}
  \mathrm{D}_\phi F_{\nu, t, q}(S, \phi)[h] \;=\; \int_{z_0}^t \int_{\calU} \rho_\nu(z)\Bigl[\frac{c_{x, \nu}}{\Delta}h(x) + \frac{c_{y, \nu}}{\Delta}h(y) - H_{R, \nu}\,B_h(p)\Bigr]\,\mathrm{d}\nu_z(p, x, y)\,\mathrm{d}z.
  \label{eq:needle-score-variation}
\end{equation}
Because $\Delta$ stays uniformly positive on a norm neighborhood of $\calK_\lambda$ and the ODE coefficients and $T_{\lambda, S}$ have uniformly bounded first and second Fr\'echet derivatives with respect to $v = (S, \phi)$ independently of $(\nu, t, q)$, differentiating the linear ODE for $Y_\nu$ once and twice gives a uniform Taylor expansion
\[
  F_{\nu, t, q}(v + \tau h_v) \;=\; F_{\nu, t, q}(v) + \tau\,\mathrm{D}F_{\nu, t, q}(v)[h_v] + O\bigl(\tau^2\|h_v\|^2\bigr)
\]
with first-derivative bound and quadratic remainder uniform over all $(\nu, t, q)$.

\paragraph{Constructing a feasible local perturbation.}
Suppose toward a contradiction that the lemma fails on a subset of $(c, t_e]$ of positive measure. Then there exist a Borel set $E \subset [c + d_c, t_e]$ of positive measure $b := |E| > 0$ (disjoint from $I \setminus L$) and constants $d_c, \eta, \delta > 0$ such that
\[
  \lambda - r - \phi(r) \;\ge\; \eta
  \qquad\text{and}\qquad
  m(r) \;\ge\; \delta
  \qquad\text{for all } r \in E.
\]
Set
\[
  a_E \;:=\; 1 - \frac{\lambda}{1 - \ell} \;>\; 0,
  \qquad
  h \;:=\; \1_E,
  \qquad
  s \;:=\; -\frac{b}{a_E} \;<\; 0,
  \qquad
  (S_\tau, \phi_\tau) \;:=\; (S + \tau s,\, \phi + \tau h).
\]
The endpoint condition $J(t_e) = Q_\lambda(t_e) - S$ is affine in $(S, \phi)$ and has variation $a_E s + b = 0$, so it holds for every $\tau$. Decreasing $S$ preserves $S_\tau \le S_+(\lambda)$ (and the lower bound $S_\tau \ge S_-$ follows automatically from the other constraints in \cref{eq:low-constraints}), while $0 \le \phi_\tau(r) \le \lambda - r$ holds for all $0 < \tau \le \eta$. Furthermore, the change in $J - C_{\lambda, S}$ is
\[
  \delta\bigl(J - C_{\lambda, S}\bigr)(r) \;=\; \int_\ell^r h(u)\,\mathrm{d}u + s\,\frac{\lambda(r - \ell)}{(1 - r)(1 - \ell)}.
\]
For $r \in [\ell, c + d_c]$, we have $\int_\ell^r h(u)\,\mathrm{d}u = 0$ and $s < 0$, so $\delta(J - C_{\lambda, S})(r) \le 0$. On $[c + d_c, t_e]$, the continuous function $C_{\lambda, S} - J$ has a strictly positive minimum because $c$ is the last contact of $J$ with $C_{\lambda, S}$. Hence $(S_\tau, \phi_\tau) \in \calK_\lambda$ for all small enough $\tau > 0$.

For any $(\nu, t, q)$ with $F_{\nu, t, q}(S, \phi) \ge -\varepsilon$, every triple in $\calU$ with $x \in E$ or $y \in E$ has $\Gamma(z, p, x, y) \ge m(r) \ge \delta$, so \cref{eq:needle-gap-budget} bounds the $\nu_z\,\mathrm{d}z$-measure of such triples by $C\varepsilon/\delta$. Since $0 \le B_h(p) \le b$ and $H_{R, \nu} \ge 0$ outside a set of measure $O(\varepsilon)$, \cref{eq:needle-scalar-sign,eq:needle-score-variation} give
\[
  \mathrm{D}_\phi F_{\nu, t, q}(S, \phi)[h] \;\le\; \frac{C\varepsilon}{\delta} + C\varepsilon b
  \qquad\text{and}\qquad
  \mathrm{D}F_{\nu, t, q}(S, \phi)[(s, h)] \;\le\; -\frac{d b}{a_E} + C\varepsilon\Bigl(\frac{b}{a_E} + b + \frac{1}{\delta}\Bigr).
\]
Fix $\varepsilon > 0$ small enough that the right-hand side is at most $-d b/(2a_E)$ for every $(\nu, t, q)$ with $F_{\nu, t, q}(S, \phi) \ge -\varepsilon$. By the uniform quadratic remainder, a small enough step $\tau > 0$ makes $F_{\nu, t, q}(S_\tau, \phi_\tau) < 0$ simultaneously for all such $(\nu, t, q)$, while all $(\nu, t, q)$ with $F_{\nu, t, q}(S, \phi) < -\varepsilon$ remain strictly negative by the uniform first-derivative bound. Taking the supremum over $(\nu, t, q)$ gives $j_\lambda(S_\tau, \phi_\tau) < 0$, a contradiction. Thus $m(r) = 0$ at almost every $r \in (c, t_e]$ where $\phi(r) < \lambda - r$. Since $m(r) = +\infty$ outside $I$ and any triple with $\Gamma \le 0.005$ satisfies $\phi(x), \phi(y) > \varepsilon_0$, this proves all claims of the lemma.
\end{proof}

\begin{proof}[Proof of {\cref{prop:shape}}]
Throughout the proof, write $\lambda = \lambda^\star$, $C = C_{\lambda^\star, S}$, and for $(p, x, y) \in \calU$, recall
\[
  R = J(p), \qquad P = Q_\lambda(1 - p) - S, \qquad K = Q_\lambda(z - p) - S - Y(z),
\]
\[
  h = 2\lambda - x, \qquad k = p + y - \lambda, \qquad U = z - \lambda, \qquad \Delta = h R - k P,
\]
\[
  \theta \;=\; \frac{R\phi(x) + P\phi(y)}{\Delta}, \qquad \alpha \;=\; \frac{k\phi(x) + h\phi(y)}{\Delta},
\]
so that $H(z, Y(z); p, x, y) = -\theta(z - \lambda) - \alpha K$.

\paragraph{Step 1: Almost-everywhere activity above $c$.}
Let $c \in [\ell, t_e]$ be the last rank where $J(c) = C(c)$. Because $C'(r) - Q'_\lambda(r) = r - \lambda/2 - \lambda S/(1 - r)^2$ is strictly increasing on $[\ell, t_e]$, contact cannot occur at a point where $C'(r) > Q'_\lambda(r)$: since $J' \le Q'_\lambda$, that would make $C - J$ negative immediately to the left of $c$. Thus $0 < C'(r) \le \lambda - r$ on $[\ell, c]$. By \cref{lem:first-order}, almost every rank $r \in (c, t_e]$ either satisfies $\phi(r) = \lambda - r$ or belongs to $I = [x_-, x_+]$ and has a sequence of triples in $\calU$ with $x = r$ or $y = r$ along which $\Gamma(z, p, x, y) \to 0$, and in addition $\phi(r) \ge \varepsilon_0$ almost everywhere on $(c, t_e]$ and $\phi(r) = \lambda - r$ almost everywhere on $(x_+, t_e]$. In Steps~2 and~3 below, we work with these sequences of approximate triples until continuity of $\phi$ is proved.

\paragraph{Step 2: Setting $J = C$ on $[\ell, c]$.}
Let $(S, \phi_0)$ be a minimizing pair with primitive $J_0$ and continuation $Y_0$, and let $c$ be the last contact of $J_0$ with $C$. Since $0 < C'(r) \le \lambda - r$ on $[\ell, c]$ and $J_0(c) = C(c)$, replacing $J_0$ by $C$ on $[\ell, c]$ and keeping $\phi_1 = \phi_0$ on $(c, t_e]$ defines a pair $(S, \phi_1) \in \calK_\lambda$ with primitive $J_1 \ge J_0$.

To see that the continuation of $(S, \phi_1)$ lies at or above $Y_0$, note first that if $R, P, K, y, \phi(y)$ are held fixed and $x, \phi(x)$ are changed by $\delta x, \delta\phi(x)$, clearing denominators in the definitions of $H$ and $\theta$ gives the exact identities
\begin{equation}
  H_{\mathrm{new}} - H \;=\; -\frac{c_x}{\Delta_{\mathrm{new}}}\bigl(\delta\phi(x) + \theta\,\delta x\bigr),
  \qquad
  \theta_{\mathrm{new}} - \theta \;=\; \frac{R}{\Delta_{\mathrm{new}}}\bigl(\delta\phi(x) + \theta\,\delta x\bigr),
  \label{eq:control-change-identities}
\end{equation}
and the same identities hold for a change in $y, \phi(y)$ with $(c_x, R)$ replaced by $(c_y, P)$. Now evaluate the right-hand side of \cref{eq:high-ode} for $(S, \phi_1)$ at the old continuation $Y_0(z)$. If that essential infimum equals $\lambda - z$, it is automatically at least $Y'_0(z)$. Otherwise, choose an $\varepsilon$-minimizing triple $(p, x, y) \in \calU$ for $(S, \phi_1)$ at $(z, Y_0(z))$ with $0 < \varepsilon \le 10^{-4}$ small enough that $H_1(z, Y_0(z); p, x, y) < \lambda - z$. The bounds $\phi_1(x) \le \lambda - x$, $\phi_1(y) \le \lambda - y$, and $\lambda R_1 - p P > 0$ give $0 \le \theta_1 < 1$. Comparing $x$ with the average of $\phi_1$ over $[0.34, 0.35]$ using \cref{eq:control-change-identities} and $\Delta_1/c_x < 0.116/0.0145 = 8$ (\cref{lem:variational-estimates}) gives
\[
  \phi_1(x) \;\ge\; \frac{1}{0.01}\int_{0.34}^{0.35} \phi_1(t)\,\mathrm{d}t - 0.035 - 8\varepsilon \;>\; 0.1106 - 8\varepsilon.
\]
Combining this with $\phi_1(y) > 0.0332$, $R_1 > 0.0826$, $P > 0.11379$, and $\Delta_1 < 0.104$ from \cref{lem:variational-estimates} gives
\[
  \theta_1 \;>\; \frac{(0.1106 - 0.0008)\cdot 0.0826 + 0.0332 \cdot 0.11379}{0.104} \;=\; \frac{3211827}{26000000} \;>\; 0.116,
\]
whereas across the entire low-rank interval $[\ell, t_e]$, we have
\[
  \sup_{[\ell, t_e]} |C''| \;\le\; \frac{2 \cdot 0.562330 \cdot 0.023354}{(1 - 0.385)^3} \;<\; 0.113 \;<\; 0.116 \;<\; \theta_1.
\]
If $x < c$ (or $y < c$), where $\phi_1(r) = C'(r)$, choose Lebesgue points $t_n \uparrow c$ of the old function $\phi_0$ with $\limsup_{n \to \infty} \phi_0(t_n) \ge C'(c)$ (which exist because $J_0(c) - J_0(c - h) \ge C(c) - C(c - h)$ for $h > 0$). Since $\theta_1 > \sup_{[\ell, c]} |C''|$, for large enough $n$ we have
\[
  \phi_0(t_n) + \theta_1 t_n \;>\; C'(r) + \theta_1 r \;=\; \phi_1(r) + \theta_1 r.
\]
By \cref{eq:control-change-identities}, replacing each coordinate $r < c$ in $(x, y)$ by $t_n$ (one coordinate at a time) strictly lowers $H$, raises $\theta$, and preserves $y \ge p$. Finally, replacing $R_1 = J_1(p)$ by $R_0 = J_0(p) \le R_1$ lowers $H$ further since $H_R > 0$. Thus $H_1(z, Y_0(z); p, x, y) \ge H_0(z, Y_0(z); p, x', y') \ge Y'_0(z)$. Letting $\varepsilon \downarrow 0$ shows that the continuation $Y_1$ of $(S, \phi_1)$ satisfies $Y_1 \ge Y_0$, so $(S, \phi_1)$ is also a minimizing pair with $j_\lambda(S, \phi_1) = 0$. Reapplying \cref{lem:first-order} to $(S, \phi) := (S, \phi_1)$, we may henceforth assume $J = C$ on $[\ell, c]$.

In addition, $c \ge x_-$: if $c < x_-$, then \cref{lem:first-order} would give $\phi(r) = \lambda - r = Q'_\lambda(r)$ almost everywhere on $(c, x_-)$, which is impossible because on $[\ell, x_-]$ we have
\[
  C'(r) - Q'_\lambda(r) \;\le\; 0.285 - \frac{0.562325}{2} - \frac{0.562325\,S_-}{0.72^2} \;<\; 0,
\]
so $J(r) = C(c) + \int_c^r Q'_\lambda(u)\,\mathrm{d}u > C(r)$ immediately to the right of $c$.

\paragraph{Step 3: Supporting lines and Lipschitz continuity.}
For almost every $r \in (c, x_+)$ with $\phi(r) < \lambda - r$, \cref{lem:first-order} gives a sequence of triples in $\calU$ with $x = r$ or $y = r$ along which $\Gamma(z, p, x, y) \to 0$. Passing to a subsequence, their slopes $\theta$ converge to a limit in $[\theta_\circ, 1)$ (where $\theta_\circ = 0.116$ as in Step~2, and $\theta < 1$ uniformly since $\Delta - (R\phi(x) + P\phi(y)) \ge \lambda R - p P > 0$). If $r$ is the first position ($x = r$), comparing with any Lebesgue point $s \in I$ using \cref{eq:control-change-identities} and $H_{\mathrm{new}} \ge Y'(z)$ while the old gap tends to zero gives
\[
  \phi(s) + \theta s \;\le\; \phi(r) + \theta r.
\]
If $r$ is the second position ($y = r$), the same comparison holds for every Lebesgue point $s \in [r, x_+]$ because $s \ge r \ge p$ is admissible as a second position; at ranks where $\phi(r) = \lambda - r$, it holds for $s \ge r$ with $\theta = 1$ since $\phi(s) \le \lambda - s$. Thus almost every $r \in (c, x_+)$ has a right supporting line of slope $-\theta \in [-1, -\theta_\circ]$, and modifying $\phi$ on a null set to $\phi(r) := \operatorname{ess\,sup}_{s \in [r, x_+]} \phi(s)$ makes $\phi$ nonincreasing on $(c, x_+)$.

This right supporting line already rules out $c \le c_{\mathrm{safe}} := 0.2854 > x_-$. Indeed, if $c \le c_{\mathrm{safe}}$, then $\phi(c+) \le C'(c)$ (since $J(c) = C(c)$ and $J \le C$ to the right of $c$), and the right supporting line at $r \downarrow c$ gives $\phi(r) \le \phi(c+) - \theta_\circ(r - c) \le C'(c) - \theta_\circ(r - c)$ for almost every $r \in (c, c_{\mathrm{ref}}]$. Since $|C''| \le K_\circ < \theta_\circ$ on $[\ell, c_{\mathrm{ref}}]$, integrating from $c$ to $c_{\mathrm{ref}}$ gives
\[
  C(c_{\mathrm{ref}}) - J(c_{\mathrm{ref}}) \;\ge\; \frac{\theta_\circ - K_\circ}{2}(c_{\mathrm{ref}} - c_{\mathrm{safe}})^2,
\]
contradicting the reference bound on $C(c_{\mathrm{ref}}) - J(c_{\mathrm{ref}})$ in \cref{lem:variational-estimates}. Hence $c > c_{\mathrm{safe}} > x_-$.

To rule out downward jumps and prove a left Lipschitz bound on $(c, x_+)$, consider a rank $r \in (c, x_+)$ where the approximate triples have $y = r$, and compare with a Lebesgue point $s \in [x_-, r)$ along the coupled shift $(p, y) \mapsto (\min\{p, s\}, s)$, which remains in $\calU$ because $s \ge x_- > p_-$. First shift the ranks $(p, y)$ to $(\min\{p, s\}, s)$ while holding the value at the second position fixed at $\phi(r)$, and then replace $\phi(r)$ by $\phi(s)$. Because $J$ is Lipschitz, all coefficients are bounded, and $\Delta > 0.0519$, the geometric shift changes $H$ by $O(r - s)$, while the change in $\phi$ multiplies $\phi(s) - \phi(r)$ by $-c_y/\Delta$ (which is bounded away from zero). Since the resulting triple satisfies $H \ge Y'(z)$ and the original gap tends to zero, letting the gap tend to zero gives $\phi(s) - \phi(r) \le L(r - s)$ for a uniform constant $L$. Explicitly, using $P_p = 1 - p - \lambda$, $K_p = -(\lambda - z + p)$, and $H_k = -(c_y/\Delta)\theta$, the geometric partial derivatives of $H$ with respect to $p$ and $y$ (at fixed values of $\phi(x)$ and $\phi(y)$) are
\begin{equation}
  \frac{H_{p, \mathrm{geom}}}{c_y/\Delta} \;=\; \alpha\Bigl[\phi(p) - \frac{c_x}{c_y}(1 - p - \lambda) + (\lambda - z + p)\frac{\Delta}{c_y}\Bigr] - \theta,
  \qquad
  \frac{H_{y, \mathrm{geom}}}{c_y/\Delta} \;=\; -\theta,
  \label{eq:mass-geometric-derivative}
\end{equation}
which are bounded in \cref{app:variational-estimates}. At ranks where $x = r$, the two-sided supporting line $\phi(s) + \theta s \le \phi(r) + \theta r$ for $s < r$ directly bounds $\phi(s) - \phi(r) \le \theta(r - s)$, and at ranks where $\phi(r) = \lambda - r$, the bound $\phi(s) \le \lambda - s$ gives $\phi(s) - \phi(r) \le r - s$. Together with $\phi(s) \ge \phi(r)$ for $s < r$, this proves that the monotone representative $\phi$ is Lipschitz on $(c,x_+)$ and has a Lipschitz extension to its closure, and since $c > x_-$, taking $s < c$ in the left Lipschitz bound rules out any downward jump at $c$, so $\phi(c) = C'(c)$ and the two pieces are Lipschitz across $c$. For the control comparisons we use this left trace at $x_+$. The continuation on $(x_+,t_e]$ already equals $\lambda-r$ almost everywhere.

Because this Lipschitz representative agrees almost everywhere with the original $\phi$, it agrees at every interior Lebesgue point. Compactness of $[z_0, z_1] \times \calU$ and continuity of $(z, p, x, y) \mapsto H(z, Y(z); p, x, y) - Y'(z)$ now upgrade the sequences of approximate triples from Step~1 to attained tight triples satisfying $H(z, Y(z); p, x, y) = Y'(z)$.

\paragraph{Step 4: Slope bounds between $c$ and the line $\lambda - r$.}
By Rademacher's theorem and Step~3, almost every $r \in (c, x_+)$ with $\phi(r) < \lambda - r$ is a differentiability point of $\phi$ and belongs to an attained tight triple $(p, x, y) \in \calU$ at some $z \in [z_0, z_1]$. Applying Fermat's condition at $r$ gives three cases (excluding the isolated endpoint values of $p, x, y$):
\begin{itemize}
  \item At almost every rank active as $x$, or as $y$ with $y > p$, interior maximization of $\phi(r) + \theta r$ gives $\phi'(r) = -\theta$.
  \item At a boundary second position $y = p$ with $x > p = y$, since $x$ maximizes $\phi(r) + \theta r$ on $[x_-, x_+]$ and $y = p$ maximizes $\phi(r) + \theta r$ on $[p, x_+]$ (which contains $x$), we have $\phi(p) + \theta p \ge \phi(x) + \theta x \ge \phi(r) + \theta r$ for all $r \in [x_-, x_+]$; thus $p$ is also an interior maximizer on $[x_-, x_+]$ and $\phi'(p) = -\theta$.
  \item At a boundary second position $y = p$ with $x \le p = y$, differentiating $H(z, Y(z); p, x, p)$ with respect to $p$ along the coupled boundary $y = p$ using \cref{eq:mass-geometric-derivative} gives
  \begin{equation}
    \phi'(p) \;=\; L_{\mathrm{diag}} \;:=\; \alpha\Bigl[\phi(p) - \frac{c_x}{c_y}(1 - p - \lambda) + (\lambda - z + p)\frac{\Delta}{c_y}\Bigr] - 2\theta,
    \label{eq:diagonal-slope}
  \end{equation}
  while the right supporting line on $[p, x_+]$ gives $\phi'(p) \le -\theta$.
  \item Where $\phi(r) = \lambda - r$, we have $\phi'(r) = -1$.
\end{itemize}
At $r = c$, the slack $\sigma(r) = \lambda - r - \phi(r)$ satisfies $0 \le \sigma(c) < \sigma_0$. By \cref{lem:variational-estimates}, $L_{\mathrm{diag}} > -1$ whenever $0 \le \sigma(p) \le \sigma_0$ for $p \in [x_-, p_b]$. Hence $\sigma'(r) = -1 - \phi'(r) \le 0$ almost everywhere, so $\sigma$ is nonincreasing and remains in $[0, \sigma_0]$ up to $p_b$. On $[p_{\mathrm{ref}}, p_+]$ (where $p_{\mathrm{ref}} < p_b$), monotonicity of $\phi$ implies $\sigma(r + u) \ge \sigma(r) - u$, so any value $\sigma(r) \ge \epsilon$ produces a triangle of slack area at least $\epsilon^2/2 > \varepsilon_{\mathrm{area}}$, contradicting the reference bound $\int_{p_{\mathrm{ref}}}^{t_e} \sigma(r)\,\mathrm{d}r < \varepsilon_{\mathrm{area}}$ from \cref{lem:variational-estimates}.

Using $0 \le \sigma \le \sigma_0$, we sharpen the lower bound on $\theta = 1/\eta$: for any known lower bound $t_0 < c$, replacing $\max_x \{x + \eta\phi(x)\}$ and $\max_{y \ge p} \{y + \eta\phi(y)\}$ in \cref{eq:eta-balance} by their values at $t_0$ (where $\phi(t_0) = C'(t_0) =: f_0$) and $\max\{t_0, p\}$ (using $\phi(p) \ge \lambda - p - \sigma_0$ when $p \ge t_0$) gives, with $w = R/P \le \lambda/(1 - p)$,
\[
  \theta \;\ge\;
  \begin{cases}
    \dfrac{f_0(w + 1)}{(2\lambda - t_0)w + \lambda - p - t_0}, & p \le t_0, \\[10pt]
    \dfrac{f_0 w + \lambda - p - \sigma_0}{(2\lambda - t_0)w + \lambda - 2p}, & p \ge t_0.
  \end{cases}
\]
Because $\sigma'(r) \ge -(1 - \theta_{\min})$, the slack area between $c$ and the line $\lambda - r$ satisfies
\[
  S - Q_\lambda(c) + C(c) \;=\; \int_c^{t_e} \sigma(r)\,\mathrm{d}r \;\ge\; \frac{\sigma(c)^2}{2(1 - \theta_{\min})}.
\]
Iterating between this area bound and the lower bound on $\theta$ (\cref{app:variational-estimates}) gives $c > \underline{c}$ and $\underline{\theta} < \theta < \overline{\theta} < 1$. The improved slope bound $\sigma' \ge -(1 - \underline{\theta})$ and the integral bound on $[p_{\mathrm{ref}}, t_e]$ then imply $\sigma(p_{\mathrm{ref}}) < \varepsilon_{\mathrm{low}}$, on which $L_{\mathrm{diag}} > -1 + \gamma_{\mathrm{tail}}$ with $\gamma_{\mathrm{tail}} > 0$. Thus $\sigma(r)$ decreases at rate at least $\gamma_{\mathrm{tail}}$ whenever positive, hits $0$ by $x_c^+ := p_{\mathrm{ref}} + \varepsilon_{\mathrm{low}}/\gamma_{\mathrm{tail}} < p_+$, and remains zero up to $p_+$. Above $p_+$ every active second position lies strictly above the mass, so at an unsaturated differentiability point $\phi'=-\theta>-1$. Thus $\sigma$ cannot increase there either. The score agrees with $\lambda-r$ through $x_+$ and joins its prescribed value on $(x_+,t_e]$ continuously. Finally, $\int_c^{t_e} \sigma(r)\,\mathrm{d}r \le \sigma(c)(p_{\mathrm{ref}} - c) + \varepsilon_{\mathrm{area}}$ rules out $c \ge \overline{c}$.

\paragraph{Step 5: Continuing into the high ranks.}
For each mass $p \in [p_-, p_+]$, let $\eta, \beta$ be defined by \cref{eq:eta-balance,eq:beta-def}, and set $\theta = 1/\eta$ and $\alpha = \beta/\eta$. At the maximizing positions $x, y$ of $\phi(r) + \theta r$, we have $\phi(x) = \theta(2\lambda - x) - \alpha P$ and $\phi(y) = \alpha R - \theta(p + y - \lambda)$. Multiplying the supporting inequalities $\phi(r_1) + \theta r_1 \le \phi(x) + \theta x$ and $\phi(r_2) + \theta r_2 \le \phi(y) + \theta y$ by $c_x > 0$ and $c_y > 0$ shows that the minimum of $H(z, Y(z); p, r_1, r_2)$ over $(r_1, r_2)$ is
\[
  H_p(z) \;:=\; -\theta(z - \lambda) - \alpha\bigl[Q_\lambda(z - p) - S - Y(z)\bigr],
\]
where $\theta$ and $\alpha$ depend only on $p$. By \cref{lem:variational-estimates}, $\theta + \alpha p > 1 + \gamma_{\mathrm{high}}$ with $\gamma_{\mathrm{high}} > 0$, so along $Y(z)$ (where $Y'(z) \le \lambda - z$),
\[
  \frac{\mathrm{d}}{\mathrm{d}z} H_p(z) \;=\; -\theta + \alpha\bigl[Y'(z) - (\lambda - z + p)\bigr] \;\le\; -\theta - \alpha p \;<\; -1 - \gamma_{\mathrm{high}}.
\]
Because $Y'$ is Lipschitz and every $H_p(z)$ decreases at rate at least $1 + \gamma_{\mathrm{high}}$ while $\lambda - z$ decreases at rate $1$, their point-wise minimum $Y'(z) = \min\bigl\{\lambda - z, \inf_{p \in [p_-, p_+]} H_p(z)\bigr\}$ satisfies $Y''(z) \le -1$ almost everywhere. Thus each $H_p$ satisfies $H''_p(z) = \alpha(Y''(z) + 1) \le 0$ almost everywhere and is therefore concave on $[z_0, z_1]$, so $Y'(z)$ is also concave on $[z_0, z_1]$ and cannot return to $\lambda - z$ once it leaves. If $Y'(z)$ left $\lambda - z$ at or before $z_- := 0.782$, the slope bound $Y'(z) - (\lambda - z) \le -\gamma_{\mathrm{high}}(z - z_-)$ would give $Q_\lambda(z_{\mathrm{ref}}) - S - Y(z_{\mathrm{ref}}) \ge \frac{1}{2}\gamma_{\mathrm{high}}(z_{\mathrm{ref}} - 0.782)^2$, contradicting the reference bound in \cref{lem:variational-estimates}; conversely, the fixed control in \cref{lem:completion-bounds} shows that $Y'(z)$ leaves $\lambda - z$ before $0.788066162 < z_+ < z_B$.
\end{proof}

\subsection{Comparing Every Feasible Function with a Reference Function}
\label{sec:reference-comparison}

We now compare any pair $(S, \phi)$ satisfying the low-rank and election constraints with an explicit finite-dimensional reference family $J_{\mathrm{ref}, \lambda}$ (constructed at each tested value of $\lambda$ in \cref{app:calibration}). This single comparison excludes feasibility at $\underline{\lambda} = 0.56232844$ and gives the localization bounds at $\lambda_{\mathrm{ref}} = 0.5623285$ used in \cref{prop:attainment,prop:shape}. Because it is used before \cref{prop:shape}, the argument below relies only on the variational and election inequalities (and does not assume any regularity of $\phi$); the fixed test split $q_{\mathrm{ref}}$ from \cref{app:calibration} is independent of the minimizer's contact split $q_\star$. When $\lambda$ is fixed, we drop the subscript $\lambda$ from $J_{\mathrm{ref}, \lambda}$ and $u_\lambda$.

\paragraph{Step 1: Equal-primitive rank matching and logarithmic ODE.}
Fix $\lambda$ and a competing primitive $J$ with derivative $\phi = J'$ satisfying $Q_\lambda - S \le J \le C_{\lambda, S}$, $0 \le \phi(r) \le \lambda - r$, $J(t_e) = Q_\lambda(t_e) - S$, and the two-class and three-class inequalities. Write $a := c_{\mathrm{ref}}$ and $b := p_{\mathrm{ref}}$. Using the gauge function $u = u_\lambda$ from \cref{app:calibration}---which agrees with $-\lambda/(1 - r)$ on $[\ell, a]$ and equals $-1$ at $b$---define the shifted reference primitive
\[
  \Xi_S(r) \;:=\; J_{\mathrm{ref}}(r) + (S - S_{\mathrm{ref}})\,u(r),
  \qquad
  m_S(r) \;:=\; \Xi'_S(r) \;>\; 0,
\]
and define the matched rank $\gamma(r)$ as the first rank in $[a, b]$ where $J$ reaches $\Xi_S(r)$:
\[
  \gamma(r) \;:=\; \inf\bigl\{s \ge a : J(s) \ge \Xi_S(r)\bigr\}.
\]
We take the right-continuous representative of $\gamma$ on $(a, b)$, with endpoint traces $\gamma(a-) := a$ and $\gamma(b) := \lim_{r \uparrow b} \gamma(r)$. Because $J(a) \le C_{\lambda, S}(a) = \Xi_S(a)$ and $J(b) \ge Q_\lambda(b) - S = \Xi_S(b)$, we have $a \le \gamma(r) \le b$ for all $r \in [a, b]$ and $\gamma(b) \le b$; any initial jump $\gamma(a+) - a \ge 0$ (corresponding to $J(a) < \Xi_S(a)$ or an initial plateau of $J$) is kept as an atom of the derivative measure $\mathrm{D}\gamma$ at $a$.

Because $m_S(r) \ge m_0 > 0$ on $[a, b]$ and $J$ has Lipschitz constant $M_0 < \infty$, we have $M_0\bigl(\gamma(r_2) - \gamma(r_1)\bigr) \ge J(\gamma(r_2)) - J(\gamma(r_1)) = \Xi_S(r_2) - \Xi_S(r_1) \ge m_0(r_2 - r_1)$ for all $a \le r_1 < r_2 \le b$. Thus $\gamma$ is strictly increasing with a Lipschitz inverse, its absolutely continuous derivative $\gamma'_{\mathrm{ac}}(r)$ is bounded below by $m_0/M_0 > 0$, and the monotone change-of-variables theorem gives, for almost every $r \in [a, b]$,
\begin{equation}
  \phi\bigl(\gamma(r)\bigr) \;=\; \frac{m_S(r)}{\gamma'_{\mathrm{ac}}(r)},
  \label{eq:reciprocal-score}
\end{equation}
while the singular part $\mathrm{D}\gamma_{\mathrm{sing}}$ is kept separately. Below, we write $\gamma'(r)$ for $\gamma'_{\mathrm{ac}}(r)$.

To apply the three-class bound along the one-dimensional control paths $y = x$ and $y = p$ before regularity of $\phi$ is known, observe that whenever $x, y \in (x_-, x_+)$ are Lebesgue points of $\phi$ with $y \ge p$, we can approximate $(p, x, y)$ by full-dimensional sets of interior triples in $\calU$ whose values of $\phi$ converge to $\phi(x)$ and $\phi(y)$ (including when $y = x$ or $y = p$). Continuity of $J(p)$ and uniform positivity of $\Delta$ then pass the essential-infimum bound $Y'(z) \le H(z, Y(z); p, x, y)$ to $(p, x, y)$. Because $\gamma^{-1}$ is Lipschitz and the reference control maps $\pi(s), \upsilon(s)$ below are piecewise affine and strictly increasing, the preimage of the null set of non-Lebesgue points of $\phi$ has measure zero, so the three-class differential inequality holds almost everywhere along $y = x = \gamma(\upsilon(s))$ and $y = p = \gamma(\pi(s))$.

Along the reference high-rank interval $[\zeta_0, \zeta_1]$, we fix strictly increasing, piecewise-affine control maps $\pi(s)$ (for the class mass) and $\upsilon(s)$ (for the first low rank). For $\pi(s) \ge c_{\mathrm{ref}}$, we test the competitor at class mass $p = \gamma(\pi(s))$ and first low rank $x = \gamma(\upsilon(s))$; for $\pi(s) < c_{\mathrm{ref}}$, we keep $p = \pi(s)$ fixed and replace $R = J(p)$ by its upper bound $C_{\lambda, S}(p)$ (which weakens the three-class inequality since $\partial E_3/\partial R = U\phi(x) + h\phi(z) \le -\lambda U < 0$). Similarly, in the high-rank two-class inequality at $q_{\mathrm{ref}} < c_{\mathrm{ref}}$, replacing $J(q_{\mathrm{ref}})$ by $C_{\lambda, S}(q_{\mathrm{ref}})$ weakens the lower bound.

As $S$ varies, we rescale the high-rank variable $s \in [\zeta_0, \zeta_1]$ affinely by
\[
  z_S(s) \;:=\; s - \frac{s - \zeta_0}{\zeta_1 - \zeta_0}(S - S_{\mathrm{ref}}),
  \qquad
  j_S \;:=\; z'_S(s) \;=\; 1 - \frac{S - S_{\mathrm{ref}}}{\zeta_1 - \zeta_0} \;>\; 0.
\]
Choosing $M = 0.135$ so that $Z(s) := M - S - Y(z_S(s)) > 0$, we make the logarithmic change of variable $w(s) := \log Z(s)$, so $Y'(z_S(s)) = -\frac{\mathrm{e}^{w(s)}}{j_S}w'(s)$. Put $P = Q_\lambda(1 - p) - S$ and $K_0 = Q_\lambda(z_S(s) - p) - M$, so that $K = Q_\lambda(z_S(s) - p) - S - Y(z_S(s)) = K_0 + \mathrm{e}^{w(s)}$. Substituting $\phi(\gamma(r)) = m_S(r)/\gamma'(r)$ and $K = K_0 + \mathrm{e}^{w(s)}$ into the three-class differential inequality $\Delta\,Y'(z_S(s)) \le -(U R + k K)\phi(x) - (U P + h K)\phi(y)$ and dividing by $-\Delta\,\mathrm{e}^{w(s)}/j_S$ splits the lower bound on $w'(s)$ into $1/\gamma'$ and $\mathrm{e}^{-w(s)}/\gamma'$ terms:
\begin{itemize}
  \item On subintervals where the two low positions coincide ($y = x = \gamma(\upsilon)$, so $k = p + \gamma(\upsilon) - \lambda$, $h + k = \lambda + p$, and $\Delta = h R - k P$), the differential inequality for $w'(s)$ is a sum of two positive terms:
  \[
    \frac{j_S(\lambda + p)\,m_S(\upsilon)}{\Delta\,\gamma'(\upsilon)}
    \;+\;
    \frac{j_S\bigl[U(R + P) + (\lambda + p)K_0\bigr]\,m_S(\upsilon)}{\Delta\,\mathrm{e}^{w(s)} \gamma'(\upsilon)}.
  \]
  \item On subintervals where the second low position equals the class mass ($y = p = \gamma(\pi)$, so $k = 2p - \lambda$), $w'(s)$ is a sum of four positive terms with numerators
  \[
    j_S k\,m_S(\upsilon), \qquad j_S h\,m_S(\pi), \qquad j_S(U R + k K_0)\,m_S(\upsilon), \qquad j_S(U P + h K_0)\,m_S(\pi),
  \]
  divided by $\Delta\,\gamma'(\upsilon)$, $\Delta\,\gamma'(\pi)$, $\Delta\,\mathrm{e}^w \gamma'(\upsilon)$, and $\Delta\,\mathrm{e}^w \gamma'(\pi)$, respectively.
\end{itemize}
Let $\widehat{w}(s)$ solve this differential equation with initial condition $\widehat{w}(\zeta_0) = \log\bigl(M - Q_\lambda(\zeta_0)\bigr)$, and set $\widehat{Z}(s) := \mathrm{e}^{\widehat{w}(s)}$. Because the competing function $Z_{\mathrm{actual}}(s) = M - S - Y(z_S(s))$ starts at $Z_{\mathrm{actual}}(\zeta_0) \ge \widehat{Z}(\zeta_0)$ and has derivative at least that prescribed by the equation, ODE comparison gives $Z_{\mathrm{actual}}(\zeta_1) \ge \widehat{Z}(\zeta_1)$. At the right endpoint $z_S(\zeta_1)$, the two-class inequality at $q_{\mathrm{ref}}$ requires $Z_{\mathrm{actual}}(\zeta_1) \le Z_{\mathrm{end}}(S)$, where
\[
  Z_{\mathrm{end}}(S) \;:=\; M - S - \frac{z_S(\zeta_1)}{\lambda - q_{\mathrm{ref}}}\bigl[Q_\lambda(1 - z_S(\zeta_1) + q_{\mathrm{ref}}) - S - C_{\lambda, S}(q_{\mathrm{ref}})\bigr].
\]
Thus every feasible competitor must satisfy
\begin{equation}
  \Lambda(S, \gamma) \;:=\; \log\frac{\widehat{Z}(\zeta_1)}{Z_{\mathrm{end}}(S)} \;\le\; 0.
  \label{eq:log-comparison}
\end{equation}

\paragraph{Step 2: Quadratic error bounds from the reciprocal derivative.}
Every term on the right-hand side of the differential equation for $\widehat{w}$ is a positive coefficient multiplied by either $\mathrm{e}^{-w}/\gamma'$ or $1/\gamma'$. To extract a coercive quadratic error in $\gamma' - 1$, we use the inequality
\begin{equation}
  \frac{\mathrm{e}^v}{\xi} - 1 - v + \xi - 1 \;\ge\; \chi(\xi) \;:=\; \xi - 1 - \log\xi \;\ge\; \frac{8}{9}\,\frac{(\xi - 1)^2}{\xi + 1}
  \qquad (\xi > 0,\ v \in \R),
  \label{eq:chi-inequality}
\end{equation}
where the first inequality minimizes over $v$ at $v = \log\xi$, and the second follows because $\chi(\xi) - \frac{8}{9}\frac{(\xi - 1)^2}{\xi + 1}$ vanishes at $\xi = 1$ and has derivative $\frac{(\xi - 1)(\xi - 3)^2}{9\xi(\xi + 1)^2}$.

Write $e(r) := \gamma(r) - r$ for the rank displacement, and define the variation functional
\[
  \calJ \;:=\; \int_{c_{\mathrm{ref}}}^{p_{\mathrm{ref}}} \frac{(\gamma'(r) - 1)^2}{\gamma'(r) + 1}\,\mathrm{d}r + |\mathrm{D}\gamma_{\mathrm{sing}}|\bigl([c_{\mathrm{ref}}, p_{\mathrm{ref}}]\bigr),
\]
where $\mathrm{D}\gamma_{\mathrm{sing}}$ is the singular part of the derivative measure $\mathrm{D}\gamma$ (including the initial atom at $a = c_{\mathrm{ref}}$). To see that $\calJ$ controls $\|e\|_\infty$ on $[a, b]$ with the initial atom at $a$ included, set $\mathfrak{m} := \mathrm{d}r + \mathrm{D}\gamma$ on $[a, b]$. Then $\mathrm{d}e = (\gamma'(r) - 1)\,\mathrm{d}r + \mathrm{D}\gamma_{\mathrm{sing}}$, so
\[
  \int_{[a, b]} \Bigl|\frac{\mathrm{d}e}{\mathrm{d}\mathfrak{m}}\Bigr|^2 \mathrm{d}\mathfrak{m} \;=\; \calJ
  \qquad\text{and}\qquad
  \mathfrak{m}\bigl([a, b]\bigr) \;=\; (b - a) + \gamma(b) - a \;\le\; 2(b - a).
\]
Since $e(a-) = 0$, the Cauchy--Schwarz inequality with respect to $\mathfrak{m}$ on $[a, r]$ gives $\|e\|_\infty^2 \le 2(b - a)\calJ$.

Now write the ODE for $\widehat{w}(s)$ in the common form
\[
  \widehat{w}'(s) \;=\; \sum_i a_i\bigl(s, S, \gamma(r_1(s)), \gamma(r_2(s))\bigr)\,\frac{\mathrm{e}^{-\delta_i \widehat{w}(s)}}{\gamma'(r_i(s))},
  \qquad \delta_i \in \{0, 1\},\quad a_i > 0,
\]
where $a_i$ are the positive coefficients from the two branches in Step~1. Let $w_0(s)$ be the reference solution (at $S = S_{\mathrm{ref}}$ and $\gamma(r) = r$), and define the reference weights and integrating factor
\[
  b_i(s) \;:=\; a_{i, 0}(s)\,\mathrm{e}^{-\delta_i w_0(s)},
  \qquad
  B(s) \;:=\; \sum_i \delta_i b_i(s),
  \qquad
  \rho(s) \;:=\; \exp\!\Bigl(-\int_s^{\zeta_1} B(t)\,\mathrm{d}t\Bigr) \;>\; 0,
\]
where $a_{i, 0}(s) := a_i(s, S_{\mathrm{ref}}, r_1(s), r_2(s))$. Applying \cref{eq:chi-inequality} to each term with $v_i = \log(a_i/a_{i, 0}) - \delta_i(\widehat{w} - w_0)$ and $\xi_i = \gamma'(r_i(s))$ gives
\[
  (\widehat{w} - w_0)'(s) \;\ge\; -B(s)\bigl(\widehat{w}(s) - w_0(s)\bigr) + \sum_i b_i(s)\Bigl[\log\frac{a_i}{a_{i, 0}} - (\xi_i - 1) + \chi(\xi_i)\Bigr].
\]
Multiplying by $\rho(s)$ and integrating over $[\zeta_0, \zeta_1]$ eliminates the linear state term $-B(\widehat{w} - w_0)$. Next, Taylor-expand each $\log(a_i/a_{i, 0})$ to second order along the segment from $(S_{\mathrm{ref}}, r_1(s), r_2(s))$ to $(S, \gamma(r_1(s)), \gamma(r_2(s)))$. Combining the certified second-derivative enclosures with the second derivative of $-\log Z_{\mathrm{end}}(S)$ gives a positive quadratic term $\kappa(S - S_{\mathrm{ref}})^2$ and a rank remainder bounded by $K_{\mathrm{rank}}\,\calJ$.

For the first-order terms, push the positive weights $\rho(s)b_i(s)$ forward along the strictly increasing piecewise-affine maps $r_i(s)$ to form the weight density $W(r) > 0$ on $[a, b]$; push forward the first-order rank partial derivatives of $\log a_i$ to form $f(r)$, and set $P_1(r) := \int_r^b f(t)\,\mathrm{d}t$. Write $A_S$ for the total first derivative with respect to $S$ (including $-\mathrm{d}\log Z_{\mathrm{end}}/\mathrm{d}S$ at $S_{\mathrm{ref}}$). Because $e(a-) = 0$ and $P_1(b) = 0$, integration by parts for functions of bounded variation gives $\int_a^b f(r)e(r)\,\mathrm{d}r = \int_{[a, b]} P_1(r-)\,\mathrm{d}e(r)$, while splitting $\mathrm{d}e(r) = (\gamma'(r) - 1)\,\mathrm{d}r + \mathrm{D}\gamma_{\mathrm{sing}}(r)$ gives the exact measure identity
\[
  -\int_a^b W(r)\bigl(\gamma'(r) - 1\bigr)\,\mathrm{d}r
  \;=\; -\int_{[a, b]} W(r-)\,\mathrm{d}e(r) + \int_{[a, b]} W(r-)\,\mathrm{d}\mathrm{D}\gamma_{\mathrm{sing}}(r),
\]
where at the left endpoint $a$ we set $W(a-) := W(a+)$ and $P_1(a-) := P_1(a)$, and at $r \in (a, b]$ we take the left limit $r-$. Combining these terms gives
\begin{equation}
  \begin{aligned}
    \Lambda(S, \gamma) \;\ge\; \Lambda_{\mathrm{ref}} + A_S(S - S_{\mathrm{ref}})
    &+ \int_{[c_{\mathrm{ref}}, p_{\mathrm{ref}}]} \bigl(P_1(r-) - W(r-)\bigr)\,\mathrm{d}e(r)
    + \int_{c_{\mathrm{ref}}}^{p_{\mathrm{ref}}} W(r)\,\chi\bigl(\gamma'(r)\bigr)\,\mathrm{d}r \\
    &+ \int_{[c_{\mathrm{ref}}, p_{\mathrm{ref}}]} W(r-)\,\mathrm{d}\mathrm{D}\gamma_{\mathrm{sing}}(r)
    + \kappa(S - S_{\mathrm{ref}})^2 - K_{\mathrm{rank}}\,\calJ.
  \end{aligned}
  \label{eq:adjoint-expansion}
\end{equation}

Choose constants $W_0 > 0$ and $k_0 < 0$ such that $W(r) \ge W_0$ and $a(r) := P_1(r) - k_0 - W_0 > 0$ on $[c_{\mathrm{ref}}, p_{\mathrm{ref}}]$. Splitting $\chi(\gamma') = \frac{W_0}{W}\chi(\gamma') + \frac{w_1}{W}\chi(\gamma')$ with $w_1(r) := W(r) - W_0 \ge 0$ and adding and subtracting $k_0\,\mathrm{d}e$, the identity $\log\xi = \log(a\xi/w_1) + \log(w_1/a)$ gives, when $w_1 > 0$ and $\xi = \gamma'(r) > 0$,
\[
  w_1\,\chi(\xi) + (a - w_1)(\xi - 1)
  \;=\; w_1\,\chi\!\Bigl(\frac{a\xi}{w_1}\Bigr) + w_1\,\chi\!\Bigl(\frac{w_1}{a}\Bigr) - \frac{(a - w_1)^2}{a}
  \;\ge\; -\frac{(P_1 - W - k_0)^2}{P_1 - k_0 - W_0},
\]
and when $w_1 = 0$, the same lower bound holds directly since $a(\xi - 1) \ge -a$. Let $\varepsilon$ be an upper bound on the integral of $\frac{(P_1 - W - k_0)^2}{P_1 - k_0 - W_0}$ over $[c_{\mathrm{ref}}, p_{\mathrm{ref}}]$. On the singular part $\mathrm{D}\gamma_{\mathrm{sing}}$, after adding and subtracting $k_0\,\mathrm{d}e$ we have $(P_1(r-) - k_0)\,\mathrm{d}\mathrm{D}\gamma_{\mathrm{sing}}(r) = \bigl(W_0 + a(r-)\bigr)\,\mathrm{d}\mathrm{D}\gamma_{\mathrm{sing}}(r) \ge W_0\,\mathrm{d}\mathrm{D}\gamma_{\mathrm{sing}}(r) \ge \frac{8}{9}W_0\,\mathrm{d}\mathrm{D}\gamma_{\mathrm{sing}}(r)$, while the boundary term from $\int_{[a, b]} k_0\,\mathrm{d}e$ is $k_0\,e(b) = (-k_0)\bigl(b - \gamma(b)\bigr) \ge 0$. Completing the square in $S - S_{\mathrm{ref}}$ with $A_0 \ge |A_S|$ therefore gives
\begin{equation}
  \Lambda(S, \gamma) \;\ge\; \Lambda_{\mathrm{ref}} - \varepsilon - \frac{A_0^2}{4\kappa} + \Bigl(\frac{8}{9}W_0 - K_{\mathrm{rank}}\Bigr)\calJ + (-k_0)\bigl[b-\gamma(b)\bigr].
  \label{eq:coercive-master}
\end{equation}
Furthermore, if the competitor has an initial high-rank deficit $E_{\mathrm{in}} := Q_\lambda(\zeta_0) - S - Y(\zeta_0) \ge 0$ at $\zeta_0$, then $Z_{\mathrm{actual}}(\zeta_0) = \widehat{Z}(\zeta_0) + E_{\mathrm{in}}$, and since the linear ODE for $Z$ has a nonnegative homogeneous coefficient, ODE comparison gives $Z_{\mathrm{actual}}(\zeta_1) \ge \widehat{Z}(\zeta_1) + E_{\mathrm{in}}$, strengthening \cref{eq:log-comparison} to
\[
  \Lambda(S, \gamma) \;\le\; \log\Bigl(1 - \frac{E_{\mathrm{in}}}{Z_{\mathrm{end}}(S)}\Bigr) \;\le\; -\frac{E_{\mathrm{in}}}{Z_{\mathrm{end}}(S)}.
\]

By \cref{lem:numerical-calibration} in \cref{app:calibration}, $\kappa > 68.429$ and $\frac{8}{9}W_0 - K_{\mathrm{rank}} > 0.95775 > 0$, which absorbs the linear terms without requiring the reference curve to be an exact stationary point:
\begin{itemize}
  \item At $\underline{\lambda} = 0.56232844$, the constant $\Lambda_{\mathrm{ref}} - \varepsilon - A_0^2/(4\kappa)$ in \cref{eq:coercive-master} exceeds $9.1128 \times 10^{-7} > 0$, contradicting \cref{eq:log-comparison} and proving that no feasible function exists at or below $\underline{\lambda}$.
  \item At $\lambda_{\mathrm{ref}} = 0.5623285$ (and hence for any minimizer at $\lambda^\star \le \lambda_{\mathrm{ref}}$ using the shift of \cref{app:attainment}), \cref{eq:log-comparison,eq:adjoint-expansion,eq:coercive-master} bound $\calJ < 7.15 \times 10^{-6}$ and $|S - S_{\mathrm{ref}}| < 3.18 \times 10^{-4}$. For $r\in[\ell,x_+]$, the transfer increases the cap gap (see \cref{app:variational-estimates}) and satisfies $\widetilde S-Q_{\lambda+\delta}(r)+\widetilde J(r)=S-Q_\lambda(r)+J(r)+\delta(x_+-r)$. Thus upper bounds on both low-rank deficits transfer back to the original pair. Keeping the favorable boundary term $(-k_0)\bigl[p_{\mathrm{ref}} - \gamma(p_{\mathrm{ref}})\bigr]$, the initial atom at $c_{\mathrm{ref}}$, and the initial deficit $E_{\mathrm{in}}$ gives the tube bounds on $C(c_{\mathrm{ref}}) - J(c_{\mathrm{ref}})$, $\int_{p_{\mathrm{ref}}}^{t_e} \sigma$, and $Q_\lambda(z_{\mathrm{ref}}) - S - Y(z_{\mathrm{ref}})$ used in the proof of \cref{prop:shape}.
\end{itemize}

\section{Verifying the Upper Bound Across All Masses and Ranks}
\label{app:upper-all}

\subsection{Normalized Verifier Conditions, Packing Bounds, and Three-Function Test}
\label{app:verifier-proof}

We first complete the endpoint cases ($p_t = 0$, $p_t = 1$, and $A = 0$) in the proof of \cref{lem:verifier}, restricting attention to voters of positive weight:
\begin{itemize}
  \item If $p_t = 0$, then $q_v(t) = z_v(t) = 0$ for every voter $v$, and setting $L_v(t) := 0$ satisfies \cref{eq:telescoping-target} with $\E_v[L_v(t) - q_v(t)] = 0$.
  \item If $A > 0$ and $p_t = 1$, set
  \[
    \beta(1) \;:=\; 0,
    \qquad
    c(1, x) \;:=\; \lambda - x,
    \qquad
    R \;:=\; \frac{\E_v\bigl[\1_{\{s_v < \lambda\}}(\lambda - s_v)\bigr]}{A} \;\ge\; 1.
  \]
  Then the averaged potential is $W(z) = \lambda - z - R\phi(z)$ for $z > \lambda$, with $W(1) = \lambda - 1 + R \ge \lambda$. Define
  \[
    L_v(t) \;:=\; q_v(t) +
    \begin{cases}
      \lambda - s_v, & s_v < \lambda, \\
      R\phi(s_v), & s_v \ge \lambda,
    \end{cases}
  \]
  so that $\E_v\bigl[L_v(t) - q_v(t)\bigr] = R\,\E_v[\phi(s_v)] \ge 0$. For the first inequality in \cref{eq:telescoping-target}, if $s_v < \lambda$, then $q_v(t) \le s_v < \lambda$ gives $W(q_v(t)) = 0$ and $L_v(t) + W(q_v(t)) = q_v(t) + \lambda - s_v \le \lambda$; if $s_v \ge \lambda$, then $L_v(t) + W(q_v(t)) \le \lambda$ follows from $\phi(s_v) \le 0$ when $q_v(t) \le \lambda$ and from $\phi(s_v) \le \phi(q_v(t))$ when $q_v(t) > \lambda$. For the second inequality in \cref{eq:telescoping-target}, the range $z_v(t) \le z \le s_v$ is nonempty only when $z_v(t) = z = s_v = 1$ (note that $p_t = 1$ does not by itself imply $s_v = 1$, because a candidate of lottery mass zero may lie at an interior rank), in which case $L_v(t) + W(q_v(t)) \le \lambda \le W(1)$ gives $L_v(t) + W(q_v(t)) - W(1) \le 0$.
  \item If $A = 0$, then $\phi(r) > 0$ on $[0, \lambda)$ and $\E_v[\phi(s_v)] \ge 0$ imply $s_v = \lambda$ for every voter $v$. Set $W := 0$. For $0 < p_t < 1$, setting $L_v(t) := q_v(t) + \beta(p_t)\,k_v(t)$ has $\E_v[L_v(t) - q_v(t)] \ge 0$; the first inequality in \cref{eq:telescoping-target} follows from \cref{eq:ver-3} with $q = q_v(t) \le s_v = \lambda$, and the second follows from \cref{eq:ver-4} whenever $z_v(t) \le z \le s_v = \lambda$ is nonempty. For $p_t = 1$, setting $L_v(t) := q_v(t) \le \lambda$ satisfies the first inequality in \cref{eq:telescoping-target}, and the range $z_v(t) \le z \le \lambda$ is empty because $z_v(t) \ge p_t = 1 > \lambda$.
\end{itemize}

To verify \cref{eq:ver-1,eq:ver-2,eq:ver-3,eq:ver-4} of \cref{lem:verifier} across all $p \in (0, 1)$, we rewrite the high-rank conditions in normalized form and record several bounds on the packing integrals $\delta_p(q)$ and $K_p(q, z)$~\cite{CaiEtAl26b}. For each low rank $x \in [0, \lambda)$, define the normalized ratios
\[
  r_x \;:=\; \frac{\lambda - x}{\phi(x)} \;\ge\; 1,
  \qquad
  u_x \;:=\; \frac{c(p, x)}{\phi(x)} \;\ge\; 0.
\]
Dividing \cref{eq:ver-3,eq:ver-4} by $\phi(x) > 0$ and substituting the definition~\eqref{eq:potential-def} of $W_x$ splits the two conditions according to where $q$ and $z$ lie relative to $\lambda$:
\begin{itemize}
  \item In \cref{eq:ver-3}, when $0 \le q \le \min\{p, \lambda\}$, we have $W_x(q) = 0$ and $[-\phi(q)]_+ = 0$, so \cref{eq:ver-3} reduces to $q + \beta(p)\,\delta_p(q) \le \lambda$. When $\lambda < q \le p$ (which occurs only if $p > \lambda$), we have $W_x(q) = \lambda - q - r_x\phi(q)$ and $[-\phi(q)]_+ = -\phi(q)$, so \cref{eq:ver-3} becomes
  \begin{equation}
    \beta(p)\,\delta_p(q) + (r_x - u_x)\bigl[-\phi(q)\bigr] \;\le\; 0
    \qquad (\lambda < q \le p).
    \label{eq:ver-3-norm}
  \end{equation}
  \item In \cref{eq:ver-4}, when $0 \le q \le p \le z \le \lambda$, we have $W_x(q) = W_x(z) = 0$ and $[-\phi(z)]_+ = 0$, so \cref{eq:ver-4} reduces to $q + \beta(p)\,K_p(q, z) \le 0$. When $0 \le q \le \min\{p, \lambda\}$ and $\max\{p, \lambda\} \le z \le 1$, we have $W_x(q) = 0$, $W_x(z) = \lambda - z - r_x\phi(z)$, and $[-\phi(z)]_+ = -\phi(z)$, so \cref{eq:ver-4} becomes
  \begin{equation}
    z + q - \lambda + \beta(p)\,K_p(q, z) + (r_x + u_x)\phi(z) \;\le\; 0
    \qquad \bigl(0 \le q \le \min\{p, \lambda\},\ \max\{p, \lambda\} \le z \le 1\bigr).
    \label{eq:ver-4-lowq}
  \end{equation}
  When $\lambda < q \le p \le z \le 1$, we have $W_x(q) - W_x(z) = z - q - r_x\bigl[\phi(q) - \phi(z)\bigr]$, so the initial $q$ in \cref{eq:ver-4} cancels $-q$ and \cref{eq:ver-4} becomes
  \begin{equation}
    z + \beta(p)\,K_p(q, z) - r_x\bigl[\phi(q) - \phi(z)\bigr] + u_x\phi(z) \;\le\; 0
    \qquad (\lambda < q \le p \le z \le 1).
    \label{eq:ver-4-highq}
  \end{equation}
\end{itemize}

Because $\phi$ is nonincreasing with $\phi(\lambda) = 0$ and $\phi(1) = -1$, its primitive $J$ is concave and nonnegative on $[0, 1]$ with $J(0) = J(1) = 0$. Monotonicity of $\phi$ gives
\begin{equation}
  \begin{aligned}
    \delta_p(q) &\;=\; -\int_0^q \phi(u)\,\mathrm{d}u - \int_{1 - p + q}^1 \phi(u)\,\mathrm{d}u \;\le\; p - q && (0 \le q \le p), \\
    \delta_p(q) &\;=\; \int_q^{1 - p + q} \phi(u)\,\mathrm{d}u \;\le\; (1 - p)\phi(q) \;\le\; 0 && (\lambda \le q \le p), \\
    J(q) &\;\ge\; q\phi(q) && (0 \le q \le \lambda), \\
    K_p(q, z) &\;\le\; -(p - q)\phi(z) - J(q) && (0 \le q \le p \le z \le 1).
  \end{aligned}
  \label{eq:elementary-packing-bounds}
\end{equation}
Furthermore, concavity of $J$ gives $J(q) \ge (q/p)J(p)$ for $0 \le q \le p$, and hence
\begin{equation}
  K_p(q, z) \;\le\; -\frac{q}{p}J(p) - (p - q)\phi(z)
  \qquad (0 \le q \le p \le z \le 1),
  \label{eq:Kp-secant}
\end{equation}
whose right-hand side is affine in $q$ at each fixed $(p, z)$.

For any class mass $p \in (0, \lambda]$, to verify all four inequalities \cref{eq:ver-1,eq:ver-2,eq:ver-3,eq:ver-4} of \cref{lem:verifier}, it is enough to find $\beta, \eta \ge 0$ such that the three test functions
\begin{equation}
  \begin{aligned}
    \Psi_1(q, x) &\;:=\; x + q - 2\lambda + \beta\,\delta_p(q) + \eta\phi(x) && \bigl(0 \le x \le \lambda,\ 0 \le q \le \min\{p, x\}\bigr), \\
    \Psi_2(q, y) &\;:=\; y + q - \lambda + \beta\,K_p(q, y) + \eta\phi(y) && (0 \le q \le p \le y \le \lambda), \\
    \Psi_3(q, z) &\;:=\; z + q - \lambda + \beta\,K_p(q, z) + \eta\phi(z) && (0 \le q \le p,\ \lambda \le z \le 1)
  \end{aligned}
  \label{eq:L-R-H-test}
\end{equation}
are nonpositive on their respective domains, and then set $c(p, x) := [\eta\phi(x) - (\lambda - x)]_+$:
\begin{itemize}
  \item In \cref{eq:ver-1}, when $c(p, x) = \eta\phi(x) - (\lambda - x) > 0$, the left-hand side minus $\lambda$ is $\Psi_1(q, x) \le 0$; when $c(p, x) = 0$, evaluating at $x = \lambda$ (where $\phi(\lambda) = 0$) gives $q + \beta\,\delta_p(q) - \lambda = \Psi_1(q, \lambda) \le 0$. Because $q \le p \le \lambda$, this also proves \cref{eq:ver-3}.
  \item In \cref{eq:ver-2}, when $c(p, y) = \eta\phi(y) - (\lambda - y) > 0$, the left-hand side is $\Psi_2(q, y) \le 0$; when $c(p, y) = 0$, monotonicity of $K_p(q, \cdot)$ gives $q + \beta\,K_p(q, y) \le q + \beta\,K_p(q, \lambda) = \Psi_2(q, \lambda) \le 0$ (which at $q = p$ also gives $\beta J(p) \ge p$). This also proves \cref{eq:ver-4} for $z \le \lambda$.
  \item In \cref{eq:ver-4} for $z \ge \lambda$, we have $r_x + u_x = \frac{\lambda - x + c(p, x)}{\phi(x)} = \max\{r_x, \eta\} \ge \eta$ and $\phi(z) \le 0$, so $(r_x + u_x)\phi(z) \le \eta\phi(z)$ and the left-hand side of \cref{eq:ver-4-lowq} is at most $\Psi_3(q, z) \le 0$.
\end{itemize}

\paragraph{Two reductions for $0 < p \le \lambda$.}
First, suppose $0 < p \le \lambda$, $\beta \ge 0$, and $\eta \le 1 + \beta p$. For any fixed $y \in [p, \lambda]$, \cref{eq:Kp-secant} bounds $\Psi_2(q, y)$ for $q \in [0, p]$ by the affine interpolation between its endpoint bounds at $q = 0$ and $q = p$:
\[
  A_0 \;:=\; y - \lambda + (\eta - \beta p)\phi(y) \;\le\; y - \lambda + \phi(y) \;\le\; 0,
  \qquad
  A_p \;:=\; y + p - \lambda + \eta\phi(y) - \beta J(p).
\]
Thus $A_p \le 0$ for all $y \in [p, \lambda]$ implies $\Psi_2(q, y) \le 0$ for all $q \in [0, p]$ and $y \in [p, \lambda]$.

Second, because the completed function $\phi$ satisfies $\phi(z) = \lambda - z$ on $[\lambda, z_-]$ and $\phi'(r) \ge -1$ almost everywhere on $[0, z_-]$, for every fixed $q \in [0, p]$ we have
\[
  \frac{\partial^2}{\partial z^2}\Psi_3(q, z) \;=\; \beta\bigl[\phi'(z - p + q) + 1\bigr] \;\ge\; 0
  \qquad\text{for almost every } z \in [\lambda, z_-].
\]
Thus $z \mapsto \Psi_3(q, z)$ is convex on $[\lambda, z_-]$ for every fixed $q \in [0, p]$, so $\Psi_3(q, z) \le 0$ on $[0, p] \times [\lambda, z_-]$ follows from $\Psi_3(q, \lambda) = \Psi_2(q, \lambda) \le 0$ and $\Psi_3(q, z_-) \le 0$.

\subsection{Completing \texorpdfstring{$\phi^\star$}{phi*} Near Zero and One}
\label{app:joining-curve}

We now verify \cref{lem:completed-score} for the completed minimizing function $\phi = \phi^\star$ at $\lambda = \lambda^\star$. By \cref{prop:shape} and \cref{lem:completion-bounds} (\cref{app:completion-bounds}), $\phi$ follows $C'_{\lambda, S}$ up to $c \in [c_-, \overline{c}]$, transitions with derivative in $[-1, -\underline{\theta}]$ to the line $\lambda - r$ at a rank below $x_c^+$, and leaves $\lambda - r$ at a high rank in $[z_-, z_+]$.

\paragraph{Joining to $r = 0$.}
At $x_h = 0.2635 < \ell$, define
\[
  \phi_0 \;:=\; \frac{2C_{\lambda, S}(x_h)}{x_h} - C'_{\lambda, S}(x_h),
  \qquad
  s_h \;:=\; \frac{C'_{\lambda, S}(x_h) - \phi_0}{x_h},
\]
and on $[0, x_h]$ set
\begin{equation}
  \phi(r) \;=\; \phi_0 + s_h r,
  \qquad
  J(r) \;=\; \phi_0 r + \frac{1}{2}s_h r^2.
  \label{eq:head-affine}
\end{equation}
Then $J(0) = 0$, and both $J$ and $\phi = J'$ match $C_{\lambda, S}$ at $x_h$.

\paragraph{Joining to $r = 1$.}
At the high-rank end, \cref{prop:attainment} gives an active contact $Y(z_\star) = T_{\lambda, S}(z_\star, q_\star)$. The boundary exclusions in \cref{app:checking-binary,app:completion-bounds} place $(z_\star, q_\star)$ strictly in the interior of $[z_B, z_1] \times \calQ$, where $J(1 - z + q) = Q_\lambda(1 - z + q) - S$, $J(q) = C_{\lambda, S}(q)$, and both first-order partial derivatives of $Y(z) - T_{\lambda, S}(z, q)$ vanish. Writing $h = 1 - z$ and expanding $Q_\lambda(q + h) = Q_\lambda(q) + (\lambda - q)h - \frac{1}{2}h^2$ expresses the two-class bound as
\[
  T_{\lambda, S}(z, q) \;=\; \frac{z}{\lambda - q}\bigl[Q_\lambda(q + h) - S - C_{\lambda, S}(q)\bigr]
  \;=\; z\Bigl(h - \frac{\frac{1}{2}h^2 + E(q)}{\lambda - q}\Bigr),
\]
where $E(q) := S - Q_\lambda(q) + C_{\lambda, S}(q)$ has derivative $E'(q) = C'_{\lambda, S}(q) - (\lambda - q) = -\sigma(q)$. At the contact point $(z_\star, q_\star)$, write
\[
  \sigma_\star \;:=\; \lambda - q_\star - C'_{\lambda, S}(q_\star),
  \qquad
  E_\star \;:=\; S - Q_\lambda(q_\star) + C_{\lambda, S}(q_\star),
  \qquad
  h_\star \;:=\; 1 - z_\star.
\]
Setting $\frac{\partial T_{\lambda, S}}{\partial q}(z_\star, q_\star) = 0$ gives $\frac{\frac{1}{2}h_\star^2 + E_\star}{\lambda - q_\star} = \sigma_\star$; substituting this identity back into $T_{\lambda, S}(z_\star, q_\star)$ and into $\frac{\partial T_{\lambda, S}}{\partial z}(z_\star, q_\star)$ (recalling $\mathrm{d}h/\mathrm{d}z = -1$) gives the coupled tangency identities
\begin{equation}
  h_\star^2 \;=\; 2(\lambda - q_\star)\sigma_\star - 2E_\star,
  \qquad
  Y_\star \;:=\; Y(z_\star) \;=\; z_\star(h_\star - \sigma_\star),
  \qquad
  H_\star \;:=\; Y'(z_\star) \;=\; 1 - 2z_\star - \sigma_\star + \frac{z_\star h_\star}{\lambda - q_\star}.
  \label{eq:tangency-identities}
\end{equation}
For $w = (r - z_\star)/h_\star \in [0, 1]$, we complete $\phi$ on $[z_\star, 1]$ by the quadratic polynomial
\begin{equation}
  \phi(r) \;=\; H_\star + c_1 w + c_2 w^2,
  \qquad
  c_1 \;:=\; -\frac{6Y_\star}{h_\star} - 4H_\star + 2,
  \qquad
  c_2 \;:=\; 3H_\star - 3 + \frac{6Y_\star}{h_\star}.
  \label{eq:tail-quadratic}
\end{equation}
This completion matches $\phi(z_\star) = H_\star$, ends at $\phi(1) = -1$, and integrates to $J(1) = Y_\star + h_\star(H_\star + c_1/2 + c_2/3) = 0$. To compare the one-sided slopes at $z_\star$, choose a fixed control $H_p$ attaining $Y'(z_\star)$. The line branch has already been left. Since $H_p\ge Y'$ on the continuation and they agree at $z_\star$, left difference quotients give
\[
 (Y')'_-(z_\star)\ \ge\ H'_p(z_\star)
 =-\theta+\alpha[z_\star-\lambda-p+H_\star]\ >\ c_1/h_\star ,
\]
where the last inequality is the joining-slope bound in \cref{lem:completion-bounds}. The tail starts with slope $c_1/h_\star<0$ and has $c_2<0$. Hence $\phi$ remains strictly decreasing and concave on $[\lambda,1]$.

In backward coordinates $1 - r$, \cref{eq:tail-quadratic} takes the form
\begin{equation}
  \phi(r) \;=\; -1 + v(1 - r) - \kappa_{\mathrm{tail}}(1 - r)^2,
  \qquad
  J(r) \;=\; (1 - r) - \frac{v}{2}(1 - r)^2 + \frac{\kappa_{\mathrm{tail}}}{3}(1 - r)^3,
  \label{eq:tail-backward}
\end{equation}
with $v > 0$ and $\kappa_{\mathrm{tail}} > 0$ bounded in \cref{lem:completion-bounds}. For any fixed-$p$ differential upper bound $H_p(r) = -\theta(r - \lambda) - \alpha\bigl[Q_\lambda(r - p) - S - J(r)\bigr]$, the derivative of the gap $H_p(r) - \phi(r)$ on $[z_\star, 1]$ is
\begin{equation}
  \bigl(H_p(r) - \phi(r)\bigr)' \;=\; v - \theta - \alpha(p + \lambda) + \bigl[\alpha(v - 1) - 2\kappa_{\mathrm{tail}}\bigr](1 - r) - \alpha\kappa_{\mathrm{tail}}(1 - r)^2.
  \label{eq:gap-derivative}
\end{equation}
For $\alpha\ge0$, the right-hand side of \cref{eq:gap-derivative} is concave in $1-r$ and positive at both endpoints $r = z_\star$ and $r = 1$ for every $p \in [x_h, p_+]$ (\cref{lem:completion-bounds}). If $\alpha<0$, the tail satisfies $\phi(r)\le\lambda-r$ since it starts below this line and has slope below $-1$. Integrating the ceiling from the middle line gives $Q_\lambda(r-p)-S-J(r)\ge p(r-\lambda-p/2)>0$ on the tail. Thus $H_p(r)\ge-\theta(r-\lambda)\ge\lambda-r$. In either case the differential upper bound $\phi(r)\le H_p(r)$ continues to hold across $[z_\star, 1]$ for all $p \in [x_h, p_+]$ (the range in which completing $[0, x_h]$ leaves $J(p)$ unchanged).

\subsection{Checking the Two-Class Inequalities}
\label{app:checking-binary}

Write $h = 1 - p$. For $0 \le q < \lambda$, the two-class inequality~\eqref{eq:two-class} is
\begin{equation}
  \frac{J(q + h) - J(q)}{\lambda - q} \;\le\; \frac{J(p)}{p}.
  \label{eq:binary-ratio}
\end{equation}
Differentiating the left-hand side with respect to $q$ gives a fraction with denominator $(\lambda - q)^2$ and numerator
\begin{equation}
  N(q, h) \;:=\; (\lambda - q)\bigl[\phi(q + h) - \phi(q)\bigr] + J(q + h) - J(q),
  \qquad
  \frac{\partial N}{\partial q}(q, h) \;=\; (\lambda - q)\bigl[\phi'(q + h) - \phi'(q)\bigr].
  \label{eq:N-derivatives}
\end{equation}
At any interior stationary point $q$ where $N(q, h) = 0$ and $q + h \in [t_e, z_0]$ (so that $\phi(q + h) = \lambda - q - h$), \cref{eq:N-derivatives} simplifies the left-hand side of \cref{eq:binary-ratio} to
\[
  \frac{J(q + h) - J(q)}{\lambda - q} \;=\; \phi(q) - \phi(q + h) \;=\; \bigl[\lambda - q - \sigma(q)\bigr] - \bigl[\lambda - q - h\bigr] \;=\; h - \sigma(q).
\]
We use $N(q, h)$ to locate the maximizing split rank $q$ across three regimes of $p$:

\paragraph{1. Small and middle masses ($p \le z_0$).}
For $p \le 0.65$, we show that $F_p(q) := q + \frac{p}{J(p)}\delta_p(q)$ is nonincreasing in $q$ on $[0, \min\{p, \lambda\}]$. Since $\phi(q) \ge \phi_0 - q$ on $[0, \lambda]$ and $\phi(q + 1 - p) \le \lambda - q - 1 + p$, we have:
\begin{itemize}
  \item For $p \le 1 - \lambda$, using $J(p) \le p\phi_0$ gives
  \[
    F'_p(q) \;=\; 1 + \frac{p}{J(p)}\bigl[\phi(q + 1 - p) - \phi(q)\bigr] \;\le\; 1 + \frac{\lambda - 1 + p - \phi_0}{\phi_0} \;\le\; 0.
  \]
  \item For $p \in [1 - \lambda, 0.65]$, where $J(p) = Q_\lambda(p) - S$, the numerator of $F'_p(q)$ is at most
  \[
    J(p) + p(\lambda - 1 + p - \phi_0) \;=\; -\Bigl[p(1 + \phi_0 - 2\lambda) - \frac{p^2}{2} + S\Bigr] \;<\; -0.0138 \;<\; 0
  \]
  by \cref{lem:completion-bounds}.
\end{itemize}
Thus it remains to prove $F_p(0)\le\lambda$, or $pJ(1-p)\le\lambda J(p)$. For $p\le x_h$, the score ceiling integrated from the middle line gives $J(1-p)\le Q_\lambda(1-p)-S$. Substituting \cref{eq:head-affine} into $\lambda J(p)/p-[Q_\lambda(1-p)-S]$ gives the positive head margin in \cref{app:completion-bounds}. For $x_h\le p\le0.64$, use $J(p)\ge Q_\lambda(p)-S$ and $J(1-p)=Q_\lambda(1-p)-S$; the resulting polynomial lower bound on $\lambda J(p)-pJ(1-p)$ is positive by the same certificate. For $0.64\le p\le0.65$, the cap at $1-p\in[0.35,0.36]$ gives
\[
 pJ(1-p)\le pC_{\lambda,S}(1-p)=\lambda[Q_\lambda(p)-S]=\lambda J(p).
\]

For $p \in [0.65, z_0]$, \cref{lem:completion-bounds} gives
\[
  J(p)\,F'_p(0.18) \;=\; J(p) + p\bigl[\phi(0.18 + 1 - p) - \phi(0.18)\bigr] \;<\; -0.006 \;<\; 0.
\]
For $0.18 \le q \le \lambda$, the rank $q + 1 - p \ge 0.405$ lies in $[t_e, 1]$ where $\phi'(q + 1 - p) \le -1 \le \phi'(q)$ almost everywhere, so $F'_p(q)$ is nonincreasing and stays negative on $[0.18, \lambda]$. On the strip $0 \le q \le 0.18 < x_h$, writing $r = 1 - p + q$ and using the identity $C_{\lambda, S}(r) - C_{\lambda, S}(1 - p) = q\bigl[\frac{\lambda}{2} - \frac{\lambda S}{p(p - q)}\bigr]$ together with $J(r)\le C_{\lambda,S}(r)+\varepsilon_C(r)$ gives
\begin{equation}
\begin{aligned}
 &p[J(r)-J(q)]-(\lambda-q)J(p)\\
 &\qquad\le q\Bigl[J(p)+p\Bigl(\frac{\lambda}{2}-\phi_0-\frac{s_h q}{2}\Bigr)
               -\frac{\lambda S}{p-q}\Bigr]+p\,\varepsilon_C(r)\le0 .
\end{aligned}
 \label{eq:linear-q-strip}
\end{equation}
Here $\varepsilon_C(r)\le0$ is an upper bound on $J(r)-C_{\lambda,S}(r)$. The cap suffices through $r=\overline c=0.305$, including the head, which is tangent to the cap at $x_h$ and has $s_h\le C''_{\lambda,S}$. Above $\overline c$ take the minimum of the cap and the two upper bounds
\[
\begin{aligned}
 J(r)&\le C_{\lambda,S}(\overline c)+C'_{\lambda,S}(\overline c)(r-\overline c)
                         -\tfrac12\underline\theta(r-\overline c)^2,\\
 J(r)&\le Q_\lambda(r)-S+\epsilon_E+\epsilon_\sigma(\overline p_{\mathrm{ref}}-r)_+
                         +\tfrac{k}{2}(\overline p_{\mathrm{ref}}-r)_+^2 .
\end{aligned}
\]
The first follows by integrating the score bound from $\overline c$, using $\phi(\overline c)\le C'_{\lambda,S}(\overline c)$ since $\phi(c)=C'_{\lambda,S}(c)$ and $\phi'\le-\underline\theta<C''_{\lambda,S}$ between $c$ and $\overline c$. The second follows by integrating the slack bound backwards from $p_{\mathrm{ref}}$, as in \cref{lem:completion-bounds}. The central binary certificate checks the last inequality in \cref{eq:linear-q-strip} with these bounds, keeping the factor $q$ at $q=0$.

\paragraph{2. Locating the active contact on the original continuation ($p \in [z_0, z_1]$).}
To show that the active contact $(z_\star, q_\star)$ lies in the interior of $[z_B, z_1] \times \calQ$ before attaching the quadratic tail on $[z_\star, 1]$, we compare the original continuation $Y(p)$ on $[z_0, z_1]$ with the completed low-rank primitive $J$. Put $\kappa_H=-s_h$ and use the head polynomial to define the stationary barrier up to its entry into the cap. On the stationary head branch with $q+h$ on the middle line, $\phi(q)=\phi_0-\kappa_H q$, $J(q) = \phi_0 q - \frac{1}{2}\kappa_H q^2$, $\phi(q + h) = \lambda - q - h$, and $J(q+h)=Q_\lambda(q+h)-S$. Hence $2N(q,h)=(1-\kappa_H)(\lambda-q)^2-h^2+J_H$, and the stationary rank is
\[
  q \;=\; \lambda - \sqrt{\frac{h^2 - J_H}{1 - \kappa_H}},
  \qquad
  J_H \;:=\; 2\lambda(\lambda - \phi_0) - 2S - (1 - \kappa_H)\lambda^2,
\]
with two-class bound $T_H(p) := p\bigl[h - \sigma(q)\bigr]$. Write $q_H(p)$ for this stationary rank and $J_H^{\mathrm{pol}}(u)=\phi_0u+s_hu^2/2$ for the head polynomial. For every $u\in\calQ$, expansion at $q_H(p)$ gives the identity
\[
 T_H(p)-T_{\lambda,S}(p,u)=\frac{p}{\lambda-u}
 \left[\frac{1+s_h}{2}(u-q_H(p))^2+C_{\lambda,S}(u)-J_H^{\mathrm{pol}}(u)\right].
\]
The cap and the polynomial are tangent at $x_h$, and $C''_{\lambda,S}\ge s_h$ on $\calQ$. Thus $T_H(p)$ dominates every original terminal barrier on this branch.

Near $z_0$ ($p \le 0.807$), a quadratic bound on $Q_\lambda(p) - S - Y(p)$ certifies $Y(p) > T_H(p)$ (\cref{app:completion-bounds}). On $[0.807, 0.85]$, fixing the control triple $p_f = 0.29$ and $x = y = c_-$ gives the strict differential inequality
\begin{equation}
  Y'(p) \;\le\; -\theta(p - \lambda) + \alpha\bigl[Y(p) - Q_\lambda(p - p_f) + S\bigr] \;<\; T'_H(p) - \alpha\bigl[T_H(p) - Y(p)\bigr].
  \label{eq:barrier-ode-comparison}
\end{equation}
At the entry mass $p_{\mathrm{ent}}\in(0.849,0.85)$ the stationary rank equals $x_h$. Hence $T_H(p_{\mathrm{ent}})=T_{\lambda,S}(p_{\mathrm{ent}},x_h)\le Y(p_{\mathrm{ent}})$. Backward integration of \cref{eq:barrier-ode-comparison}, together with the early quadratic bound, gives $Y(p)>T_H(p)$ for $z_0\le p<p_{\mathrm{ent}}$. The displayed square identity now rules out contact with any original terminal barrier before $p_{\mathrm{ent}}$.

On $\calQ$ the cap stationary level
$2(\lambda-q)\sigma_C(q)-2E(q)$, with $\sigma_C=\lambda-q-C'_{\lambda,S}(q)$, has derivative $2(\lambda-q)\sigma_C'(q)<0$. At $x_h$ the level equals $(1-p_{\mathrm{ent}})^2$, and at $0.28$ it is strictly below $(1-z_1)^2$ by the endpoint sign in \cref{app:completion-bounds}. Thus for $p_{\mathrm{ent}}\le p\le z_1$ the original barrier attains its maximum in $q$ at its stationary rank in $[x_h,0.28)$. At $z_1$ the continuation slope exceeds $-0.388$, whereas the barrier slope at this rank is below $-0.396$. Equality there would force the left derivative of their nonnegative difference to be nonpositive. Consequently $z_B<z_\star<z_1$ and $0.26<x_h\le q_\star<0.28$. All these comparisons concern the original continuation $Y$.

It remains to identify the largest binary ratio before $z_\star$. On the head, $N(q,h)$ is first nondecreasing and then nonincreasing. Indeed, in \cref{eq:N-derivatives}, the slope at $q+h$ equals $s_h$ on the head, is greater than $s_h$ on the cap, and is smaller than $s_h$ on the free and middle pieces. For $x_h\le q<\lambda$, $q+h\ge x_h+1-z_1>x_c^+$, so $\phi'(q+h)\le-1\le\phi'(q)$ almost everywhere. Thus $N$ is nonincreasing there. For $p\le0.85$, the certificate gives $N(x_c^+-h,h)>0$ and the ratio at $x_c^+-h$ exceeds its value at zero. For $0.85\le p\le z_\star$, both $N(0,h)$ and $N(x_h,h)$ are positive. In the former case the maximizing rank has $q+h\ge x_c^+$; in the latter it has $q\ge x_h$. The endpoint sign $N(0.28,h)<0$ places the maximizer below $0.28$, hence $q+h<z_0$ in both cases. The maximizing rank therefore follows the stationary head branch up to $p_{\mathrm{ent}}$ and the cap branch thereafter. The head comparison and \cref{eq:high-binary-compatibility} give \cref{eq:binary-ratio} up to $z_\star$.

\paragraph{3. Above the active contact ($p \in [z_\star, 1]$).}
For $p \ge z_\star$, while the maximizing split $q(p)$ remains on the rational piece $[x_h, c]$ where $J(q) = C_{\lambda, S}(q)$, write
\[
  T_{\mathrm{bin}}(p) \;:=\; \max_q \frac{p}{\lambda - q}\bigl[J(q + 1 - p) - J(q)\bigr] \;=\; p\bigl[1 - p - \sigma(q(p))\bigr],
\]
which has first derivative $T'_{\mathrm{bin}}(p) = 1 - 2p - \sigma(q) + \frac{p(1 - p)}{\lambda - q}$ along the stationary curve $(1 - p)^2 = 2(\lambda - q)\sigma(q) - 2E(q)$ (where $E'(q) = -\sigma(q)$). Differentiating the stationary condition gives
\[
  \frac{\mathrm{d}q}{\mathrm{d}p} \;=\; \frac{1 - p}{(\lambda - q)(-\sigma'(q))},
  \qquad
  -\sigma'(q) \;=\; 1 - \frac{2\lambda S}{(1 - q)^3} \;>\; 0,
\]
and differentiating $T'_{\mathrm{bin}}(p)$ gives
\begin{equation}
  T''_{\mathrm{bin}}(p) \;=\; -2 + \frac{2 - 3p}{\lambda - q} + \frac{p(1 - p)^2}{(\lambda - q)^3\bigl[1 - 2\lambda S/(1 - q)^3\bigr]},
  \label{eq:binary-curvature}
\end{equation}
which stays strictly below $J''(p)$ while $q(p) \in [q_\star, c_-]$ (\cref{lem:completion-bounds}). Because $J$ and $T_{\mathrm{bin}}$ have matching values and first derivatives at $z_\star$, integrating twice gives $J(p) \ge T_{\mathrm{bin}}(p)$ until the rational maximizer $q(p)$ reaches $c_-$ at mass $p_a \in (0.91257, 0.91339)$.

For $p \ge 0.9125$ (which overlaps $[z_\star, p_a]$), we first verify that no split $q < c_-$ can maximize the two-class ratio once $p \ge p_a$. Write $h_a := 1 - p_a$. When $h = 1 - p \ge 0.06765$, any $q \ge c_-$ has $q + h \ge x_c^+ = 0.35945$, so $\phi'(q + h) \le -1 \le \phi'(q)$ and $N(q, h)$ is nonincreasing in $q$; on $[x_h,c_-]$ the score is $C'_{\lambda,S}$. If $q+h$ is still on the cap, $C''_{\lambda,S}(q+h)\le C''_{\lambda,S}(q)$; beyond the cap the slope is at most $-\underline\theta<C''_{\lambda,S}(q)$. Thus $N(q,h)$ is nonincreasing on $[x_h,c_-]$, also for the smaller widths used below. Together with the check on $[0, x_h]$ above, this confirms that the rational stationary curve is the global maximizer up to $p_a$. For $p \ge p_a$, when $0.06765 \le h \le h_a$, the middle-rank formula gives $N(c_-, h) = (h_a^2 - h^2)/2 \ge 0$; when $0.02 \le h \le 0.06765$, the endpoint check in \cref{app:completion-bounds} shows $N(c_-, h) > 0$ (accounting for any slack at $c_- + h$). Because $N(\cdot, h)$ increases and then decreases on $[0, x_h]$ and decreases on $[x_h, c_-]$ with $N(c_-, h) \ge 0$, no split $q < c_-$ can outperform $q = c_-$ for $0.02 \le h \le h_a$.

For splits $q \ge c_-$, when $c_- \le q < c$ the interval checker uses the exact rational formula $E(q) = S - Q_\lambda(q) + C_{\lambda, S}(q)$. Once $q \ge c$, we bound $E(q) = J(q) - Q_\lambda(q) + S = \int_q^{t_e} \sigma(r)\,\mathrm{d}r$ from below using the total slack area $E_c := \int_c^{t_e} \sigma(r)\,\mathrm{d}r = S - Q_\lambda(c) + C_{\lambda, S}(c)$ and the slope bound $\sigma'(r) \in [-k, 0]$ with $k = 1 - \underline{\theta}$. Writing $s = \sigma(c)$, $x_0 = c + E_c/s + s/(2k) \le x_c^+$, and $x_1 = x_0 - s/k$, minimizing $\max\bigl\{E_c - \int_c^q \sigma,\, \int_q^{t_e} \sigma\bigr\}$ over $v_q = \sigma(q)$ (using $\int_c^q \sigma \le s(q - c) - (s - v_q)^2/(2k)$ and $\int_q^{t_e} \sigma \ge v_q^2/(2k)$) gives
\begin{equation}
  E(q) \;\ge\; \frac{k}{2}\bigl[(x_0 - q)_+^2 - (x_1 - q)_+^2\bigr].
  \label{eq:remaining-slack-envelope}
\end{equation}
Since $\phi(r) = \lambda - r - \sigma(r)$, integrating over $[q, q + h]$ gives $J(q + h) - J(q) = (\lambda - q)h - \frac{1}{2}h^2 - \int_q^{q+h} \sigma(r)\,\mathrm{d}r$. Setting $\kappa(h) := \frac{h - J(1 - h)/(1 - h)}{h^2}$ (so that $(1 - h)\kappa(h)$ is affine in $h$ on the quadratic piece) and dividing by $h^2 / (\lambda - q)$ rewrites \cref{eq:binary-ratio} as
\begin{equation}
  h^{-2}\int_q^{q+h} \sigma(r)\,\mathrm{d}r - \kappa(h)(\lambda - q) + \frac{1}{2} \;\ge\; 0.
  \label{eq:tail-binary-integral}
\end{equation}
For $h \in [0.02, 0.0875]$, \cref{eq:remaining-slack-envelope} (together with the exact rational $E(q)$ on $[c_-, c]$ and an endpoint bound for $q \ge 0.337$) and the upper bound on $\int_{p_{\mathrm{ref}}}^{t_e} \sigma$ certify \cref{eq:tail-binary-integral} (\cref{app:completion-bounds}). For $h\le0.02$ and $0.25\le q\le0.337$, $\sigma$ is nonnegative and nonincreasing throughout $[q,q+0.02]$. Thus $h^{-2}\int_q^{q+h}\sigma(r)\,\mathrm{d}r=h^{-1}\int_0^1\sigma(q+ht)\,\mathrm{d}t$ is nonincreasing in $h$, and \cref{lem:completion-bounds} gives $\kappa(h) \le \kappa(0) \le 2.122 + 140(S - 0.02271790)$. It is therefore enough to check \cref{eq:tail-binary-integral} at $h = 0.02$ using this upper bound on $\kappa(h)$, the exact rational $\sigma(r)$ for $r \le c$, and the pointwise lower bound
\begin{equation}
  \sigma(r) \;\ge\; \frac{s(x_0 - r)_+}{s/k + x_c^+ - r}
  \qquad (c \le r < x_c^+),
  \label{eq:sigma-pointwise-floor}
\end{equation}
which follows from $-k \le \sigma' \le 0$: writing $v = \sigma(r)$, we have $\int_c^r \sigma \le s(r - c) - (s - v)^2/(2k)$ and $\int_r^{x_c^+} \sigma \le v(x_c^+ - r)$, so adding the two bounds and dropping $-v^2/(2k)$ gives $s(x_0 - c) - s^2/(2k) = E_c \le s(r - c) - s^2/(2k) + v(s/k + x_c^+ - r)$.

For $h\le0.02$, write $B=2.122+140(S-0.02271790)$. If $q\le0.25$, then $q+h\le0.27$ and the first term of \cref{eq:tail-binary-integral} is at least $\sigma(0.27)/0.02$; the certified scalar bound is $\sigma(0.27)/0.02-B\lambda+1/2>0$. If $0.337\le q\le\lambda$, use $\sigma\ge0$ and $1/2-B(\lambda-0.337)>0$.

Finally, for $q \ge \lambda$ (when $p > \lambda$), concavity of $\phi$ on $[\lambda, 1]$ with $\phi(\lambda) = 0$ and $\phi(1) = -1$ gives $\phi(r) \ge -(r - \lambda)/(1 - \lambda)$ on $[\lambda, 1]$. Integrating from $p = 1 - h$ to $1$ and using $\lambda \ge 1/2$ gives
\[
  J(p) \;=\; -\int_p^1 \phi(r)\,\mathrm{d}r \;\le\; h - \frac{h^2}{2(1 - \lambda)} \;\le\; h - h^2 \;=\; p(1 - p).
\]
Together with $\phi(q + h) - \phi(q) \le -h$ (since $\phi' \le -1$ on $[\lambda, 1]$), this shows that $q + \frac{p}{J(p)}\bigl[J(q + h) - J(q)\bigr]$ is nonincreasing in $q$ on $[\lambda, p]$, reducing the two-class inequality for $q \ge \lambda$ to its already verified value at $q = \lambda$.

\subsection{Verifier Coefficients for Masses \texorpdfstring{$p \le \lambda^\star$}{p <= lambda*}}
\label{app:small-middle-masses}

For $0 < p \le \lambda$, we set $c(p, x) := [\eta\phi(x) - (\lambda - x)]_+$ and verify the three test inequalities in \cref{eq:L-R-H-test}. By the first reduction in \cref{app:verifier-proof}, to verify $\Psi_2(q, y) \le 0$ for all $0 \le q \le p \le y \le \lambda$ and $\Psi_1(0, x) \le 0$ for all $0 \le x \le \lambda$, it is enough to check $\eta \le 1 + \beta p$ and the two low-rank supporting-line inequalities
\begin{equation}
  x + \eta\phi(x) \;\le\; 2\lambda - \beta J(1 - p) \quad (0 \le x \le \lambda),
  \qquad
  y + \eta\phi(y) \;\le\; \lambda - p + \beta J(p) \quad (p \le y \le \lambda).
  \label{eq:upper-low-support}
\end{equation}
Furthermore, by the second reduction in \cref{app:verifier-proof}, once $\Psi_2(q, \lambda) \le 0$ is verified, convexity of $z \mapsto \Psi_3(q, z)$ on $[\lambda, z_-]$ reduces $\Psi_3(q, z) \le 0$ on $[0, p] \times [\lambda, 1]$ to checking $z \in [z_-, 1]$. We specify $(\beta(p), \eta(p))$ across three sub-ranges of $(0, \lambda]$:

\paragraph{1. Variational masses ($p \in [x_h, p_+]$).}
For $p \in [x_h, p_+]$, we take $\eta(p)$ and $\beta(p)$ from the balance equations~\eqref{eq:eta-balance} and~\eqref{eq:beta-def}. By \cref{prop:shape}, $\theta(p) = 1/\eta(p) \in (0.691, 0.78)$. Because $r + \eta\phi(r)$ increases on $[0, c]$ and decreases once $\phi(r) = \lambda - r$, its global maxima over $x \in [0, \lambda]$ and over $y \in [p, \lambda]$ are attained inside $[x_-, x_+]$ and $[\max\{x_-, p\}, x_+]$, so \cref{eq:eta-balance,eq:beta-def} give both inequalities in \cref{eq:upper-low-support}. Evaluating the second inequality in \cref{eq:upper-low-support} at $y = \lambda$ gives $\beta J(p) \ge p$; since $J(p) \le p\phi_0$ and $\phi_0 < 0.443$ (\cref{lem:completion-bounds}), we have
\[
  \beta p \;\ge\; \frac{p}{\phi_0} \;\ge\; \frac{0.2635}{0.443} \;>\; \frac{1}{0.691} - 1 \;\ge\; \eta - 1,
\]
so $\eta \le 1 + \beta p$. For $z \in [z_-, 1]$, the bounds $\phi(q) \ge \phi_0 - q$, $\phi(z - p + q) \le \lambda - z + p - q$, and $\beta(p)(\phi_0 - \lambda + z_- - p) > 1$ (\cref{lem:completion-bounds}) give
\[
  \frac{\partial \Psi_3}{\partial q}(q, z) \;=\; 1 + \beta(p)\bigl[\phi(z - p + q) - \phi(q)\bigr] \;\le\; 1 - \beta(p)(\phi_0 - \lambda + z_- - p) \;<\; 0,
\]
and evaluating at $z = 1$ gives $\frac{\partial \Psi_1}{\partial q}(q, x) < 0$. This reduces $\Psi_1(q, x) \le 0$ and $\Psi_3(q, z) \le 0$ to $q = 0$, where $\Psi_3(0, z) \le 0$ holds on $[z_-, z_\star]$ by the ODE~\eqref{eq:high-ode} and on $[z_\star, 1]$ by \cref{eq:gap-derivative}.

\paragraph{2. Small masses ($p \in (0, x_h]$).}
Write $R = J(p)$ and $\overline{P} = Q_\lambda(1 - p) - S$, and define
\begin{equation}
  \beta(p) \;:=\; \frac{\lambda + p}{R + \overline{P}},
  \qquad
  m(p) \;:=\; 2\lambda - \beta(p)\overline{P} \;=\; \lambda - p + \beta(p)R.
  \label{eq:small-mass-beta}
\end{equation}
We split $(0, x_h]$ into three subintervals (\cref{lem:completion-bounds}):
\begin{itemize}
  \item For $p \in (0, 0.04]$, we replace $\overline{P}$ in \cref{eq:small-mass-beta} by the exact value $P = J(1 - p)$ and set $\eta(p) := 1 + \beta(p)p$. On this range, \cref{lem:completion-bounds} gives $R/p \ge 0.39$, $0.84 \le P/p \le 1$, $\beta p < 0.49$, and $m - \lambda > 0.1289$. Checking the endpoints of the affine piece $[0, x_h]$ and using $\phi(x) \le \lambda - x$ above $x_h$ gives $x + \eta\phi(x) \le \lambda + 0.26\beta p < m(p)$ for all $x \in [0, \lambda]$, proving \cref{eq:upper-low-support} and $\beta R \ge p$. Substituting \cref{eq:Kp-secant} and $\eta = 1 + \beta p$ directly into $\Psi_3(q, z)$ gives, for all $q \in [0, p]$ and $z \in [\lambda, 1]$,
  \[
    \Psi_3(q, z) \;\le\; q\Bigl(1 - \frac{\beta R}{p}\Bigr) + \bigl(z - \lambda + \phi(z)\bigr) + \beta q\phi(z) \;\le\; 0,
  \]
  while $\phi(q) \ge \phi_0 - p \ge 0.37$, $\phi(q + 1 - p) \le 0$, and $R + P \le 1.45p$ give $\frac{\partial \Psi_1}{\partial q}(q, x) = 1 + \beta\bigl[\phi(q + 1 - p) - \phi(q)\bigr] < 0$.
  \item For $p \in [0.04, 0.14]$, we use \cref{eq:small-mass-beta} (where $\overline{P} = J(1 - p)$ for $p \ge 1 - z_0 = 0.225$ and $\overline{P} \ge J(1 - p)$ everywhere) and set
  \[
    \eta(p) \;:=\; \min\Bigl\{1 + \beta(p)p,\, 1 + \frac{m(p) - \lambda}{C'_{\lambda, S}(\ell)}\Bigr\}.
  \]
  Because $x + \eta\phi(x) \le \lambda + (\eta - 1)\phi(\ell) \le m(p)$ for $x \in [\ell, \lambda]$ and $x + \eta\phi(x)$ increases on $[x_h, \ell]$, checking the two endpoints $x = 0$ and $x = x_h$ of the affine piece proves \cref{eq:upper-low-support} for all $x \in [0, \lambda]$. For $z \in [z_-, 1]$, write $e(z) := \lambda - z - \phi(z) \ge 0$ and $d(z) := Q_\lambda(z) - S - J(z) = \int_{z_-}^z e(u)\,\mathrm{d}u$. By concavity of $\phi$ on $[\lambda, 1]$, $e(z)$ is convex on $[z_-, 1]$ with $e(z_-) = 0$, so $d(z) \le \frac{1}{2}(z - z_-)e(z)$. Together with $J(z - p) \le Q_\lambda(z - p) - S$, this bounds $\Psi_3(0, z)$ by
  \begin{equation}
    \Psi_3(0, z) \;\le\; (z - \lambda)(1 + \beta p - \eta) - \frac{1}{2}\beta p^2 + \Bigl[\frac{1}{2}\beta(z - z_-) - \eta\Bigr]e(z).
    \label{eq:H-small-mass-bound}
  \end{equation}
  Since $\frac{1}{2}\beta(z - z_-) - \eta < 0$ and $e(z) \ge 0.26(z - 0.789)_+$ (\cref{prop:shape}), substitution of this lower bound for $e(z)$ gives an upper bound in \cref{eq:H-small-mass-bound} that is affine on $[z_-, 0.789]$ and convex on $[0.789, 1]$, and is negative at $z_-$, $0.789$, and $1$ (\cref{lem:completion-bounds}), while $\partial_q \Psi_1 < 0$ and $\partial_q \Psi_3 < 0$ reduce $q > 0$ to $q = 0$.
  \item For $p \in [0.14, x_h]$, we use \cref{eq:small-mass-beta} and set
  \[
    \eta(p) \;:=\; \inf_{0 \le x < \lambda}\frac{m(p) - x}{\phi(x)},
  \]
  which satisfies \cref{eq:upper-low-support} by construction and $\eta(p) \le 1 + \beta(p)p$ by \cref{lem:completion-bounds}. Because $\partial_q \Psi_1 < 0$ and $\partial_q \Psi_3 < 0$ on $[0.14, x_h]$, it remains to check $\Psi_3(0, z) \le 0$ for $z \in [z_-, 1]$. Rather than bounding $\Psi_3(0, z)$ from scratch, we transport the already verified bound $\Psi_3(0, z) \le 0$ at $p = x_h$ leftward to $p = 0.14$ by showing $\frac{\partial}{\partial p}\Psi_3(0, z) \ge 0$ for almost every $p \in [0.14, x_h]$ and all $z \in [z_-, 1]$. Here $m(p)$ and $\eta(p)$ are Lipschitz, with $\beta'(p) < 0$, $m'(p) < 0$, and every minimizing rank $x_*(p)$ has $\phi(x_*(p))\le A_{\max}$, where $A_{\max}:=\phi_0$ for $p\le0.20$ and $A_{\max}:=C'_{\lambda,S}(\ell)$ for $p\ge0.20$ by \cref{lem:completion-bounds}. Testing the infimum at the same rank at $p+t$ gives $\eta(p+t)-\eta(p)\le[m(p+t)-m(p)]/\phi(x_*(p))$ for $t>0$. At a differentiability point this yields $\eta'(p)\le m'(p)/A_{\max}$, since $m'(p)<0$. With $a := -m'(p)/A_{\max} > 0$, $b := \beta'(p) < 0$, and $v := a + b p + \beta(p)$, using $d(z) \le \frac{1}{2}(z - z_-)e(z)$ bounds $\frac{\partial}{\partial p}\Psi_3(0, z)$ from below by
  \[
    B(z) + \Bigl[a + \frac{1}{2}b(z - z_-)\Bigr]e(z),
    \qquad
    B(z) \;:=\; v(z - \lambda) - \frac{1}{2}b p^2 - \beta(p)p.
  \]
  Because $0 \le e(z) \le \lambda$ on $[z_-, 1]$, checking that this expression is nonnegative at $e = 0$ and $e = \lambda$ for $z \in \{z_-, 1\}$ proves $\frac{\partial}{\partial p}\Psi_3(0, z) \ge 0$ across $[0.14, x_h] \times [z_-, 1]$.
\end{itemize}

\paragraph{3. Masses between $p_+$ and $\lambda^\star$ ($p \in [p_+, \lambda]$).}
Write $p_0 := p_+ = 0.36$, let $(\beta_0, \eta_0)$ be the coefficients at $p_0$ from \cref{eq:eta-balance,eq:beta-def}, and set $m_0 := 2\lambda - \beta_0 J(1 - p_0)$ and
\[
  t_{\mathrm{sup}} \;:=\; \frac{\eta_0\lambda - m_0}{\eta_0 - 1}.
\]
To see that $t_{\mathrm{sup}} \in [0.3058, 0.36]$, let $r_0 \in (0.3058, 0.313)$ solve $C'_{\lambda, S}(r_0) = \lambda - r_0$. Above $c$, $\phi' \le -0.691 < C''_{\lambda, S}$ up to $r_0$, so $\phi \le C'_{\lambda, S}$ there, and $r + \eta_0\phi(r)$ is bounded by its maximum over $[0, r_0]$ where $\phi \le C'_{\lambda, S}$ and over $[r_0, \lambda]$ where $\phi \le \lambda - r$; thus $m_0 \le \eta_0\lambda - (\eta_0 - 1)r_0$, giving $t_{\mathrm{sup}} \ge r_0 > 0.3058$, while $p_0 + \eta_0(\lambda - p_0) \le m_0$ gives $t_{\mathrm{sup}} \le p_0 = 0.36$.

For $p \in [p_0, \lambda]$, where $R := J(p) = Q_\lambda(p) - S$ and $P := J(1 - p) = Q_\lambda(1 - p) - S$, define
\begin{equation}
  \begin{gathered}
    A \;:=\; \lambda - t_{\mathrm{sup}}, \qquad N \;:=\; A R + (\lambda - p)P, \qquad D \;:=\; (2\lambda - t_{\mathrm{sup}})R - (2p - \lambda)P, \\
    \beta(p) \;:=\; \frac{\lambda^2 - t_{\mathrm{sup}} p}{N}, \qquad \eta(p) \;:=\; \frac{D}{N}.
  \end{gathered}
  \label{eq:transport-middle}
\end{equation}
These coefficients solve the linear system $\beta P + \eta A = \lambda + A$ and $\beta R - \eta(\lambda - p) = 2p - \lambda$, and satisfy $\beta > 0$, $1 < \eta \le \eta_0$, and $\eta \le 1 + \beta p$ (\cref{lem:completion-bounds}; note that $D - N = \lambda R - p P \ge 0$ gives $\eta \ge 1$ directly). For $x \in [0, t_{\mathrm{sup}}]$, combining $x + \eta_0\phi(x) \le m_0 = \eta_0\lambda - (\eta_0 - 1)t_{\mathrm{sup}}$ with $\eta \le \eta_0$ gives $x + \eta\phi(x) \le \eta\lambda - (\eta - 1)t_{\mathrm{sup}} = 2\lambda - \beta P$, while for $x \in [t_{\mathrm{sup}}, \lambda]$, $\phi(x) \le \lambda - x$ gives the same bound. For $y \in [p, \lambda]$, we have $\phi(y) = \lambda - y$, so $y + \eta\phi(y) \le p + \eta(\lambda - p) = \lambda - p + \beta R$. Thus both inequalities in \cref{eq:upper-low-support} hold.

Furthermore, along the boundary $q = p$, \cref{eq:transport-middle} and $\phi(z) \le \lambda - z$ give, for all $z \in [\lambda, 1]$,
\[
  \Psi_3(p, z) \;=\; z + p - \lambda - \beta R + \eta\phi(z) \;=\; z - p + \eta\bigl[\phi(z) - (\lambda - p)\bigr] \;\le\; (1 - \eta)(z - p) \;\le\; 0.
\]
Hence, on any box where $\frac{\partial \Psi_3}{\partial q}(q, z) \ge 0$ for $q \in [q_{\mathrm{lo}}, \min\{q_{\mathrm{hi}}, p\}]$, the maximum over $q$ is attained either at $q = p$ (where $\Psi_3(p, z) \le 0$ is already proved) or at $q = q_{\mathrm{hi}} \le p$. On $[0.36, 0.38]$, $\frac{\partial \Psi_3}{\partial q} < 0$ reduces $\Psi_3(q, z) \le 0$ to $q = 0$, which is checked on $[z_-, z_\star]$ against the differential bound at $p_0 = 0.36$ and preserved on $[z_\star, 1]$ by \cref{eq:gap-derivative} using the bounds $\theta = 1/\eta < 0.81$ and $\alpha = \beta/\eta < 3$ (\cref{lem:completion-bounds}). On $[0.38, \lambda]$, the interval checker verifies $\Psi_3(q, z) \le 0$ across all $q \in [0, p]$ on $[z_-, z_\star]$ (using $Q_\lambda(z_\star) - S - Y_\star < 0.00290$) and on $[z_\star, 1]$ (using the quadratic tail values at the fixed anchor $z_A = 0.8624$ from \cref{lem:completion-bounds}).

\subsection{Verifier Coefficients for Masses \texorpdfstring{$p > \lambda^\star$}{p > lambda*}}
\label{app:large-masses}

For $p > \lambda$, condition~\eqref{eq:ver-2} of \cref{lem:verifier} is vacuous. Throughout $(\lambda, 1)$, we set $\beta(p) := p/J(p) > 0$, so the two-class inequality~\eqref{eq:two-class} gives $q + \beta(p)\,\delta_p(q) \le \lambda$ for all $0 \le q \le \lambda$, which proves \cref{eq:ver-3} for $0 \le q \le \lambda$. We specify $c(p, x)$ on two sub-ranges of $(\lambda, 1)$ and then treat high splits $q \in (\lambda, p]$ uniformly at the end.

\paragraph{1. Intermediate masses ($p \in (\lambda, 0.65]$).}
For $p \in (\lambda, 0.65]$, we set
\begin{equation}
  \beta(p) \;:=\; \frac{p}{J(p)},
  \qquad
  c(p, x) \;:=\; c_I(p) \;:=\; \lambda - \beta(p)J(1 - p) \;\ge\; 0,
  \label{eq:constant-correction}
\end{equation}
which is independent of $x$. Because $q + \beta(p)\,\delta_p(q)$ is nonincreasing in $q$ on $[0, \lambda]$ for $p \le 0.65$ (\cref{app:checking-binary}), we have $q + \beta(p)\,\delta_p(q) + c_I(p) \le \beta(p)J(1 - p) + c_I(p) = \lambda$, proving \cref{eq:ver-1}. For \cref{eq:ver-3-norm} ($q \in (\lambda, p]$), since $\delta_p(q) \le (1 - p)\phi(q) \le 0$ by \cref{eq:elementary-packing-bounds}, it is enough to check $r_x - u_x \le \beta(p)(1 - p)$, or equivalently
\[
  x + \beta(p)(1 - p)\phi(x) - \beta(p)J(1 - p) \;\ge\; 0
  \qquad (0 \le x < \lambda).
\]
Let $r_c \le x_c^+ < 0.36$ be the rank where $\phi(r)$ reaches $\lambda - r$, and write $h = 1 - p$. For $x \in [0, r_c]$, monotonicity of $\phi$ and $\phi' \ge -1$ give $J(h) - h\phi(x) \le \frac{1}{2}x^2$ (if $h \ge x$, drop $\int_x^h [\phi(r) - \phi(x)]\,\mathrm{d}r \le 0$; if $h < x$, then $J(h) - h\phi(x) \le h x - \frac{1}{2}h^2 \le \frac{1}{2}x^2$), so $x + \beta h\phi(x) - \beta J(h) \ge x(1 - \frac{1}{2}\beta x) \ge 0$ since $x < 0.36$ and $\beta(p) < 4.966 < 2/0.36$ (\cref{lem:completion-bounds}). On $[r_c, \lambda]$, where $\phi(x) = \lambda - x$, the expression $x + \beta h(\lambda - x) - \beta J(h)$ is affine in $x$ and nonnegative at both endpoints $x = r_c$ and $x = \lambda$ (where it equals $c_I(p) \ge 0$).

For \cref{eq:ver-4-lowq} ($0 \le q \le \lambda < p \le z \le 1$), set $\eta_I(p) := 1 + c_I(p)/\phi(x_h)$. For $x \in [x_h, \lambda)$, we have $r_x \ge 1$ and $\phi(x) \le \phi(x_h)$, so $r_x + u_x = r_x + c_I(p)/\phi(x) \ge \eta_I(p)$; on the affine piece $[0, x_h]$, the inequality $(\lambda - x + c_I(p)) - \eta_I(p)\phi(x) \ge 0$ is affine in $x$ and holds at $x = 0$ and $x = x_h$ (\cref{lem:completion-bounds}). Since $\phi(z) \le 0$, this bounds the left-hand side of \cref{eq:ver-4-lowq} by
\[
  q + z - \lambda + \beta(p)\,K_p(q, z) + \eta_I(p)\phi(z) \;\le\; 0,
\]
which is convex in $z$ on $[p, z_-]$ (where it equals $q - \lambda + \eta_I(p)(\lambda - p) \le 0$ at $z = p$) and is verified on $[z_-, 0.875]$ and $[0.875, 1]$ by the interval checkers (\cref{lem:completion-bounds}).

\paragraph{2. Large masses ($p \in [0.65, 1)$).}
For $p \in [0.65, 1)$, write $h := 1 - p$ and $t(p) := \beta(p)h = p(1 - p)/J(p)$, and set
\begin{equation}
  \beta(p) \;:=\; \frac{p}{J(p)},
  \qquad
  c(p, x) \;:=\; \bigl[\lambda - x - t(p)\phi(x)\bigr]_+.
  \label{eq:large-mass-correction}
\end{equation}
Then $u_x = [r_x - t(p)]_+$ and $r_x - u_x = \min\{r_x, t(p)\} \le t(p) = \beta(p)(1 - p)$, which combines with $\delta_p(q) \le (1 - p)\phi(q) \le 0$ to prove \cref{eq:ver-3-norm} for all $q \in (\lambda, p]$.

For \cref{eq:ver-1}, when $p \in [0.65, 0.88]$, \cref{lem:completion-bounds} gives $r_x < 1.29 < t(p)$ for all $x \in [0, \lambda)$, so $c(p, x) = 0$ and \cref{eq:ver-1} reduces to the two-class inequality $q + \beta(p)\,\delta_p(q) \le \lambda$. When $p \in [0.88, 1)$ and $c(p, x) > 0$, write $v := x - q \ge 0$ and $w := v - \bigl[\phi(q) - \phi(x)\bigr] \ge 0$ (since $\phi' \ge -1$ below $\lambda$). Integrating by parts gives $\delta_p(q) = J(q + h) - J(q) = h\phi(q) - \int_0^h (h - s)\bigl(-\phi'(q + s)\bigr)\,\mathrm{d}s$. Because $0 \le -\phi'(q + s) \le 1$ on $[0, v]$ with integral $\int_0^v \bigl(-\phi'(q + s)\bigr)\,\mathrm{d}s = \phi(q) - \phi(x) = v - w$ and $-\phi'(q + s) \ge 0$ above $v$, the integral against the decreasing weight $(h - s)_+$ is minimized by placing the mass $v - w$ at the right end $[w, v]$ of $[0, v]$:
\[
  \int_0^h (h - s)\bigl(-\phi'(q + s)\bigr)\,\mathrm{d}s \;\ge\; \int_w^v (h - s)_+\,\mathrm{d}s \;=\; \frac{1}{2}\bigl[(h - w)_+^2 - (h - v)_+^2\bigr].
\]
Substituting this bound and $c(p, x) = \lambda - x - t\phi(x)$ into \cref{eq:ver-1} gives
\[
  q + \beta(p)\,\delta_p(q) + c(p, x) - \lambda \;\le\; (t - 1)v - t w - \frac{\beta(p)}{2}\bigl[(h - w)_+^2 - (h - v)_+^2\bigr].
\]
Checking the two cases $v \le h$ and $v \ge h$ shows that the right-hand side is nonpositive for all $w \in [0, v]$ whenever $t \le 2$ and $(t - 1)v \le \frac{1}{2}t h$. By \cref{lem:completion-bounds}, $-6.7 \le \phi' \le -2$ on $[0.88, 1]$ implies $t(p) \le (1 - h)/(1 - 3.35h) < 2$, which satisfies $(t - 1)v \le \frac{1}{2}t h$ for all $v \le 0.18$. If instead $v = x - q \ge 0.18$, then $c(p, x) > 0$ implies $x < x_c^+ < 0.36$ (since $r_x = 1 \le t(p)$ on $[x_c^+, \lambda)$) and hence $q = x - v < 0.18$, where \cref{eq:ver-1} becomes
\begin{equation}
  q - x + t(p)\Bigl[\frac{J(q + h) - J(q)}{h} - \phi(x)\Bigr] \;\le\; 0,
  \label{eq:long-interval-ver1}
\end{equation}
which is certified by the interval checker (\cref{lem:completion-bounds}).

For \cref{eq:ver-4-lowq} ($0 \le q \le \lambda < p \le z \le 1$), since $r_x + u_x \ge 1$ and $\phi(z) \le 0$, it is enough to verify
\[
  E(q, z) \;:=\; q + z - \lambda + \beta(p)\,K_p(q, z) + \phi(z) \;\le\; 0.
\]
For $p \in [0.65, 0.925]$, the interval checkers verify $E(q, z) \le 0$ on $[p, 0.875]$ and $[0.875, 1]$ (\cref{lem:completion-bounds}), using on $[0.88, 0.925]$ the lower bound on $e(z) = \lambda - z - \phi(z)$ from the fixed control $p_f = x = y = 0.34854$: with $m_f := \theta - 1 + \alpha p_f > 0$, $\ell_f(z) := m_f(z - \lambda) - \frac{1}{2}\alpha p_f^2$, and $z_c$ being the zero of $\ell_f$, integrating $d'(z) = e(z) \ge \ell_f(z) + \alpha\,d(z)$ forward from $z_c$ (where $d(z_c) \ge 0$) gives
\begin{equation}
  e(z) \;\ge\; \frac{m_f}{\alpha}\Bigl(\mathrm{e}^{\alpha(z - z_c)_+} - 1\Bigr) \;\ge\; \bigl[\ell_f(z)\bigr]_+ + \frac{\alpha\bigl[\ell_f(z)\bigr]_+^2}{2m_f}.
  \label{eq:e-z-lower-bound}
\end{equation}
For $p \in [0.925, 1)$, where $h = 1 - p \le 0.075$, set $M_H := 8$. On $[0.925, 1]$, we have $\phi' \ge -M_H$ and $J(p) \le h$, so $\beta(p) \ge (1 - h)/h$ and $\phi(u) \le -1 + M_H(1 - u)$. In addition, for $u \in [z, 1]$ and $q \le \lambda$, we have $u - p + q \le \lambda + h$, where $\lambda + h < z_-$ lies on the line $\lambda - r$, so $\phi(u - p + q) \ge \phi(\lambda + h) = -h$. Integrating $\frac{\partial E}{\partial u}(q, u) = 1 + \beta(p)\bigl[\phi(u - p + q) - \phi(u)\bigr] + \phi'(u)$ from $z$ to $1$ therefore gives
\begin{equation}
  E(q, z) - E(q, 1) \;\le\; (1 - z)\Bigl[M_H - 1 - \frac{(1 - h)\bigl(1 - (M_H/2 + 1)h\bigr)}{h}\Bigr] \;\le\; -\frac{17}{24}(1 - z) \;\le\; 0,
  \label{eq:near-one-analytic}
\end{equation}
and at $z = 1$ we have $E(q, 1) = q + \beta(p)\,\delta_p(q) - \lambda \le 0$ by the two-class inequality.

\paragraph{3. High splits ($q \in (\lambda, p]$) for all $p > \lambda$.}
Fix $x \in [0, \lambda)$, $p \in (\lambda, 1)$, and $z \in [p, 1]$, set $s := z - p \ge 0$, and let $F(q)$ be the left-hand side of \cref{eq:ver-4-highq} for $q \in [\lambda, p]$. Since $\beta(p)J(p) = p$ and $K_p(q, z) = -J(p) + \int_q^p \bigl[\phi(u) - \phi(u + s)\bigr]\,\mathrm{d}u$, we have
\[
  F(q) \;=\; s + \beta(p)\int_q^p \bigl[\phi(u) - \phi(u + s)\bigr]\,\mathrm{d}u - r_x\bigl[\phi(q) - \phi(z)\bigr] + u_x\phi(z).
\]
\begin{itemize}
  \item If $\beta(p)s \le r_x$, concavity of $\phi$ on $[\lambda, 1]$ gives $\int_q^p \bigl[\phi(u) - \phi(u + s)\bigr]\,\mathrm{d}u \le s\int_q^p \bigl(-\phi'(u + s)\bigr)\,\mathrm{d}u = s\bigl[\phi(q + s) - \phi(z)\bigr]$. Dropping $u_x\phi(z) \le 0$ and using $\beta(p)s \le r_x$, $\phi' \le -1$, and $r_x \ge 1$ gives
  \[
    F(q) \;\le\; s - r_x\bigl[\phi(q) - \phi(q + s)\bigr] \;\le\; (1 - r_x)s \;\le\; 0.
  \]
  \item If $\beta(p)s \ge r_x$, differentiating $F(q)$ with respect to $q$ and using $\phi(q + s) - \phi(q) \le s\phi'(q)$ (by concavity of $\phi$ on $[\lambda, 1]$) gives, for almost every $q \in [\lambda, p]$,
  \[
    F'(q) \;=\; \beta(p)\bigl[\phi(q + s) - \phi(q)\bigr] - r_x\phi'(q) \;\le\; \bigl(\beta(p)s - r_x\bigr)\phi'(q) \;\le\; 0.
  \]
  Thus $F(q) \le F(\lambda) \le 0$, where $F(\lambda) \le 0$ is the already verified case $q = \lambda$ of \cref{eq:ver-4-lowq}.
\end{itemize}
This completes the verification of \cref{lem:completed-score} across all masses and ranks.

\section{Numerical Verification and Certificates}
\label{app:numerical}

\subsection{Rank Intervals and Scalar Bounds}
\label{app:rank-intervals}

Unless explicitly marked as a rounded display, every terminating decimal below represents an exact rational number; the constant-flux witnesses used by the interval code are exact dyadic rationals given in \cref{eq:certificate-flux-witnesses}. In \cref{sec:variational,app:variational}, we use the parameters
\[
  \begin{gathered}
    \lambda_0 = 0.562325, \qquad \lambda_1 = 0.562330, \\
    p_- = 0.25, \quad \ell = 0.28, \quad x_- = 0.285, \quad p_+ = 0.36, \quad x_+ = 0.38, \quad t_e = 0.385, \\
    z_0 = 0.775, \quad z_B = 0.84, \quad z_1 = 0.875, \quad \calQ = [0.26, 0.28], \quad q_0 = 0.27, \\
    S_+(\lambda) \;=\; 0.0266503148804352 - \frac{0.562330 - \lambda}{2}.
  \end{gathered}
\]
Evaluating $Q_\lambda(r) - S \le J(r) \le C_{\lambda, S}(r)$ at $r = 0.307$ implies, for all $\lambda \in [\lambda_0, \lambda_1]$,
\[
  S \;\ge\; \frac{(1 - r)Q_\lambda(r) - \lambda Q_\lambda(1 - r)}{1 - r - \lambda} \;\ge\; S_- \;:=\; \frac{28711863873}{1306700000000} \;>\; 0.021972.
\]
For every $(p, x, y) \in \calU$, $z \in [z_0, z_1]$, and $(z_T, q) \in [z_B, z_1] \times \calQ$, the complementary ranks satisfy
\[
  1 - p \in [0.64, 0.75],
  \qquad
  z - p \in [0.415, 0.625],
  \qquad
  1 - z_T + q \in [0.385, 0.44],
\]
all of which lie inside the middle interval $[t_e, z_0] = [0.385, 0.775]$ where $J(r) = Q_\lambda(r) - S$. Throughout $\calK_\lambda \times \calU$, the determinant satisfies $\Delta > 0.0519$, and $J(p) \ge Q_\lambda(p) - S$ gives
\[
  \lambda J(p) - p\bigl[Q_\lambda(1 - p) - S\bigr] \;>\; 0.01446 \;>\; 0,
\]
so the $q = 0$ two-class inequality $p\,J(1 - p) \le \lambda J(p)$ is strictly slack on $\calK_\lambda$.

\subsection{Estimates Used in the Variational Proofs}
\label{app:variational-estimates}

\begin{lemma}[Estimates for the variational proofs]
\label{lem:variational-estimates}
For every $\lambda \in [0.562325, 0.562330]$, the following bounds hold:
\begin{enumerate}[label=(\roman*)]
  \item \textbf{Slack removal on the middle ranks (\cref{app:reduction}):} Any primitive $J$ and derivative $\phi = J'$ satisfying \cref{lem:election-1}, \cref{lem:election-2}, and \cref{lem:election-3} has $\sigma(r) = \lambda - r - \phi(r)$ satisfying
  \[
    \int_{0.345}^{0.785} \sigma(r)\,\mathrm{d}r \;<\; \mu_0 \;:=\; 0.00025509854,
    \qquad
    0.02196 \;<\; Q_\lambda(0.5) - J(0.5) \;\le\; S_+(\lambda).
  \]
  Throughout the linear interpolation in \cref{app:reduction} on $\calU \times [z_0, z_1]$, we have $0.0826 < R < 0.136$, $P < 0.135$, $\Delta > 0.0519$, $c_x = U R + k K > 0$, $c_y = U P + h K > 0$, and wherever $H \le \lambda - z$,
  \[
    0 < \alpha := \frac{k\phi(x) + h\phi(y)}{\Delta} < \overline{\alpha} := 4.53,
    \;\;
    |H_P| < \overline{p} := 6.19,
    \;\;
    \omega := \frac{6.19}{4.53}\bigl(\mathrm{e}^{0.453} - 1\bigr) < 0.784.
  \]
  With $\kappa_f = 2.98$, the required slack removal below $0.245$ is at most $0.000761$, whereas available slack below $0.245$ exceeds $0.014975$. Moreover, $[Q_\lambda(0.38)-S-C_{\lambda,S}(0.36)]/0.02>0.09$, $C_{\lambda, S}(x_+) - (Q_\lambda(x_+) - S) > 0.002449158$ and $Q_\lambda(t_e) - S - C_{\lambda, S}(\ell) > 0.0227$, and wherever $H \le \lambda - z$, we have $\phi(y) > 0.0332$, $U R - K(\lambda_1 - x_-) > 0.00030254668 > 0$ (so $0 < \theta < 1$), and $\bigl[Q_\lambda(0.35) - S - C_{\lambda, S}(0.34)\bigr]/0.01 - 0.035 > 0.0618$.
  \item \textbf{Uniform bounds on $\calK_\lambda$ (\cref{app:attainment,app:shape}):} The quadratic trial function $\overline{\phi}$ in \cref{app:attainment} satisfies $L_C(\overline{S}, \overline{\phi}) \le -0.0033 < 0$ on $(\ell, t_e]$. On $\calK_\lambda \times \calU \times [z_0, z_1]$ with $Y \ge Q_\lambda - Q_\lambda(1)$, we have $\Delta < 0.116$, $c_x > 0.0145$, $c_y > 0.024$, and for any triple within $0.005$ of the essential infimum in \cref{eq:high-ode}, $c_x/c_y < 0.74 < v_p$, $\phi(x) > 0.0218 > \varepsilon_0 := 0.02$, and $\phi(y) > 0.0289 > \varepsilon_0$. At $(z_1, q_0) = (0.875, 0.27)$, $T_{\lambda, S}(z_1, q_0) - [Q_\lambda(z_1) - Q_\lambda(1)] > 0.019$, and $|\partial T_{\lambda, S}(z, q)/\partial S| < b_S := 0.765$ for all $(z, q) \in [z_B, z_1] \times \calQ$.
  \item \textbf{Tube bounds and slope estimates (\cref{app:shape}):} For any feasible pair $(S, \phi)$ at $\lambda \le 0.56232847$, the reference comparison in \cref{sec:reference-comparison,app:calibration} gives
  \[
    \begin{gathered}
      0.02271790 < S < 0.023354, \qquad C_{\lambda, S}(c_{\mathrm{ref}}) - J(c_{\mathrm{ref}}) < 1.91061404 \times 10^{-6}, \\
      \int_{p_{\mathrm{ref}}}^{t_e} \sigma(r)\,\mathrm{d}r < \varepsilon_{\mathrm{area}} := 8.688098 \times 10^{-8}, \qquad Q_\lambda(z_{\mathrm{ref}}) - S - Y(z_{\mathrm{ref}}) < 1.711874 \times 10^{-7},
    \end{gathered}
  \]
  where $c_{\mathrm{ref}} = 0.2950935005377889$, $p_{\mathrm{ref}} = 0.34854007977695545$, and $z_{\mathrm{ref}} = 0.7835342428237777$. These tube bounds apply to every feasible pair on this parameter range; the algebraic coefficient and slope bounds below hold on their stated domains, and their contact and shape implications apply to the minimizing profile in the proof of \cref{prop:shape} in the order proved there.

  Across the entire low-rank interval $[\ell, t_e]$, we have
  \[
    \sup_{[\ell, t_e]} |C''_{\lambda, S}| \;\le\; \frac{2 \cdot 0.562330 \cdot 0.023354}{(1 - 0.385)^3} \;<\; 0.113 \;<\; 0.116.
  \]
  On this interval of $S$, active first positions $x$ satisfy $\phi(x) > 0.1106$ and $\theta > \theta_\circ := 0.116$; on $[\ell, c_{\mathrm{ref}}]$, $|C''_{\lambda, S}| \le K_\circ := 0.074987345523$, and $\frac{1}{2}(\theta_\circ - K_\circ)(c_{\mathrm{ref}} - c_{\mathrm{safe}})^2 > 1.92685556 \times 10^{-6}$ at $c_{\mathrm{safe}} = 0.2854$. For $\epsilon = 0.000418$, we have $\epsilon^2/2 > \varepsilon_{\mathrm{area}}$. The diagonal slope $L_{\mathrm{diag}}$ in \eqref{eq:diagonal-slope} satisfies $L_{\mathrm{diag}} > -1$ on $0.285 \le x \le p \le p_b := 0.35$, $0 \le \sigma(p) \le \sigma_0 := 0.022$, and satisfies $L_{\mathrm{diag}} > -1 + \gamma_{\mathrm{tail}}$ ($\gamma_{\mathrm{tail}} = 0.022$) on $p_{\mathrm{ref}} \le p \le 0.36$, $0.2918 \le x \le p$, $0 \le \sigma(p) \le \varepsilon_{\mathrm{low}} := 0.00024$. For the minimizing profile, iterating the supporting-line and area bounds gives
  \[
    \begin{aligned}
      (c > 0.2854) &\;\Rightarrow\; \theta > 0.6809 \;\Rightarrow\; c > 0.2913 \;\Rightarrow\; \theta > 0.6907 \\
      &\;\Rightarrow\; c > \underline{c} := 0.2918 \;\Rightarrow\; \theta > \underline{\theta} := 0.691,
    \end{aligned}
  \]
  together with $c < \overline{c} := 0.305$, $\theta < \overline{\theta} := 0.78$, $\alpha > 2.30$, and $\theta + \alpha p_- > 1.26 = 1 + \gamma_{\mathrm{high}}$ ($\gamma_{\mathrm{high}} := 0.26$).
\end{enumerate}
\end{lemma}

\paragraph{Justification of the finite LP certificate in Part~(i).}
The linear-programming certificate in Part~(i) uses a uniform grid of $N = 200$ intervals with $t_i = i/N$ ($0 \le i \le N$) and variables
\[
  J_i \;=\; J(t_i),
  \qquad
  V_i \;=\; \int_0^{t_i} J(r)\,\mathrm{d}r.
\]
All LP rows are necessary conditions at $L := \lambda_1 = 0.562330$ (since any $\phi$ satisfying the election inequalities at $\lambda \le L$ also satisfies them at $L$). The derivative bound $J'(r) \le L - r$ gives $\max\{0, Q_L(r) - Q_L(1)\} \le J(r) \le Q_L(r)$ and corresponding bounds on $V_i$. Writing $h = 1/N$ and $t = t_i$, integrating $J'(r) \le L - r$ forward and backward over $[t_i, t_{i+1}]$ gives
\[
  N(J_{i+1} - J_i) \;\le\; L - t - \frac{h}{2},
  \qquad
  N(V_{i+1} - V_i) - J_i \;\le\; (L - t)\frac{h}{2} - \frac{h^2}{6},
  \qquad
  J_{i+1} - N(V_{i+1} - V_i) \;\le\; (L - t)\frac{h}{2} - \frac{h^2}{3}.
\]
The two-class rows evaluate \cref{eq:two-class} either at grid ranks or integrated over a grid bin of $q$ in $[0, \min\{p, L\}]$ (replacing $L - q$ by its bin average), so they are linear in $(J_i, V_i)$ with exact rational coefficients.

The certificate then appends $66$ three-class rows sequentially using a linearizing ratio substitution. Averaging \cref{eq:three-class} independently over a low-rank bin for $x$, a low-rank bin for $y$, and a high-rank bin for $z$ with bin centers $a, b, z$ and bin averages $A, B, C$ of $\phi$ (where $k := p + b - L > 0$ on all bins in the certificate), set
\[
  U \;=\; z - L,
  \qquad
  h_* \;=\; 2L - a,
  \qquad
  R \;=\; J(p),
  \qquad
  X \;=\; \frac{J(1 - p)}{R},
  \qquad
  Y \;=\; \frac{\operatorname{avg}_z\bigl[J(z - p) - J(z)\bigr]}{R}.
\]
Given certified nonnegative lower bounds $x \le X$ and $y \le Y$ from earlier dual certificates, the linearized row is
\begin{equation}
  (U + k y)A + (U x + h_* y)B + (h_* - k x)C \;\le\; 0.
  \label{eq:lp-moving-row}
\end{equation}
To see that \cref{eq:lp-moving-row} follows from the bin-averaged three-class inequality $(U + k Y)A + (U X + h_* Y)B + (h_* - k X)C \le 0$, write $M := k A + h_* B$. If $M \le 0$, then $A \le L - a$, $B \le L - b$, $C \le -U$, and $x \le X \le L/p$ give $(U + k y)A + (U x + h_* y)B + (h_* - k x)C \le U(-L + p x) + y M \le 0$. If $M > 0$, then $h_* - k X \ge h_* - k L/p > 0$ (by the $q = 0$ two-class inequality and $a, b < L$), and the upper bound
\[
  H(x, y) \;:=\; -\frac{(U + k y)A + (U x + h_* y)B}{h_* - k x}
\]
has negative partial derivatives $H_x = -\frac{M(U + k y)}{(h_* - k x)^2} < 0$ and $H_y = -\frac{M}{h_* - k x} < 0$, so $C \le H(X, Y) \le H(x, y)$, proving \cref{eq:lp-moving-row}. By induction over the sequence of $66$ rows, every row holds for every admissible $\phi$.

For the resulting linear system $A w \le b$ with box bounds $l \le w \le u$ and any rational dual multiplier vector $\pi \le 0$, weak LP duality gives
\[
  c \cdot w \;\ge\; \pi \cdot b + \sum_j \min\bigl\{e_j l_j,\, e_j u_j\bigr\},
  \qquad
  e \;=\; c - A^\top \pi,
\]
which the verifier checks in exact rational arithmetic, giving
\[
  \int_{0.345}^{0.785} \sigma(r)\,\mathrm{d}r \;\le\; \frac{159436587170745658497841277055894957}{625000000000000000000000000000000000000} \;<\; \mu_0,
\]
and the lower and upper bounds on $J(1/2)$ that define $S_-$ and $S_+(\lambda)$.
The lower polynomial for $\Delta$ has Bernstein coefficients $0.052023192$, $0.05493359928$, $0.055721837976666\dots$, and $0.05505340809$, so even after the second-phase adjustment $\Delta > 0.0519$ holds throughout. The bound on $[Q_\lambda(0.38)-S-C_{\lambda,S}(0.36)]/0.02$ follows by substitution: this expression decreases with both $S$ and $\lambda$ on the stated box, and its value at $(\lambda_1,S_+(\lambda_1))$ exceeds $0.099$. To verify $Q_\lambda(t_e) - S - C_{\lambda, S}(\ell) > 0.0227$, note that
\[
  Q_\lambda(t_e) - S - C_{\lambda, S}(\ell) \;=\; \lambda\Bigl[t_e - \lambda + \frac{1 - \ell}{2}\Bigr] - \frac{t_e^2}{2} - S\Bigl(1 - \frac{\lambda}{1 - \ell}\Bigr)
\]
has negative derivatives with respect to both $\lambda$ and $S$ on the allowed domain, and evaluating at $(\lambda_1, S_+(\lambda_1))$ gives $\frac{21349005266668987}{937500000000000000} > 0.0227$.

For the two-rank Lipschitz comparison in Step~3 of the proof of \cref{prop:shape}, we have $0.024 < c_y < 0.14$, $0.2 < c_y/\Delta < 2.7$, $|\alpha| < 7.4$, $|1 - p - \lambda| < 0.188$, and $|\lambda - z + p| < 0.148$, giving the geometric slope bound
\[
  7.4\Bigl(0.313 + 0.74 \cdot 0.188 + 0.148\,\frac{0.116}{0.024}\Bigr) + 1 \;<\; 9.7,
\]
and the coupled Lipschitz bound
\[
  \frac{2.7}{0.2}\Bigl[7.4\Bigl(0.313 + 0.74 \cdot 0.188 + 0.148\,\frac{0.116}{0.024}\Bigr) + 2\Bigr] \;<\; 144.
\]
For the slope and contact transitions in Steps~4 and~5 of the proof of \cref{prop:shape},
\[
  \frac{\epsilon^2}{2} > \varepsilon_{\mathrm{area}}, \qquad p_+ + \epsilon < x_+, \qquad \frac{\varepsilon_{\mathrm{low}}^2}{2(1 - \underline{\theta})} > \varepsilon_{\mathrm{area}}, \qquad p_{\mathrm{ref}} + \frac{\varepsilon_{\mathrm{low}}}{\gamma_{\mathrm{tail}}} < 0.3594492 < p_+,
\]
and $\frac{0.26}{2}(z_{\mathrm{ref}} - 0.782)^2 > 3.06 \times 10^{-7} > 1.711874 \times 10^{-7}$. Finally, for the parameter shift in Step~2 of the proof of \cref{prop:attainment}, the difference $\bigl[C_{\lambda+\delta, \widetilde{S}}(r) - \widetilde{J}(r)\bigr] - \bigl[C_{\lambda, S}(r) - J(r)\bigr]$ equals
\[
  \delta\Bigl[2\lambda - 0.6788 + \frac{r}{2} - (r - 0.38)_+ - \frac{S + 0.2012\lambda}{1 - r}\Bigr] + \delta^2\Bigl(1 - \frac{0.2012}{1 - r}\Bigr) \;>\; 0.338\delta \;>\; 0,
\]
while the derivative of $E_3$ with respect to $\lambda$ is below $-0.0027010815 < 0$. The verifier scripts for \cref{lem:variational-estimates} are listed in \cref{tab:checkers}.

\subsection{Calibration and the Two Numerical Endpoints}
\label{app:calibration}

The reference curve used in \cref{sec:reference-comparison} has the exact rational parameters
\[
  \begin{gathered}
    \lambda_{\mathrm{ref}} = 0.5623285, \qquad S_{\mathrm{ref}} = 0.02303623505254907, \qquad M = 0.135, \\
    c_{\mathrm{ref}} = 0.2950935005377889, \qquad p_{\mathrm{ref}} = 0.34854007977695545, \\
    \zeta_0 = 0.7835342428237777, \qquad \zeta_1 = 0.8624179577667147, \qquad q_{\mathrm{ref}} = 0.27010770106016474,
  \end{gathered}
\]
together with piecewise rational-cubic pieces for $\phi_{\mathrm{ref}}$ on $[c_{\mathrm{ref}}, p_{\mathrm{ref}}]$ and piecewise-affine control maps $\pi(s), \upsilon(s)$ on $[\zeta_0, \zeta_1]$. All second-derivative enclosures of the logarithmic ODE coefficients are certified over the full scalar interval $0.02197280467819 \le S \le 0.026650314880436$ and the bounding box of segments between $\gamma(r)$ and $r$.

At each tested value of $\lambda \in \{\lambda_{\mathrm{ref}}, \underline{\lambda}\}$, the certificate reader constructs $J_{\mathrm{ref}, \lambda}$ and the gauge $u_\lambda$ from the same rational rank nodes and raw curvature/control data in \path{research/closure/root/warped_cubic_ordered_reference_242.json} (or its \texttt{Llower} companion at $\underline{\lambda}$) so that the endpoint values and integrals match at $\lambda$ exactly. Specifically, writing $a = c_{\mathrm{ref}}$, $b = p_{\mathrm{ref}}$, $h = b - a$, and $k(r) > 0$ for the piecewise-affine raw curvature on $[a, b]$, set
\[
  f_a \;=\; C'_{\lambda, S_{\mathrm{ref}}}(a),
  \qquad
  f_b \;=\; \lambda - b,
  \qquad
  I \;=\; Q_\lambda(b) - S_{\mathrm{ref}} - C_{\lambda, S_{\mathrm{ref}}}(a),
\]
\[
  t \;=\; \frac{2I - h(f_a + f_b)}{\int_a^b (r - a)(b - r)k(r)\,\mathrm{d}r},
  \qquad
  v \;=\; \frac{f_b - f_a + t\int_a^b (b - r)k(r)\,\mathrm{d}r}{h}.
\]
The reference derivative and primitive at $\lambda$ are then the piecewise-cubic and piecewise-quartic rationals
\[
  f(r) \;=\; f_a + v(r - a) - t\int_a^r (r - s)k(s)\,\mathrm{d}s,
  \qquad
  J_{\mathrm{ref}, \lambda}(r) \;=\; C_{\lambda, S_{\mathrm{ref}}}(a) + \int_a^r f(s)\,\mathrm{d}s,
\]
which satisfy $f(a) = f_a$, $f(b) = f_b$, and $J_{\mathrm{ref}, \lambda}(b) = Q_\lambda(b) - S_{\mathrm{ref}}$. Similarly, given the raw gauge polynomial $d(w)$ in $w = (r - a)/h \in [0, 1]$, define the endpoint and integral residuals
\[
  e_0 \;=\; -\frac{\lambda}{(1 - a)^2} - d(0),
  \qquad
  e_1 \;=\; -d(1),
  \qquad
  e_I \;=\; \frac{-1 + \lambda/(1 - a)}{h} - \int_0^1 d(w)\,\mathrm{d}w,
\]
set $A = e_0$, $C = 3(e_0 + e_1 - 2e_I)$, $B = e_1 - e_0 - C$, and $\widetilde{d}(w) = d(w) + A + B w + C w^2$, and define
\[
  u_\lambda(r) \;=\; -\frac{\lambda}{1 - a} + h\int_0^{(r - a)/h} \widetilde{d}(w)\,\mathrm{d}w.
\]
Then $u_\lambda(a) = -\lambda/(1 - a)$, $u_\lambda(b) = -1$, $u'_\lambda(a) = -\lambda/(1 - a)^2$, and $u'_\lambda(b) = 0$, and the certificate verifies $m_S(r) = J'_{\mathrm{ref}, \lambda}(r) + (S - S_{\mathrm{ref}})u'_\lambda(r) > 0$ across $[a, b]$ for all $S$ in the calibrated interval.

\begin{lemma}[Numerical calibration]
\label{lem:numerical-calibration}
Evaluating the second-variation coefficients and first-variation integrals in \cref{eq:log-comparison,eq:adjoint-expansion,eq:coercive-master} gives
\[
  \kappa > 68.429, \qquad K_{\mathrm{rank}} < 1.708912, \qquad W_0 = 3, \qquad \inf_{[c_{\mathrm{ref}}, p_{\mathrm{ref}}]} W > 3.24367,
\]
and hence $\frac{8}{9}W_0 - K_{\mathrm{rank}} > 0.95775534 > 0$.
\begin{enumerate}[label=(\roman*)]
  \item \textbf{Tube bounds at $\lambda_{\mathrm{ref}} = 0.5623285$:} With $k_0 = k_{\mathrm{ref}}$ from \cref{eq:certificate-flux-witnesses}, we have
  \[
    \begin{gathered}
      A_S \in [-0.000177118091, -0.000065361439], \\
      \varepsilon < 2.582780877 \times 10^{-6}, \qquad \Lambda_{\mathrm{ref}} > -4.264600252 \times 10^{-6},
    \end{gathered}
  \]
  so the constant budget $-\Lambda_{\mathrm{ref}} + \varepsilon + A_S^2/(4\kappa)$ is below $6.847495738 \times 10^{-6}$. Thus every feasible competitor satisfies
  \[
    |S - S_{\mathrm{ref}}| < 0.000317626822, \qquad \calJ < 7.149525 \times 10^{-6},
  \]
  and keeping the boundary and initial-deficit terms gives $E_{\mathrm{in}} < 1.711874 \times 10^{-7}$, $p_{\mathrm{ref}} - \gamma(p_{\mathrm{ref}}) < 4.06384409 \times 10^{-7}$, and $J(p_{\mathrm{ref}}) - Q_\lambda(p_{\mathrm{ref}}) + S < 8.688098 \times 10^{-8}$.
  \item \textbf{Lower bound $\lambda^\star > \underline{\lambda} = 0.56232844$:} Re-evaluating at $\underline{\lambda} = 0.56232844$ with $k_0 = k_{\mathrm{low}}$ from \cref{eq:certificate-flux-witnesses} gives $\Lambda_{\mathrm{ref}} > 3.4940949161 \times 10^{-6}$, $|A_S| < 0.000074509648$, $\varepsilon < 2.582784988 \times 10^{-6}$, $\kappa > 68.4307154$, and $\frac{8}{9}W_0 - K_{\mathrm{rank}} > 0.957759899$, so the constant term in \cref{eq:coercive-master} exceeds $9.11289646 \times 10^{-7} > 0$ and feasibility is impossible at or below $\underline{\lambda}$.
  \item \textbf{Upper bound $\lambda^\star < \overline{\lambda} = 0.56232847$:} At $\overline{\lambda} = 0.56232847$, an explicit rational-spline trial in $\calK_{\overline{\lambda}}$ satisfies all two-class bounds in \cref{eq:high-binary-compatibility} with strict margin exceeding $7.06118753 \times 10^{-9}$ (after subtracting all root-bracketing, layer-interpolation, and ODE integration enclosures).
\end{enumerate}
\end{lemma}

\subsection{Numerical Bounds for Completing the Curve and Checking the Verifier Lemma}
\label{app:completion-bounds}

In \cref{sec:completing-score,app:upper-all}, we use the localized domain $\lambda \in [\lambda_0, \lambda_1]$, $S \in [0.02271790, 0.023354]$, and the parameters
\[
  \begin{gathered}
    x_h = 0.2635, \quad c_- = 0.2918, \quad c < 0.305, \quad \underline{\theta} = 0.691, \quad k = 1 - \underline{\theta} = 0.309, \\
    x_c^+ = 0.35945, \quad z_- = 0.782, \quad z_+ = 0.789, \quad \overline{p}_{\mathrm{ref}} = 0.348541, \quad \epsilon_\sigma = 0.00024, \quad \epsilon_E = 8.7 \times 10^{-8}.
  \end{gathered}
\]

\begin{lemma}[Numerical bounds for completing the curve]
\label{lem:completion-bounds}
On the localized parameter domain above, the completed function $\phi^\star$ and primitive $J^\star$ satisfy:
\begin{enumerate}[label=(\roman*)]
  \item \textbf{Boundary pieces on $[0, x_h]$ and $[z_\star, 1]$ (\cref{app:joining-curve}):} The affine piece in \cref{eq:head-affine} satisfies $0.436 < \phi_0 < 0.443$, $0.67 < -s_h < 0.689$, and $C'_{\lambda, S}(x_h) < 0.26$. The fixed control triple $p_f = 0.27$, $x = y = 0.2918$ implies that departure from $\lambda - r$ occurs before $0.788066162 < z_+$. The tangency equations~\eqref{eq:tangency-identities} place the active contact in $0.8495 < z_\star < 0.875$ and $0.2635 \le q_\star < 0.2764$, and the quadratic piece in \cref{eq:tail-quadratic,eq:tail-backward} satisfies $-6.5 < \phi'(r) < -2.8$, $c_2 < -0.15$, $6.18 < v < 6.42$, and $10.77 < \kappa_{\mathrm{tail}} < 12.73$. For every $p \in [x_h, p_+] = [0.2635, 0.36]$ with $0 < \theta < 0.78$ and $0 < \alpha < 4.53$, the endpoint values of \cref{eq:gap-derivative} at $r = z_\star$ and $r = 1$ exceed $0.11539$ and $1.22184$. Writing $J_{\mathrm{tail}}$ and $\phi_{\mathrm{tail}}$ for the polynomial extensions of the quadratic tail primitive and derivative on $[z_\star, 1]$, at the fixed anchor $z_A = 0.8624$ they satisfy $J_{\mathrm{tail}}(z_A) \in [0.08793, 0.08840]$ and $\phi_{\mathrm{tail}}(z_A) \in [-0.3589, -0.3514]$ (whether $z_A \ge z_\star$ or $z_A < z_\star$), and the deficit at $z_\star$ satisfies $Q_\lambda(z_\star) - S - Y_\star < 0.00290$. On $[c_-, t_e]$, the primitive $J$ satisfies
  \begin{align*}
    J(q) &\;\ge\; Q_\lambda(q) - S + \bigl[E(c_-) - \sigma(c_-)(q - c_-)\bigr]_+ && (q \ge c_-), \\
    J(r) &\;\le\; Q_\lambda(r) - S + \epsilon_E + \epsilon_\sigma(\overline{p}_{\mathrm{ref}} - r)_+ + \frac{k}{2}(\overline{p}_{\mathrm{ref}} - r)_+^2 && (r \ge c).
  \end{align*}
  \item \textbf{Two-class inequality checks (\cref{app:checking-binary}):} For $0<p\le x_h$, the head lower bound on $[\lambda J(p)-pJ(1-p)]/p$ exceeds $0.0736$; for $x_h\le p\le0.64$, the middle lower bound on $\lambda J(p)-pJ(1-p)$ exceeds $0.000556$. For $p \in [1 - \lambda, 0.65]$, the numerator of the $q$-derivative of $q + \frac{p}{J(p)}\delta_p(q)$ is below $-0.0138$; for $p \in [0.65, 0.775]$, it is negative by $q = 0.18$, and \cref{eq:linear-q-strip} is nonpositive on $q \in [0, 0.18]$. On $[0.775, 0.807]$, $Q_\lambda(p) - S - J(p) \le 0.798(p - 0.782)_+^2 < Q_\lambda(p) - S - T_H(p)$, while at $z_1 = 0.875$ every active ODE slope exceeds $-0.388$ and the two-class derivative is below $-0.396$. On $[z_\star, p_a]$ ($p_a \in (0.91257, 0.91339)$), \cref{eq:binary-curvature} satisfies $J''(p) - T''_{\mathrm{bin}}(p) > 6.84 \times 10^{-7}$. On $h = 1 - p \in [0.02, 0.09]$, $(1 - h)\kappa(h)$ is bounded above by the affine majorant
  \[
    (1 - h)\kappa(h) \;\le\; 2.122 + 140(S - 0.02271790) - \bigl[3.774 + 700(S - 0.02271790)\bigr]h,
  \]
  and \cref{eq:tail-binary-integral} holds with \cref{eq:remaining-slack-envelope} on $q \in [0.2918, 0.337]$, $h \in [0.02, 0.0875]$, $c \in [0.2918, 0.305]$. For $h \le 0.02$, $\kappa(h) \le \kappa(0) \le 2.122 + 140(S - 0.02271790)$, and integrating \cref{eq:sigma-pointwise-floor} at $h = 0.02$ certifies $q \in [0.25, 0.337]$.
  \item \textbf{Verifier Lemma checks (\cref{app:small-middle-masses,app:large-masses}):} For $p \le 0.04$, $J(p) \ge 0.39p$, $0.84p \le J(1 - p) \le p$, $\beta p < 0.49$, and $(0.562325 \cdot 0.39 - 0.04)/1.39 > 0.1289$. For $p \in [0.14, 0.2635]$, the value of $\phi$ at a minimizing rank $x_*(p)$ of $(m(p) - x)/\phi(x)$ is at most $\phi_0$ when $p \le 0.20$, and when $p \ge 0.20$ head exclusion and monotonicity on $[x_h, \ell]$ place $x_*(p) \ge 0.28$, where $\phi(x_*(p)) \le C'_{\lambda, S}(0.28)$. On $[0.2635, 0.36]$, $\beta(p)\bigl[\phi_0 - \lambda + 0.782 - p\bigr] > 1$. For $p \in [0.36, 0.38]$, the coefficients in \cref{eq:transport-middle} satisfy $\theta = 1/\eta < 0.81$ and $\alpha = \beta/\eta < 3$, with endpoint values of \cref{eq:gap-derivative} at $r = z_\star$ and $r = 1$ exceeding $0.42451$ and $2.46301$; for $p \in [0.38, \lambda]$, the separate quadratic-tail checker certifies $\Psi_3(q, z) \le 0$ on $[z_\star, 1]$. On $[\lambda, 0.65]$, \cref{eq:constant-correction} gives $\beta < 4.966$ and $c_I < 0.1$, with margin $> 0.04$ on $[0, x_h]$. On $[0.65, 0.88]$, $r_x = (\lambda - x)/\phi(x) < 1.29 < p(1 - p)/J(p)$ for all $x < \lambda$, so $c(p, x) = 0$. On $[0.88, 1)$ ($h \le 0.12$), $-6.7 \le \phi' \le -2$ gives
  \[
    t \;=\; \beta(p)h \;\le\; \frac{1 - h}{1 - 3.35h} \;<\; 2
    \qquad\text{and}\qquad
    \frac{2.35 \cdot 0.18}{0.88} \;<\; \frac{1}{2},
  \]
  certifying $x - q \le 0.18$, while \cref{eq:long-interval-ver1} certifies $x - q \ge 0.18$. Thus, the completed function $\phi^\star$ satisfies every hypothesis of \cref{lem:completed-score}.
\end{enumerate}
\end{lemma}

\begin{table}[!t]
\centering
\footnotesize
\caption{Verification scripts in \texttt{numerical\_certificates.zip} and representative certified margins.}
\label{tab:checkers}
\vspace{4pt}
\begin{tabular}{@{}
  >{\raggedright\arraybackslash}p{0.25\linewidth}
  >{\raggedright\arraybackslash}p{0.45\linewidth}
  >{\raggedright\arraybackslash}p{0.23\linewidth}@{}}
\toprule
\textbf{Assertion} & \textbf{Checker script(s)} & \textbf{Certified margin} \\
\midrule
Variational reduction and compact-domain constants &
  \texttt{lower/verify\_global\_constants.py};\newline
  \texttt{global/verify\_layer\_extension.py};\newline
  \texttt{global/verify\_master\_compact\_constants.py} &
  $\Delta > 0.0519$, $\omega < 0.784$, $\phi(y) > 0.0332$ \\[3pt]
Slater perturbation and contact/slope bounds &
  \texttt{root/verify\_weighted\_slater.py};\newline
  \texttt{core/verify\_left\_contact\_margin.py};\newline
  \texttt{global/verify\_diagonal\_slope\_band.py};\newline
  \texttt{root/verify\_expanded\_support\_bounds.py};\newline
  \texttt{global/verify\_retained\_endpoint\_bounds.py};\newline
  \texttt{global/verify\_tail\_diagonal\_extension.py};\newline
  \texttt{global/verify\_cap\_contact\_upper.py} &
  $c \in (0.2918, 0.305)$, $\theta \in (0.691, 0.78)$, $L_{\mathrm{diag}} > -0.978$ \\[3pt]
Calibration coefficients &
  \texttt{global/interval\_warped\_hessian.py} &
  $\kappa > 68.429$, $K_{\mathrm{rank}} < 1.708912$ \\[3pt]
First variation and lower bound $\lambda^\star > 0.56232844$ &
  \texttt{root/transport\_first\_variation.py} &
  lower margin $> 9.11289646 \times 10^{-7}$ \\[3pt]
Feasible trial and upper bound $\lambda^\star < 0.56232847$ &
  \texttt{root/selected\_family\_feasibility.py} &
  trial margin $> 7.06118753 \times 10^{-9}$ \\[3pt]
Terminal tangency and quadratic piece on $[z_\star, 1]$ &
  \texttt{core/verify\_right\_terminal\_contact.py};\newline
  \texttt{core/verify\_binary\_contact\_hermite.py};\newline
  \texttt{global/verify\_fixed\_anchor\_terminal.py} &
  $z_\star \in (0.8495, 0.875)$, $c_2 < -0.15$ \\[3pt]
Differential bound preservation on $[z_\star, 1]$ &
  \texttt{root/verify\_hermite\_cap\_extension.py} &
  join $> 0.11539$, end $> 1.22184$ \\[3pt]
Two-class inequality for low and middle masses &
  \texttt{global/verify\_structural\_binary\_low.py};\newline
  \texttt{global/verify\_cap\_contact\_upper.py};\newline
  \texttt{global/verify\_structural\_binary\_central.py} &
  $60{,}396$ certified leaves \\[3pt]
Two-class inequality with rational maximizing $q$ &
  \texttt{core/verify\_binary\_high.py};\newline
  \texttt{core/verify\_binary\_rational\_tail\_fast.py} &
  curvature gap $> 6.84 \times 10^{-7}$ \\[3pt]
Two-class inequality for masses near $p = 1$ &
  \texttt{core/verify\_binary\_mid\_tail\_fast.py};\newline
  \texttt{global/verify\_near\_one\_binary.py};\newline
  \texttt{core/verify\_binary\_tail\_constants.py};\newline
  \texttt{core/verify\_tail\_deficit\_ratio.py} &
  $283{,}258$ and $838{,}476$ certified leaves \\[3pt]
Small masses $p \in (0, x_h]$ &
  \texttt{global/verify\_noncritical\_small\_layers.py};\newline
  \texttt{global/verify\_structural\_left\_extension.py} &
  all coefficient signs hold \\[3pt]
Middle masses $p \in [x_h, \lambda^\star]$ &
  \texttt{root/verify\_critical\_split\_coupled.py};\newline
  \texttt{global/verify\_right\_support\_transport.py};\newline
  \texttt{global/verify\_synthetic\_tail\_constants.py};\newline
  \texttt{global/verify\_right\_support\_tail.py} &
  upper $< -1.5200 \times 10^{-11}$ \\[3pt]
Constant correction $p \in [\lambda^\star, 0.65]$ &
  \texttt{global/verify\_intermediate\_low\_credit.py};\newline
  \texttt{global/certify\_structural\_intermediate.py} &
  upper $< -2.5149 \times 10^{-11}$ \\[3pt]
Large masses $p \in [0.65, 1)$ &
  \texttt{root/verify\_structural\_large\_hc.py};\newline
  \texttt{root/verify\_structural\_tail\_hc.py};\newline
  \texttt{root/verify\_zero\_credit\_price.py};\newline
  \texttt{global/verify\_structural\_long\_credit.py} &
  upper $< -1.9049 \times 10^{-9}$, $< -3.4704 \times 10^{-6}$ \\
\bottomrule
\end{tabular}
\end{table}

\subsection{Arithmetic and Reproducibility}
\label{app:verifier-implementation}

The companion archive \texttt{numerical\_certificates.zip} contains $43$ verification scripts, each of which certifies a finite-dimensional inequality over the parameter box specified in that script; the analytical reductions in \cref{app:variational,app:upper-all} connect those finite-dimensional inequalities to the continuum variational problem and the completed function $\phi^\star$. (In particular, the feasible trial at $\overline{\lambda}$ certifies that $\calK_{\overline{\lambda}}$ contains a pair with $j_{\overline{\lambda}} < 0$, whereas the all-election upper bound uses the minimizing function $\phi^\star$ at $\lambda^\star$ together with its proved regularity from \cref{prop:shape} and the verifier checks of \cref{app:upper-all}.)

\paragraph{Arithmetic enclosures.}
Every script uses either exact rational arithmetic (\texttt{fractions.Fraction}) or outward-rounded IEEE-754 binary64 interval arithmetic. Rational polynomial bounds use Bernstein coefficients or centered Taylor ranges. In the interval scripts, rational inputs are enclosed in directed-rounded binary64 intervals before evaluation, binary64 literals represent exact dyadic rationals, basic operations are expanded outward by one ULP using \texttt{numpy.nextafter}, square roots are enclosed either by integer square-root bounds or by \texttt{numpy.sqrt} expanded outward with \texttt{nextafter}, and exponentials are bounded by positive Taylor polynomials with explicit geometric remainders. Continuum inequalities are proved by adaptive box subdivision (splitting a box at its midpoint until the interval enclosure has strict sign, a proved monotonicity reduction applies, or a prerequisite inequality is certified violated on the box); no continuum inequality is inferred from point sampling or numerical optimization. For the reference-curve checks, the scripts also report the total number of cells, the number verified, and any uncovered rank intervals.

\paragraph{Reference flows and the feasible trial.}
All piecewise-polynomial knots and control-map breakpoints are aligned with grid nodes so that every integration cell is smooth. For the linear reference flow $Z' = \alpha Z + \beta_0$, second-order Taylor steps enclose the third derivative
\[
  Z''' \;=\; (\alpha'' + 3\alpha\alpha' + \alpha^3)Z + (2\alpha' + \alpha^2)\beta_0 + \alpha\beta'_0 + \beta''_0
\]
over each step interval, and the logarithmic Hessian bounds use interval jets together with a $2 \times 2$ LDL Schur-complement test over a box containing the entire segment between the reference curve and the competitor.

For the feasible trial at $\overline{\lambda}$, concavity of the trial function $\phi$ on $[x_-, x_+]$ sets the optimal second low rank to $y = \max\{x, p\}$, and the maximizing first rank $x \in [x_-, x_+]$ is enclosed by sign-changing derivative brackets (with endpoints checked separately). On the mass grid $p \in [p_-, p_+] = [0.25, 0.36]$---which includes the exact phase knots $c$ and $r_c$ so that any upward jump in the $p$-derivative at $p = r_c$ occurs at a grid node---each cell has a certified semiconcavity bound below $1000$, allowing linear interpolation between adjacent grid nodes with an explicit penalty of $1000\,\delta_p^2/8$ (and similarly across the terminal grid $q \in \calQ = [0.26, 0.28]$). The finite-grid ODE solution $Y_{\mathrm{grid}}(z)$ is first enclosed by interval Taylor steps (which account for switches of the active $p$-branch and certify $|Y''_{\mathrm{grid}}| < 2$ almost everywhere and $Q_{\overline{\lambda}} - Q_{\overline{\lambda}}(1) \le Y_{\mathrm{grid}} \le Q_{\overline{\lambda}} - S$, so $0 < K < 0.111$); subtracting the coefficient-rounding and $p$-grid interpolation penalties then shows that $Y_{\mathrm{grid}}$ is a subsolution of the continuum ODE~\eqref{eq:high-ode} and transfers the positive terminal margin to the continuum continuation $Y(z)$.

In \cref{sec:reference-comparison,lem:numerical-calibration}, the constant-flux absorption uses the exact dyadic rationals
\begin{equation}
  k_{\mathrm{ref}} \;=\; -\frac{4742825210076883}{281474976710656} \;\approx\; -16.8499,
  \qquad
  k_{\mathrm{low}} \;=\; -\frac{4742818173202465}{281474976710656} \;\approx\; -16.849875,
  \label{eq:certificate-flux-witnesses}
\end{equation}
for which the certificates verify $P_1(r) - k_0 - W_0 > 0$ on every rank cell.

\paragraph{Running the verification suite.}
After extracting \texttt{numerical\_certificates.zip} into a fresh directory, run the following commands from that directory with Python~3.10+ and NumPy installed (without \texttt{-O}, so that assertions remain active):
\begin{quote}\small
\begin{verbatim}
python3 research/closure/simplified/replay_certificates.py --check-hashes
python3 research/closure/simplified/replay_certificates.py --all --dry-run
python3 research/closure/simplified/replay_certificates.py --all
\end{verbatim}
\end{quote}
The \texttt{--check-hashes} step verifies the SHA-256 hashes of all $83$ archived scripts and data files against the manifest before replay (since running the suite refreshes the output JSON files), and \texttt{--all} runs all $43$ commands serially and writes individual logs and a summary report. \cref{tab:checkers} lists the scripts in \texttt{numerical\_certificates.zip} (with paths relative to \texttt{research/closure}) and their certified margins.

\begingroup
\small
\patchcmd{\thebibliography}{\section*{\refname}}{\section*{\phantomsection\refname}\addcontentsline{toc}{section}{\protect\numberline{}\refname}}{}{}
\begin{thebibliography}{CCFHRWY26}

\bibitem[ABEPS18]{AnshelevichEtAl18}
Elliot Anshelevich, Onkar Bhardwaj, Edith Elkind, John Postl, and Piotr Skowron.
\newblock \href{https://doi.org/10.1016/j.artint.2018.07.006}{Approximating optimal social choice under metric preferences}.
\newblock \emph{Artificial Intelligence}, 264:27--51, 2018.

\bibitem[ABP15]{AnshelevichBhardwajPostl15}
Elliot Anshelevich, Onkar Bhardwaj, and John Postl.
\newblock \href{https://doi.org/10.1609/aaai.v29i1.9308}{Approximating optimal social choice under metric preferences}.
\newblock In \emph{Proceedings of the 29th AAAI Conference on Artificial Intelligence (AAAI)}, pages 777--783, 2015.

\bibitem[CCFHRWY26]{CaiEtAl26b}
Ziyi Cai, Moses Charikar, Fabian Frank, Jabari Hastings, Prasanna Ramakrishnan, Kangning Wang, and Qilin Ye.
\newblock Stable voting rules on the edge of optimal metric distortion.
\newblock Working paper, 2026.

\bibitem[CJMW20]{ChengEtAl20}
Yu~Cheng, Zhihao Jiang, Kamesh Munagala, and Kangning Wang.
\newblock \href{https://doi.org/10.1145/3417750}{Group fairness in committee selection}.
\newblock \emph{ACM Transactions on Economics and Computation (TEAC)}, 8(4):23:1--23:18, 2020.

\bibitem[CLRVW25]{CharikarEtAl25}
Moses Charikar, Alexandra Lassota, Prasanna Ramakrishnan, Adrian Vetta, and Kangning Wang.
\newblock \href{https://doi.org/10.1145/3717823.3718235}{Six candidates suffice to win a voter majority}.
\newblock In \emph{Proceedings of the 57th Annual ACM Symposium on Theory of Computing (STOC)}, pages 1590--1601, 2025.

\bibitem[CR22]{CharikarRamakrishnan22}
Moses Charikar and Prasanna Ramakrishnan.
\newblock \href{https://doi.org/10.1137/1.9781611977073.116}{Metric distortion bounds for randomized social choice}.
\newblock In \emph{Proceedings of the 2022 ACM-SIAM Symposium on Discrete Algorithms (SODA)}, pages 2986--3004, 2022.

\bibitem[CRTW25]{CharikarEtAl25b}
Moses Charikar, Prasanna Ramakrishnan, Zihan Tan, and Kangning Wang.
\newblock \href{https://arxiv.org/abs/2505.13630}{Metric distortion for tournament voting and beyond}.
\newblock In \emph{Proceedings of the 26th ACM Conference on Economics and Computation (EC)}, pages 790--818, 2025.

\bibitem[CRWW24]{CharikarEtAl24}
Moses Charikar, Prasanna Ramakrishnan, Kangning Wang, and Hongxun Wu.
\newblock \href{https://doi.org/10.1145/3689625}{Breaking the metric voting distortion barrier}.
\newblock \emph{Journal of the ACM}, 71(6):42:1--42:33, 2024.

\bibitem[Fis84]{Fishburn84}
Peter~C. Fishburn.
\newblock \href{https://doi.org/10.2307/2297786}{Probabilistic social choice based on simple voting comparisons}.
\newblock \emph{Review of Economic Studies}, 51(4):683--692, 1984.

\bibitem[Kre65]{Kreweras65}
Germain Kreweras.
\newblock Aggregation of preference orderings.
\newblock In \emph{Mathematics and Social Sciences I}, pages 73--79, 1965.

\bibitem[PS21]{PulyassarySwamy21}
Haripriya Pulyassary and Chaitanya Swamy.
\newblock \href{https://arxiv.org/abs/2111.08698}{On the randomized metric distortion conjecture}.
\newblock \emph{arXiv preprint arXiv:2111.08698}, 2021.

\end{thebibliography}
\endgroup

\end{document}
